%% GENERATED by generate_papers.py -- edit the bank_*.py files, not this file.
\documentclass[11pt,a4paper]{article}
\usepackage[a4paper, top=20mm, bottom=20mm, left=18mm, right=18mm, headheight=14pt]{geometry}
\usepackage[T1]{fontenc}
\usepackage{lmodern}
\usepackage{amsmath, amssymb, mathtools}
\usepackage{braket}
\usepackage{array, tabularx, booktabs, multicol}
\usepackage{enumitem}
\usepackage[dvipsnames]{xcolor}
\usepackage[most]{tcolorbox}
\usepackage{needspace}
\usepackage{fancyhdr, lastpage}
\usepackage{pifont}
\usepackage{tikz}
\usepackage{graphicx}
\usepackage{eso-pic}
% SILICO FELLER watermark: official logo, centred, ~8 % opacity, behind all content, on every page
\AddToShipoutPictureBG{\AtPageCenter{\makebox(0,0){\tikz\node[opacity=0.08, inner sep=0pt]{\includegraphics[width=0.72\paperwidth]{../../../assets/branding/silico_feller_logo.png}};}}}
\usepackage[hidelinks]{hyperref}
\newcolumntype{L}{>{\raggedright\arraybackslash}X}
\definecolor{brand}{HTML}{1F4E79}
\definecolor{okgreen}{HTML}{1B7F3B}
\setlength{\parindent}{0pt}
\setlength{\parskip}{3pt}
\newcommand{\optlabel}[1]{\textbf{(#1)}}
\newcommand{\correctlabel}[1]{\textcolor{okgreen}{\textbf{(#1)}\,\ding{51}}}
\newenvironment{qitem}[1]{\begin{list}{}{\setlength{\leftmargin}{3.3em}\setlength{\labelwidth}{3.1em}%
  \setlength{\labelsep}{0.2em}\setlength{\topsep}{3pt}\setlength{\parsep}{3pt}}\item[\textbf{Q.#1}\hfill]}{\end{list}}
\newcommand{\qsep}{\par\vspace{2pt}{\color{black!15}\hrule}\vspace{4pt}}
\newcommand{\qmeta}[2]{{\footnotesize\color{black!60}\textsc{#1} \;$\cdot$\; #2}\par}
\newtcolorbox{passagebox}[1]{enhanced, breakable, colback=brand!4, colframe=brand!70, fonttitle=\bfseries,
  title={#1}, left=6pt, right=6pt, top=4pt, bottom=4pt}
\newtcolorbox{sectionbanner}{enhanced, colback=brand, colframe=brand, coltext=white, fontupper=\large\bfseries,
  left=8pt, right=8pt, top=6pt, bottom=6pt, sharp corners}
\newtcolorbox{infobox}[1]{enhanced, breakable, colback=black!3, colframe=black!45, fonttitle=\bfseries,
  title={#1}, left=6pt, right=6pt, top=4pt, bottom=4pt}
\newtcolorbox{solbox}[1]{enhanced, breakable, colback=okgreen!3, colframe=okgreen!60!black, fonttitle=\bfseries,
  title={Answer: (#1) \quad Detailed solution}, left=6pt, right=6pt, top=3pt, bottom=4pt, before skip=4pt, after skip=6pt}
\allowdisplaybreaks
\setlength{\emergencystretch}{3em}

\begin{document}
\pagestyle{fancy}\fancyhf{}
\fancyhead[L]{\small\color{black!70}Shor's Algorithm --- Assessment}
\fancyhead[R]{\small\color{black!70}\textbf{Answer Book}}
\fancyfoot[C]{\small Page \thepage\ of \pageref*{LastPage}}
\renewcommand{\headrulewidth}{0.4pt}

\thispagestyle{empty}
\begin{center}
{\color{brand}\rule{\linewidth}{1.2pt}}\\[6pt]
{\Huge\bfseries\color{brand} Shor's Algorithm}\\[4pt]
{\Large Number theory, order finding, phase estimation, implementation and cryptography}\\[4pt]
{\large\bfseries ANSWER BOOK}\\[2pt]
{\color{brand}\rule{\linewidth}{1.2pt}}
\end{center}
This book accompanies the separate \emph{Question Paper} (80 questions, 3 hours, 240 marks). It contains
\begin{enumerate}[leftmargin=1.4em, itemsep=0pt, topsep=2pt]
\item the \textbf{answer key}, section by section, with the topic tested by every question (page~\pageref{key}),
\item a \textbf{topic-wise map} grouping the questions by theme and a \textbf{quick-marking grid} (page~\pageref{topicmap}), and
\item the \textbf{Answer Bank}: every question restated with the correct option marked \textcolor{okgreen}{\ding{51}} and a
      complete solution with derivations (from page~\pageref{bank}).
\end{enumerate}
\begin{infobox}{Marking scheme (every question)}
Each question has \textbf{four} options (A), (B), (C) and (D). \textbf{ONLY ONE} of these four options is correct.\\[2pt]
\begin{tabular}{@{}ll@{}}
\emph{Full marks}: & $+3$ if ONLY the correct option is chosen;\\
\emph{Zero marks}: & $\phantom{+}0$ if none of the options is chosen (question unanswered);\\
\emph{Negative marks}: & $-1$ in all other cases (wrong option, or more than one option chosen).
\end{tabular}
\end{infobox}

\begin{infobox}{How to evaluate}
For each section count $N_{\mathrm c}$ = number of correct answers and $N_{\mathrm w}$ = number of questions answered wrongly (or with more than one option marked). Section marks $= 3N_{\mathrm c} - N_{\mathrm w}$; unanswered questions score $0$. The total is out of $240$.
Every numerical answer was re-computed independently (exact integer and \texttt{Fraction} arithmetic, NumPy phase-estimation formulas, checks against the code and recorded outputs of \texttt{Shor.ipynb}, and CUDA-Q state-vector and density-matrix simulations with the notebook's circuit builder) by
\texttt{verify\_answers.py} in the source folder.
\end{infobox}
\begin{infobox}{Conventions (identical to the Question Paper)}
\footnotesize
\begin{itemize}[leftmargin=1.2em, itemsep=1pt, topsep=2pt]
\item $a \bmod N$ is the remainder in $\{0, \dots, N - 1\}$; the \textbf{order} $r = \operatorname{ord}_N(a)$ is the least $r > 0$ with $a^r \equiv 1 \pmod N$ (defined for $\gcd(a, N) = 1$);
      $\varphi$ is Euler's totient, $\lambda$ the Carmichael function. ``Factors'' means non-trivial factors ($1 < p < N$).
\item \textbf{Reduction:} if $r$ is even and $a^{r/2} \not\equiv -1 \pmod N$, then $\gcd(a^{r/2} - 1, N)$ and $\gcd(a^{r/2} + 1, N)$ are non-trivial factors; otherwise the attempt fails.
\item \textbf{Circuit:} $t$ counting qubits (Hadamards), work register in $\ket1$, controlled $U_{a^{2^j} \bmod N}$ with $U_c\ket y = \ket{cy \bmod N}$ for $y < N$, inverse QFT, measurement of the
      counting register giving $y \in \{0, \dots, 2^t - 1\}$ (most significant bit first). \textbf{Post-processing:} \texttt{Fraction(y, 2**t).limit\_denominator(N)} gives $s/r'$;
      the order is $r'$ or its smallest multiple $kr'$ with $a^{kr'} \equiv 1$. $y = 0$ gives no candidate.
\item \textbf{Notebook} (\texttt{Shor.ipynb}, defaults): $N = 15$, $a = 7$, $n = $ \texttt{N.bit\_length()} work qubits, $t = 2n$ counting qubits, multipliers as exact
      permutation matrices (\texttt{cudaq.register\_operation}), $1000$ shots. \textbf{Previously recorded run}: CUDA-Q 0.16.0, CPU target \texttt{qpp-cpu}, no GPU; $12$ qubits,
      $51$ gates (plus $8$ measurements), outcomes $0, 64, 128, 192$; $0.766$ of the shots gave $15 = 3 \times 5$. No physical quantum processor was used.
\item State-vector memory: $16$ bytes per amplitude ($2^q$ amplitudes for $q$ qubits). $1$ MB $= 10^6$ bytes, $1$ GB $= 10^9$ bytes.
\end{itemize}
\end{infobox}
\newpage
\section*{Answer key}\label{key}
{\small Answer distribution: (A) 20, (B) 20, (C) 20, (D) 20.}

\needspace{14\baselineskip}
\subsection*{Section A --- Mathematical and Number Theory Foundations}
{\small\renewcommand{\arraystretch}{1.15}
\begin{tabularx}{\linewidth}{@{}r c L @{\hspace{1em}} r c L@{}}
\toprule
\textbf{Q.} & \textbf{Ans.} & \textbf{Topic tested} & \textbf{Q.} & \textbf{Ans.} & \textbf{Topic tested}\\ \midrule
1 & \textbf{A} & Euclidean algorithm & 11 & \textbf{A} & Failure: $a^{r/2} \equiv -1$ \\
2 & \textbf{A} & Modular inverse & 12 & \textbf{C} & Fermat's little theorem \\
3 & \textbf{D} & Euler's totient $\varphi(pq)$ & 13 & \textbf{D} & Trial-division bound \\
4 & \textbf{B} & Multiplicative order & 14 & \textbf{C} & Success bound $1 - 2^{-(k-1)}$ \\
5 & \textbf{B} & Repeated squaring & 15 & \textbf{C} & $\varphi$, order, $\lambda$, inverse \\
6 & \textbf{D} & Square-and-multiply cost & 16 & \textbf{D} & Classical factoring complexity \\
7 & \textbf{B} & Orders divide $\lambda(N)$ & 17 & \textbf{A} & Non-trivial square root $\Rightarrow$ factor \\
8 & \textbf{A} & Lucky gcd probability & 18 & \textbf{C} & Euler's theorem vs order \\
9 & \textbf{D} & Chinese remainder theorem & 19 & \textbf{B} & Order of $3$ mod $91$ \\
10 & \textbf{B} & Square roots of $1$ mod $15$ & 20 & \textbf{C} & Factors of $91$ by gcd \\
\bottomrule
\end{tabularx}}

\needspace{14\baselineskip}
\subsection*{Section B --- Quantum Computing Foundations}
{\small\renewcommand{\arraystretch}{1.15}
\begin{tabularx}{\linewidth}{@{}r c L @{\hspace{1em}} r c L@{}}
\toprule
\textbf{Q.} & \textbf{Ans.} & \textbf{Topic tested} & \textbf{Q.} & \textbf{Ans.} & \textbf{Topic tested}\\ \midrule
21 & \textbf{B} & Uniform superposition & 31 & \textbf{B} & CNOT cost of a Fredkin gate \\
22 & \textbf{D} & Measuring the work register & 32 & \textbf{C} & Implicit measurement \\
23 & \textbf{C} & Unitarity of $U_a$ & 33 & \textbf{A} & Notebook qubit count \\
24 & \textbf{D} & Eigenvalues of $U_a$ & 34 & \textbf{D} & State-vector memory \\
25 & \textbf{B} & $\ket1$ as sum of eigenvectors & 35 & \textbf{B} & Roles of circuit blocks \\
26 & \textbf{D} & Phase kickback & 36 & \textbf{B} & Phase gates $R_1(\lambda)$ \\
27 & \textbf{C} & Phases of $U^{2^j}$ & 37 & \textbf{C} & Kickback and controls \\
28 & \textbf{A} & Hadamard transform signs & 38 & \textbf{D} & Parallelism vs interference \\
29 & \textbf{C} & Two-path interference & 39 & \textbf{A} & QPE: exact phase \\
30 & \textbf{A} & Reversible modular exponentiation & 40 & \textbf{A} & QPE: inexact phase \\
\bottomrule
\end{tabularx}}

\needspace{14\baselineskip}
\subsection*{Section C --- Shor's Algorithm Fundamentals}
{\small\renewcommand{\arraystretch}{1.15}
\begin{tabularx}{\linewidth}{@{}r c L @{\hspace{1em}} r c L@{}}
\toprule
\textbf{Q.} & \textbf{Ans.} & \textbf{Topic tested} & \textbf{Q.} & \textbf{Ans.} & \textbf{Topic tested}\\ \midrule
41 & \textbf{D} & Steps of Shor's algorithm & 51 & \textbf{D} & Exact success, notebook default \\
42 & \textbf{B} & $N = 15$, $a = 2$ & 52 & \textbf{A} & Probability $\varphi(r)/r$ \\
43 & \textbf{C} & Failure for $a = N - 1$ & 53 & \textbf{A} & Role of the QFT in Shor \\
44 & \textbf{A} & Odd order: retry & 54 & \textbf{C} & Skipped identity multipliers \\
45 & \textbf{C} & Lucky base & 55 & \textbf{B} & Functions of \texttt{Shor.ipynb} \\
46 & \textbf{A} & Register sizes for $N = 33$ & 56 & \textbf{D} & Outcomes to fractions \\
47 & \textbf{C} & $N = 15$, $a = 11$ outcomes & 57 & \textbf{D} & Prime powers \\
48 & \textbf{D} & Continued fractions, $N = 33$ & 58 & \textbf{B} & Uninformative outcomes \\
49 & \textbf{B} & Multiples test & 59 & \textbf{A} & $N = 39$: order and factors \\
50 & \textbf{B} & Condition $2^t \ge N^2$ & 60 & \textbf{C} & $N = 39$: continued fraction \\
\bottomrule
\end{tabularx}}

\needspace{14\baselineskip}
\subsection*{Section D --- Advanced Shor's Algorithm and Implementation}
{\small\renewcommand{\arraystretch}{1.15}
\begin{tabularx}{\linewidth}{@{}r c L @{\hspace{1em}} r c L@{}}
\toprule
\textbf{Q.} & \textbf{Ans.} & \textbf{Topic tested} & \textbf{Q.} & \textbf{Ans.} & \textbf{Topic tested}\\ \midrule
61 & \textbf{B} & Beauregard qubit count & 71 & \textbf{D} & Memory for $N = 1023$ \\
62 & \textbf{C} & Published RSA-2048 estimates & 72 & \textbf{A} & Toy RSA \\
63 & \textbf{C} & $O(n^3)$ scaling & 73 & \textbf{A} & Post-quantum standards \\
64 & \textbf{D} & Semiclassical QFT & 74 & \textbf{C} & Impact on RSA, ECC, AES \\
65 & \textbf{A} & Approximate QFT & 75 & \textbf{D} & Fault-tolerance terms \\
66 & \textbf{B} & Controlled phases in the QFT & 76 & \textbf{C} & Qubit counts of constructions \\
67 & \textbf{A} & Recorded gate counts & 77 & \textbf{B} & Harvest now, decrypt later \\
68 & \textbf{D} & No hardware execution & 78 & \textbf{A} & Qubits and repeated squaring \\
69 & \textbf{C} & CNOTs of a controlled multiplier & 79 & \textbf{B} & Error-free probability vs success \\
70 & \textbf{B} & Random baseline & 80 & \textbf{D} & Interpreting noisy results \\
\bottomrule
\end{tabularx}}

\needspace{18\baselineskip}
\subsection*{Topic-wise map of the questions}\label{topicmap}
{\small Every question appears under at least one theme; several appear under two.\par}
{\small\renewcommand{\arraystretch}{1.2}
\begin{tabularx}{\linewidth}{@{}>{\raggedright\arraybackslash}p{0.42\linewidth}L@{}}
\toprule
\textbf{Theme} & \textbf{Question numbers}\\ \midrule
Modular arithmetic, gcd, inverses and totients & 1, 2, 3, 9, 12, 15, 72 \\
Multiplicative order and repeated squaring & 4, 5, 6, 7, 18, 19, 27, 78 \\
Reduction of factoring to order finding & 8, 10, 11, 14, 17, 20, 41, 42, 43, 44, 45, 57, 59 \\
Classical factoring and complexity & 13, 16, 63 \\
Superposition, entanglement and measurement & 21, 22, 30, 32, 38 \\
Unitary $U_a$, eigenvectors and phase kickback & 23, 24, 25, 26, 37 \\
Gates and interference & 28, 29, 31, 36, 66, 69 \\
Phase estimation and the QFT & 35, 39, 40, 53, 64, 65, 66 \\
Outcomes, continued fractions and post-processing & 47, 48, 49, 50, 51, 52, 56, 58, 60 \\
Register sizes, memory and simulation limits & 33, 34, 46, 50, 71, 78 \\
\texttt{Shor.ipynb}: code and recorded results & 51, 54, 55, 67, 68 \\
Gate-level circuits and noise & 31, 69, 70, 79, 80 \\
Resource estimates and fault tolerance & 61, 62, 64, 75, 76 \\
RSA, ECC and post-quantum cryptography & 72, 73, 74, 77 \\
\bottomrule
\end{tabularx}}

\subsection*{Quick-marking grid}
\begin{center}\renewcommand{\arraystretch}{1.3}
\begin{tabular}{|c|c|c|c|c|c|c|c|c|c|}\hline
1\,\textbf{A} & 2\,\textbf{A} & 3\,\textbf{D} & 4\,\textbf{B} & 5\,\textbf{B} & 6\,\textbf{D} & 7\,\textbf{B} & 8\,\textbf{A} & 9\,\textbf{D} & 10\,\textbf{B} \\
\hline
11\,\textbf{A} & 12\,\textbf{C} & 13\,\textbf{D} & 14\,\textbf{C} & 15\,\textbf{C} & 16\,\textbf{D} & 17\,\textbf{A} & 18\,\textbf{C} & 19\,\textbf{B} & 20\,\textbf{C} \\
\hline
21\,\textbf{B} & 22\,\textbf{D} & 23\,\textbf{C} & 24\,\textbf{D} & 25\,\textbf{B} & 26\,\textbf{D} & 27\,\textbf{C} & 28\,\textbf{A} & 29\,\textbf{C} & 30\,\textbf{A} \\
\hline
31\,\textbf{B} & 32\,\textbf{C} & 33\,\textbf{A} & 34\,\textbf{D} & 35\,\textbf{B} & 36\,\textbf{B} & 37\,\textbf{C} & 38\,\textbf{D} & 39\,\textbf{A} & 40\,\textbf{A} \\
\hline
41\,\textbf{D} & 42\,\textbf{B} & 43\,\textbf{C} & 44\,\textbf{A} & 45\,\textbf{C} & 46\,\textbf{A} & 47\,\textbf{C} & 48\,\textbf{D} & 49\,\textbf{B} & 50\,\textbf{B} \\
\hline
51\,\textbf{D} & 52\,\textbf{A} & 53\,\textbf{A} & 54\,\textbf{C} & 55\,\textbf{B} & 56\,\textbf{D} & 57\,\textbf{D} & 58\,\textbf{B} & 59\,\textbf{A} & 60\,\textbf{C} \\
\hline
61\,\textbf{B} & 62\,\textbf{C} & 63\,\textbf{C} & 64\,\textbf{D} & 65\,\textbf{A} & 66\,\textbf{B} & 67\,\textbf{A} & 68\,\textbf{D} & 69\,\textbf{C} & 70\,\textbf{B} \\
\hline
71\,\textbf{D} & 72\,\textbf{A} & 73\,\textbf{A} & 74\,\textbf{C} & 75\,\textbf{D} & 76\,\textbf{C} & 77\,\textbf{B} & 78\,\textbf{A} & 79\,\textbf{B} & 80\,\textbf{D} \\
\hline\end{tabular}\end{center}
\newpage
\section*{Answer Bank --- questions with detailed solutions}\label{bank}

\begin{sectionbanner}SECTION A \enspace$\cdot$\enspace Mathematical and Number Theory Foundations \hfill Q.1--20\end{sectionbanner}

\needspace{13\baselineskip}
\begin{qitem}{1}
\qmeta{Single-correct MCQ}{Euclidean algorithm}
Using the Euclidean algorithm, $\gcd(1071, 462)$ equals
\par\smallskip\noindent\begin{tabularx}{\linewidth}{@{}*{4}{L}@{}}
\correctlabel{A}\,$21$ & \optlabel{B}\,$7$ & \optlabel{C}\,$147$ & \optlabel{D}\,$42$
\end{tabularx}
\end{qitem}
\begin{solbox}{A}
\textbf{\color{brand}Correct answer.} \textcolor{okgreen}{\textbf{(A)}} $21$

\textbf{\color{brand}Concepts required.} The Euclidean algorithm replaces $(a, b)$ by $(b, a \bmod b)$ until the remainder is $0$; the last non-zero remainder is the $\gcd$.

\textbf{\color{brand}Relevant formulae.} $\gcd(a, b) = \gcd(b, a \bmod b)$: $1071 = 2\cdot462 + 147$, $462 = 3\cdot147 + 21$, $147 = 7\cdot21 + 0$.

\textbf{\color{brand}Step-by-step explanation.} $1071 = 2\cdot462 + 147$; $\ 462 = 3\cdot147 + 21$; $\ 147 = 7\cdot21 + 0$. The last non-zero remainder is $21$.
\emph{Check}: $1071 = 3^2\cdot7\cdot17$ and $462 = 2\cdot3\cdot7\cdot11$, common part $3\cdot7 = 21$. The algorithm needs $O(\log N)$ division steps,
which is why the gcd steps of Shor's algorithm are classically cheap.

\textbf{\color{brand}Why the other options are incorrect.}
\begin{itemize}[leftmargin=2.4em, itemsep=1pt, topsep=2pt]
\item[\textbf{(B)}] $7$ divides both numbers, but so does $21 = 3\cdot7$; $7$ is a common divisor, not the greatest one.
\item[\textbf{(C)}] $147$ is the first remainder of the algorithm, not its last non-zero remainder; $147 \nmid 462$.
\item[\textbf{(D)}] $42$ does not divide $1071$ (odd), so it cannot be a common divisor at all.
\end{itemize}
\textbf{\color{brand}Key takeaway.} The final $\gcd$ step of Shor's algorithm is this fast classical algorithm: $O(\log N)$ division steps.
\end{solbox}
\qsep

\needspace{13\baselineskip}
\begin{qitem}{2}
\qmeta{Single-correct MCQ}{Modular inverse}
The multiplicative inverse of $3$ modulo $20$ (the $b$ with $3b \equiv 1 \pmod{20}$, $0 < b < 20$) is
\par\smallskip\noindent\begin{tabularx}{\linewidth}{@{}*{4}{L}@{}}
\correctlabel{A}\,$7$ & \optlabel{B}\,$13$ & \optlabel{C}\,$17$ & \optlabel{D}\,$3$
\end{tabularx}
\end{qitem}
\begin{solbox}{A}
\textbf{\color{brand}Correct answer.} \textcolor{okgreen}{\textbf{(A)}} $7$

\textbf{\color{brand}Concepts required.} An inverse modulo $m$ exists iff $\gcd(b, m) = 1$; it is found by the extended Euclidean algorithm or by testing.

\textbf{\color{brand}Relevant formulae.} $3b \equiv 1 \pmod{20}$: $3\cdot7 = 21 = 20 + 1$.

\textbf{\color{brand}Step-by-step explanation.} Extended Euclid: $20 = 6\cdot3 + 2$, $3 = 1\cdot2 + 1$, so $1 = 3 - 2 = 3 - (20 - 6\cdot3) = 7\cdot3 - 20$. Hence $3\cdot7 = 21 \equiv 1 \pmod{20}$,
$b = 7$. The inverse exists because $\gcd(3, 20) = 1$; the same condition $\gcd(a, N) = 1$ makes multiplication by $a$ modulo $N$ reversible.

\textbf{\color{brand}Why the other options are incorrect.}
\begin{itemize}[leftmargin=2.4em, itemsep=1pt, topsep=2pt]
\item[\textbf{(B)}] $3\cdot13 = 39 \equiv 19 \equiv -1 \pmod{20}$: $13$ is the inverse of $-3$, not of $3$.
\item[\textbf{(C)}] $3\cdot17 = 51 \equiv 11 \pmod{20}$, which is not $1$.
\item[\textbf{(D)}] $3$ is the number itself; $3\cdot3 = 9 \not\equiv 1 \pmod{20}$.
\end{itemize}
\textbf{\color{brand}Key takeaway.} Invertibility of $a$ modulo $N$ is exactly what makes the multiplier $U_a$ a permutation (unitary).
\end{solbox}
\qsep

\needspace{13\baselineskip}
\begin{qitem}{3}
\qmeta{Single-correct MCQ}{Euler's totient $\varphi(pq)$}
Euler's totient $\varphi(91)$ equals
\par\smallskip\noindent\begin{tabularx}{\linewidth}{@{}*{4}{L}@{}}
\optlabel{A}\,$84$ & \optlabel{B}\,$90$ & \optlabel{C}\,$78$ & \correctlabel{D}\,$72$
\end{tabularx}
\end{qitem}
\begin{solbox}{D}
\textbf{\color{brand}Correct answer.} \textcolor{okgreen}{\textbf{(D)}} $72$

\textbf{\color{brand}Concepts required.} Euler's totient counts the residues coprime to $N$; it is multiplicative for coprime factors.

\textbf{\color{brand}Relevant formulae.} $\varphi(pq) = (p - 1)(q - 1)$ for distinct primes: $\varphi(7\cdot13) = 6\cdot12 = 72$.

\textbf{\color{brand}Step-by-step explanation.} $91 = 7\cdot13$ with distinct primes, so $\varphi(91) = (7-1)(13-1) = 6\cdot12 = 72$: the number of $1 \le a \le 91$ coprime to $91$.
$90 = N - 1$ would hold only for a prime $N$.

\textbf{\color{brand}Why the other options are incorrect.}
\begin{itemize}[leftmargin=2.4em, itemsep=1pt, topsep=2pt]
\item[\textbf{(A)}] $84 = 91 - 7$ removes only the multiples of $13$ ($7$ of them) and forgets the multiples of $7$.
\item[\textbf{(B)}] $90 = 91 - 1$ would be $\varphi$ of a prime; $91 = 7\cdot13$ is composite.
\item[\textbf{(C)}] $78 = 91 - 13$ removes only the multiples of $7$ and forgets the multiples of $13$.
\end{itemize}
\textbf{\color{brand}Key takeaway.} Knowing $\varphi(N)$ is equivalent to knowing the factors of an RSA modulus.
\end{solbox}
\qsep

\needspace{13\baselineskip}
\begin{qitem}{4}
\qmeta{Single-correct MCQ}{Multiplicative order}
The multiplicative order of $2$ modulo $11$ is
\par\smallskip\noindent\begin{tabularx}{\linewidth}{@{}*{4}{L}@{}}
\optlabel{A}\,$11$ & \correctlabel{B}\,$10$ & \optlabel{C}\,$2$ & \optlabel{D}\,$5$
\end{tabularx}
\end{qitem}
\begin{solbox}{B}
\textbf{\color{brand}Correct answer.} \textcolor{okgreen}{\textbf{(B)}} $10$

\textbf{\color{brand}Concepts required.} The order is the least $r > 0$ with $a^r \equiv 1$; it divides $\varphi(N)$ (Lagrange).

\textbf{\color{brand}Relevant formulae.} Powers of $2$ mod $11$: $2, 4, 8, 5, 10, 9, 7, 3, 6, 1$ --- the first $1$ is at $r = 10$.

\textbf{\color{brand}Step-by-step explanation.} Powers of 2 mod 11: $2, 4, 8, 5, 10, 9, 7, 3, 6, 1$ --- the first 1 appears at exponent $10$. The order divides $\varphi(11) = 10$, so the only
candidates are $1, 2, 5, 10$: $2^2 = 4$, $2^5 = 32 \equiv 10 \equiv -1$, so the order is $10$ ($2$ is a primitive root mod 11).

\textbf{\color{brand}Why the other options are incorrect.}
\begin{itemize}[leftmargin=2.4em, itemsep=1pt, topsep=2pt]
\item[\textbf{(A)}] $11$ is the modulus; orders modulo a prime $p$ divide $p - 1 = 10$, so $11$ is impossible.
\item[\textbf{(C)}] $2^2 = 4 \ne 1$; $2$ is the base, not its order.
\item[\textbf{(D)}] $2^5 = 32 \equiv 10 \equiv -1 \pmod{11}$, not $1$; $5$ divides $10$ but is not the order.
\end{itemize}
\textbf{\color{brand}Key takeaway.} Order finding is the hard problem; for tiny $N$ we can list powers, for RSA sizes we cannot.
\end{solbox}
\qsep

\needspace{13\baselineskip}
\begin{qitem}{5}
\qmeta{Single-correct MCQ}{Repeated squaring}
Using repeated squaring, $3^{13} \bmod 17$ equals
\par\smallskip\noindent\begin{tabularx}{\linewidth}{@{}*{4}{L}@{}}
\optlabel{A}\,$5$ & \correctlabel{B}\,$12$ & \optlabel{C}\,$3$ & \optlabel{D}\,$16$
\end{tabularx}
\end{qitem}
\begin{solbox}{B}
\textbf{\color{brand}Correct answer.} \textcolor{okgreen}{\textbf{(B)}} $12$

\textbf{\color{brand}Concepts required.} Repeated squaring computes $a^{2^k}$ by squaring and multiplies the powers selected by the binary digits of the exponent.

\textbf{\color{brand}Relevant formulae.} $13 = 8 + 4 + 1$; $3^2 = 9$, $3^4 = 81 \equiv 13$, $3^8 \equiv 169 \equiv 16$; $3^{13} \equiv 16\cdot13\cdot3 = 624 \equiv 12 \pmod{17}$.

\textbf{\color{brand}Step-by-step explanation.} $3^1 = 3$, $3^2 = 9$, $3^4 = 81 \equiv 13$, $3^8 = 169 \equiv 16 \pmod{17}$. With $13 = 8 + 4 + 1$:
$3^{13} \equiv 16\cdot13\cdot3 = 624 = 36\cdot17 + 12 \equiv 12$. \emph{Check}: $3^{16} \equiv 1$ (Fermat), so $3^{13} \equiv 3^{-3} = 27^{-1} \equiv 10^{-1} \equiv 12$
since $10\cdot12 = 120 = 7\cdot17 + 1$.

\textbf{\color{brand}Why the other options are incorrect.}
\begin{itemize}[leftmargin=2.4em, itemsep=1pt, topsep=2pt]
\item[\textbf{(A)}] $5 \equiv 3^5 \pmod{17}$ ($243 = 14\cdot17 + 5$); it uses the wrong exponent.
\item[\textbf{(C)}] $3 = 3^1$ is the starting value; it would require $3^{12} \equiv 1$, but the order of $3$ modulo $17$ is $16$.
\item[\textbf{(D)}] $16 \equiv 3^8$ is only one of the three factors needed for $13 = 8 + 4 + 1$.
\end{itemize}
\textbf{\color{brand}Key takeaway.} Repeated squaring is the classical model for the controlled-$U^{2^j}$ structure of the quantum circuit.
\end{solbox}
\qsep

\needspace{13\baselineskip}
\begin{qitem}{6}
\qmeta{Single-correct MCQ}{Square-and-multiply cost}
Square-and-multiply computes $a^{1000} \bmod N$ from the binary expansion $1000 = 1111101000_2$. The number of modular multiplications
(squarings plus extra multiplications) is
\par\smallskip\noindent\begin{tabularx}{\linewidth}{@{}*{4}{L}@{}}
\optlabel{A}\,$10$ & \optlabel{B}\,$16$ & \optlabel{C}\,$999$ & \correctlabel{D}\,$14$
\end{tabularx}
\end{qitem}
\begin{solbox}{D}
\textbf{\color{brand}Correct answer.} \textcolor{okgreen}{\textbf{(D)}} $14$

\textbf{\color{brand}Concepts required.} Square-and-multiply performs one squaring per bit after the leading one and one extra multiplication per further $1$-bit.

\textbf{\color{brand}Relevant formulae.} $1000 = 1111101000_2$: $10$ bits $\Rightarrow 9$ squarings; six $1$-bits $\Rightarrow 5$ multiplications; total $14$.

\textbf{\color{brand}Step-by-step explanation.} Left-to-right: start with $a$ (leading bit), then for each of the remaining $9$ bits square once ($9$ squarings) and multiply by $a$ when the bit is
1. The bits after the leading one are $111101000$: five 1s, so $5$ extra multiplications. Total $9 + 5 = 14$. Naively $999$ multiplications would be needed.
This is the classical idea behind computing $a^{2^j} \bmod N$ in Shor's circuit.

\textbf{\color{brand}Why the other options are incorrect.}
\begin{itemize}[leftmargin=2.4em, itemsep=1pt, topsep=2pt]
\item[\textbf{(A)}] $10$ is the number of bits; it ignores the extra multiplications for the $1$-bits (and overcounts the squarings by one).
\item[\textbf{(B)}] $16$ counts a squaring and a multiplication for each of the six $1$-bits plus extras; only $9 + 5$ operations are needed.
\item[\textbf{(C)}] $999$ multiplications is the naive method $a\cdot a\cdots a$, which is exponential in the bit length.
\end{itemize}
\textbf{\color{brand}Key takeaway.} Modular exponentiation is polynomial: $O(n)$ multiplications of $n$-bit numbers, $O(n^3)$ bit operations.
\end{solbox}
\qsep

\needspace{13\baselineskip}
\begin{qitem}{7}
\qmeta{Single-correct MCQ}{Orders divide $\lambda(N)$}
For $N = 35$, which of the following can \emph{not} be the multiplicative order of any $a$ with $\gcd(a, 35) = 1$?
\par\smallskip\noindent\begin{tabularx}{\linewidth}{@{}*{4}{L}@{}}
\optlabel{A}\,$6$ & \correctlabel{B}\,$5$ & \optlabel{C}\,$4$ & \optlabel{D}\,$12$
\end{tabularx}
\end{qitem}
\begin{solbox}{B}
\textbf{\color{brand}Correct answer.} \textcolor{okgreen}{\textbf{(B)}} $5$

\textbf{\color{brand}Concepts required.} Every order divides the Carmichael function $\lambda(N)$, the exponent of the group $\mathbb Z_N^*$.

\textbf{\color{brand}Relevant formulae.} $\lambda(35) = \operatorname{lcm}(\lambda(5), \lambda(7)) = \operatorname{lcm}(4, 6) = 12$; possible orders are the divisors of $12$: $1, 2, 3, 4, 6, 12$.

\textbf{\color{brand}Step-by-step explanation.} By the Chinese remainder theorem the order modulo $35$ is $\operatorname{lcm}(r_5, r_7)$ with $r_5 \mid 4$ and $r_7 \mid 6$, so every order divides the
Carmichael value $\lambda(35) = \operatorname{lcm}(4, 6) = 12$ (and $\varphi(35) = 24$). Orders $12$ (e.g.\ $a = 2$), $4$ (e.g.\ $a = 8$: $8^2 = 64 \equiv 29$,
$8^4 \equiv 1$) and $6$ (e.g.\ $a = 4$: $4^3 = 64 \equiv 29$, $4^6 \equiv 1$) occur; $5$ divides neither $12$ nor $24$.

\textbf{\color{brand}Why the other options are incorrect.}
\begin{itemize}[leftmargin=2.4em, itemsep=1pt, topsep=2pt]
\item[\textbf{(A)}] $6$ divides $12$ and is attained, e.g.\ by $a = 4$ ($4^3 = 64 \equiv 29 \not\equiv 1$, $4^6 \equiv 1$, $4^2 = 16$).
\item[\textbf{(C)}] $4$ divides $12$ and is attained, e.g.\ by $a = 8$ ($8^2 = 64 \equiv 29$, $8^4 \equiv 29^2 = 841 \equiv 1$).
\item[\textbf{(D)}] $12 = \lambda(35)$ is attained, e.g.\ by $a = 3$ ($3^{12} \equiv 1$ and no smaller power is $1$).
\end{itemize}
\textbf{\color{brand}Key takeaway.} Orders divide $\lambda(N)$, which divides $\varphi(N)$; $5$ does not divide $12$.
\end{solbox}
\qsep

\needspace{13\baselineskip}
\begin{qitem}{8}
\qmeta{Single-correct MCQ}{Lucky gcd probability}
For $N = 15$, a base $a$ is drawn uniformly from $\{2, 3, \dots, 13\}$. The probability that $\gcd(a, 15) > 1$, so that a factor is found classically
without any quantum step, is
\par\smallskip\noindent\begin{tabularx}{\linewidth}{@{}*{4}{L}@{}}
\correctlabel{A}\,$1/2$ & \optlabel{B}\,$8/15$ & \optlabel{C}\,$7/15$ & \optlabel{D}\,$1/3$
\end{tabularx}
\end{qitem}
\begin{solbox}{A}
\textbf{\color{brand}Correct answer.} \textcolor{okgreen}{\textbf{(A)}} $1/2$

\textbf{\color{brand}Concepts required.} Before any quantum step, Shor's algorithm computes $\gcd(a, N)$; if it exceeds $1$ a factor is found classically.

\textbf{\color{brand}Relevant formulae.} In $\{2, \dots, 13\}$ (12 values) the bases sharing a factor with $15$ are $3, 5, 6, 9, 10, 12$: $6/12 = 1/2$.

\textbf{\color{brand}Step-by-step explanation.} The multiples of 3 or 5 in $\{2, \dots, 13\}$ are $3, 5, 6, 9, 10, 12$: six of the twelve values, probability $\frac12$. (For RSA-sized $N = pq$ this
probability is about $\frac1p + \frac1q$, astronomically small --- the lucky gcd is only an artefact of tiny examples.)

\textbf{\color{brand}Why the other options are incorrect.}
\begin{itemize}[leftmargin=2.4em, itemsep=1pt, topsep=2pt]
\item[\textbf{(B)}] $8/15 = \varphi(15)/15$ is the fraction of coprime residues in $\{0, \dots, 14\}$, the complement of what is asked, on the wrong set.
\item[\textbf{(C)}] $7/15$ divides by all $15$ residues and counts $0$; the bases are drawn from $12$ values, $6$ of which are lucky.
\item[\textbf{(D)}] $1/3$ counts only the multiples of $3$ ($3, 6, 9, 12$) or of $5$ separately; both kinds must be included.
\end{itemize}
\textbf{\color{brand}Key takeaway.} For tiny $N$ lucky bases are common; for RSA moduli the probability is negligible ($\approx 2/\sqrt N$).
\end{solbox}
\qsep

\needspace{13\baselineskip}
\begin{qitem}{9}
\qmeta{Single-correct MCQ}{Chinese remainder theorem}
The unique $x \in \{0, \dots, 14\}$ with $x \equiv 2 \pmod 3$ and $x \equiv 3 \pmod 5$ is
\par\smallskip\noindent\begin{tabularx}{\linewidth}{@{}*{4}{L}@{}}
\optlabel{A}\,$11$ & \optlabel{B}\,$2$ & \optlabel{C}\,$13$ & \correctlabel{D}\,$8$
\end{tabularx}
\end{qitem}
\begin{solbox}{D}
\textbf{\color{brand}Correct answer.} \textcolor{okgreen}{\textbf{(D)}} $8$

\textbf{\color{brand}Concepts required.} The Chinese remainder theorem: for coprime moduli, the system has a unique solution modulo their product.

\textbf{\color{brand}Relevant formulae.} $x \equiv 3 \pmod 5$: $x \in \{3, 8, 13\}$; of these only $8 \equiv 2 \pmod 3$.

\textbf{\color{brand}Step-by-step explanation.} Candidates with $x \equiv 3 \pmod 5$: $3, 8, 13$; their residues mod 3 are $0, 2, 1$. So $x = 8$. The Chinese remainder theorem guarantees one solution
modulo $15$; it is the reason why orders modulo $N = pq$ are least common multiples of the orders modulo $p$ and $q$.

\textbf{\color{brand}Why the other options are incorrect.}
\begin{itemize}[leftmargin=2.4em, itemsep=1pt, topsep=2pt]
\item[\textbf{(A)}] $11 \equiv 2 \pmod 3$ but $11 \equiv 1 \pmod 5$.
\item[\textbf{(B)}] $2 \equiv 2 \pmod 3$ but $2 \not\equiv 3 \pmod 5$.
\item[\textbf{(C)}] $13 \equiv 3 \pmod 5$ but $13 \equiv 1 \pmod 3$, not $2$.
\end{itemize}
\textbf{\color{brand}Key takeaway.} CRT explains the four square roots of $1$ modulo $pq$: $\pm1$ modulo each prime independently.
\end{solbox}
\qsep

\needspace{13\baselineskip}
\begin{qitem}{10}
\qmeta{Single-correct MCQ}{Square roots of $1$ mod $15$}
The complete set of solutions of $x^2 \equiv 1 \pmod{15}$ with $0 \le x < 15$ is
\par\smallskip\noindent\begin{tabularx}{\linewidth}{@{}*{4}{L}@{}}
\optlabel{A}\,$\{1, 11\}$ & \correctlabel{B}\,$\{1, 4, 11, 14\}$ & \optlabel{C}\,$\{1, 14\}$ & \optlabel{D}\,$\{1, 4, 14\}$
\end{tabularx}
\end{qitem}
\begin{solbox}{B}
\textbf{\color{brand}Correct answer.} \textcolor{okgreen}{\textbf{(B)}} $\{1, 4, 11, 14\}$

\textbf{\color{brand}Concepts required.} Modulo $N = pq$, $x^2 \equiv 1$ iff $x \equiv \pm1$ modulo both $p$ and $q$ (CRT), giving four solutions.

\textbf{\color{brand}Relevant formulae.} $1^2 = 1$, $4^2 = 16 \equiv 1$, $11^2 = 121 \equiv 1$, $14^2 = 196 \equiv 1 \pmod{15}$.

\textbf{\color{brand}Step-by-step explanation.} By CRT, $x^2 \equiv 1$ mod 15 iff $x \equiv \pm1 \pmod 3$ and $x \equiv \pm1 \pmod 5$: four combinations, $x = 1, 4, 11, 14$
($4^2 = 16$, $11^2 = 121 = 8\cdot15 + 1$). The non-trivial roots $4$ and $11$ give factors: $\gcd(4 - 1, 15) = 3$, $\gcd(4 + 1, 15) = 5$. Shor's algorithm
finds such a root as $x = a^{r/2}$.

\textbf{\color{brand}Why the other options are incorrect.}
\begin{itemize}[leftmargin=2.4em, itemsep=1pt, topsep=2pt]
\item[\textbf{(A)}] $\{1, 11\}$ misses $4$ and $14$; $14 \equiv -1$ is always a root.
\item[\textbf{(C)}] $\{1, 14\}$ are only the trivial roots $\pm1$; the non-trivial ones $4$ and $11$ are what Shor's algorithm exploits.
\item[\textbf{(D)}] $\{1, 4, 14\}$ misses $11 = 15 - 4$; roots come in pairs $\pm x$.
\end{itemize}
\textbf{\color{brand}Key takeaway.} A non-trivial square root of $1$ immediately factors $N$ via $\gcd(x \pm 1, N)$.
\end{solbox}
\qsep

\needspace{13\baselineskip}
\begin{qitem}{11}
\qmeta{Single-correct MCQ}{Failure: $a^{r/2} \equiv -1$}
For $N = 21$ and $a = 5$, the order $r$ and the outcome of Shor's classical reduction are
\begin{itemize}[leftmargin=2.6em, labelsep=0.4em, itemsep=2pt, topsep=3pt]
\item[\correctlabel{A}] $r = 6$; the reduction fails because $5^3 \equiv -1 \pmod{21}$
\item[\optlabel{B}] $r = 3$; it fails because $r$ is odd
\item[\optlabel{C}] $r = 12$; it gives the factors $3$ and $7$
\item[\optlabel{D}] $r = 6$; it gives the factors $3$ and $7$
\end{itemize}
\end{qitem}
\begin{solbox}{A}
\textbf{\color{brand}Correct answer.} \textcolor{okgreen}{\textbf{(A)}} $r = 6$; the reduction fails because $5^3 \equiv -1 \pmod{21}$

\textbf{\color{brand}Concepts required.} The reduction needs $r$ even and $a^{r/2} \not\equiv -1$; otherwise the attempt fails and a new base is chosen.

\textbf{\color{brand}Relevant formulae.} Powers of $5$ mod $21$: $5, 4, 20, 16, 17, 1$ $\Rightarrow r = 6$; $5^3 \equiv 20 \equiv -1$.

\textbf{\color{brand}Step-by-step explanation.} $5^1 = 5$, $5^2 = 25 \equiv 4$, $5^3 \equiv 20 \equiv -1$, $5^6 \equiv 1$: $r = 6$ (even). But $a^{r/2} = 5^3 \equiv -1$, so
$\gcd(5^3 + 1, 21) = \gcd(21, 21) = 21$ and $\gcd(5^3 - 1, 21) = \gcd(19, 21) = 1$: only trivial factors. The algorithm must restart with another $a$.

\textbf{\color{brand}Why the other options are incorrect.}
\begin{itemize}[leftmargin=2.4em, itemsep=1pt, topsep=2pt]
\item[\textbf{(B)}] $5^3 \equiv 20 \ne 1$, so $3$ is not the order; the order is $6$, which is even.
\item[\textbf{(C)}] $12$ is a multiple of the order, not the order; $5^6 \equiv 1$ already.
\item[\textbf{(D)}] With $5^3 \equiv -1$ the gcds are $\gcd(19, 21) = 1$ and $\gcd(21, 21) = 21$: trivial, so no factors are obtained.
\end{itemize}
\textbf{\color{brand}Key takeaway.} Correct order finding can still give a failed attempt; report it as a failure and retry.
\end{solbox}
\qsep

\needspace{13\baselineskip}
\begin{qitem}{12}
\qmeta{Single-correct MCQ}{Fermat's little theorem}
$2^{100} \bmod 13$ equals
\par\smallskip\noindent\begin{tabularx}{\linewidth}{@{}*{4}{L}@{}}
\optlabel{A}\,$9$ & \optlabel{B}\,$1$ & \correctlabel{C}\,$3$ & \optlabel{D}\,$12$
\end{tabularx}
\end{qitem}
\begin{solbox}{C}
\textbf{\color{brand}Correct answer.} \textcolor{okgreen}{\textbf{(C)}} $3$

\textbf{\color{brand}Concepts required.} Fermat's little theorem: $a^{p-1} \equiv 1 \pmod p$ for prime $p \nmid a$; reduce the exponent modulo the order.

\textbf{\color{brand}Relevant formulae.} $2^{12} \equiv 1 \pmod{13}$, $100 = 8\cdot12 + 4$, so $2^{100} \equiv 2^4 = 16 \equiv 3$.

\textbf{\color{brand}Step-by-step explanation.} Fermat: $2^{12} \equiv 1 \pmod{13}$. $100 = 8\cdot12 + 4$, so $2^{100} \equiv 2^4 = 16 \equiv 3$. Reducing the exponent modulo the order (or
$\varphi$) is what makes $a^x \bmod N$ periodic in $x$.

\textbf{\color{brand}Why the other options are incorrect.}
\begin{itemize}[leftmargin=2.4em, itemsep=1pt, topsep=2pt]
\item[\textbf{(A)}] $9 \equiv 2^8 \pmod{13}$ ($256 = 19\cdot13 + 9$): it reduces the exponent incorrectly; $100 \bmod 12 = 4$, not $8$.
\item[\textbf{(B)}] $1$ would require $12 \mid 100$; $100 = 8\cdot12 + 4$.
\item[\textbf{(D)}] $12 \equiv -1$ is $2^6$; the exponent $100$ reduces to $4$, not $6$.
\end{itemize}
\textbf{\color{brand}Key takeaway.} Exponents can be reduced modulo the order --- which is exactly why the order is so useful.
\end{solbox}
\qsep

\needspace{13\baselineskip}
\begin{qitem}{13}
\qmeta{Single-correct MCQ}{Trial-division bound}
Trial division factors $N = 10403$ by testing divisors $d \ge 2$ in increasing order. The largest $d$ that may have to be tested is
\par\smallskip\noindent\begin{tabularx}{\linewidth}{@{}*{4}{L}@{}}
\optlabel{A}\,$5201$ & \optlabel{B}\,$102$ & \optlabel{C}\,$103$ & \correctlabel{D}\,$101$
\end{tabularx}
\end{qitem}
\begin{solbox}{D}
\textbf{\color{brand}Correct answer.} \textcolor{okgreen}{\textbf{(D)}} $101$

\textbf{\color{brand}Concepts required.} A composite $N$ has a prime factor $\le \sqrt N$, so trial division stops at $\lfloor\sqrt N\rfloor$.

\textbf{\color{brand}Relevant formulae.} $\sqrt{10403} = 101.99$; $10403 = 101\cdot103$, so $d = 101$ is reached and divides.

\textbf{\color{brand}Step-by-step explanation.} If $N$ is composite it has a factor $\le\sqrt N$. $101^2 = 10201 \le 10403 < 102^2 = 10404$, so testing up to $\lfloor\sqrt N\rfloor = 101$ suffices
--- and $10403 = 101\times103$ is found at $d = 101$. For an $n$-bit $N$ this is about $2^{n/2}$ divisions: exponential in $n$.

\textbf{\color{brand}Why the other options are incorrect.}
\begin{itemize}[leftmargin=2.4em, itemsep=1pt, topsep=2pt]
\item[\textbf{(A)}] $5201 \approx N/2$ is a far too pessimistic bound; divisors above $\sqrt N$ pair with ones below it.
\item[\textbf{(B)}] $102 > \sqrt{10403}$ is never needed (and is even); the search stops at $101$.
\item[\textbf{(C)}] $103$ is the larger factor; it is found as $N/101$, never by testing.
\end{itemize}
\textbf{\color{brand}Key takeaway.} Trial division costs $O(\sqrt N) = O(2^{n/2})$ --- exponential in the bit length.
\end{solbox}
\qsep

\needspace{13\baselineskip}
\begin{qitem}{14}
\qmeta{Single-correct MCQ}{Success bound $1 - 2^{-(k-1)}$}
For an odd $N$ with $k$ distinct prime factors (not a prime power) and $a$ uniformly random among the residues coprime to $N$, the probability
that $r$ is even and $a^{r/2} \not\equiv -1 \pmod N$ is at least $1 - 2^{-(k-1)}$. For $N = 105 = 3\cdot5\cdot7$ this lower bound is
\par\smallskip\noindent\begin{tabularx}{\linewidth}{@{}*{4}{L}@{}}
\optlabel{A}\,$1/8$ & \optlabel{B}\,$1/2$ & \correctlabel{C}\,$3/4$ & \optlabel{D}\,$7/8$
\end{tabularx}
\end{qitem}
\begin{solbox}{C}
\textbf{\color{brand}Correct answer.} \textcolor{okgreen}{\textbf{(C)}} $3/4$

\textbf{\color{brand}Concepts required.} For $N$ with $k$ distinct odd prime factors, a random coprime base works with probability at least $1 - 2^{-(k-1)}$.

\textbf{\color{brand}Relevant formulae.} $k = 3$: $1 - 2^{-2} = 3/4$.

\textbf{\color{brand}Step-by-step explanation.} $k = 3$: $1 - 2^{-2} = \frac34$. For RSA moduli ($k = 2$) the bound is $\frac12$, so on average at most two bases are needed. This bound comes from
the CRT: the reduction fails only if the 2-adic parts of the orders modulo the prime powers all agree.

\textbf{\color{brand}Why the other options are incorrect.}
\begin{itemize}[leftmargin=2.4em, itemsep=1pt, topsep=2pt]
\item[\textbf{(A)}] $1/8 = 2^{-3}$ confuses the failure bound with the success bound and uses the wrong exponent.
\item[\textbf{(B)}] $1/2$ is the bound for $k = 2$ (an RSA modulus), not for three prime factors.
\item[\textbf{(D)}] $7/8 = 1 - 2^{-3}$ uses the exponent $k$ instead of $k - 1$.
\end{itemize}
\textbf{\color{brand}Key takeaway.} For an RSA modulus at least half of the bases work, so a few repetitions suffice.
\end{solbox}
\qsep

\needspace{13\baselineskip}
\begin{qitem}{15}
\qmeta{Matching lists}{$\varphi$, order, $\lambda$, inverse}
Match each quantity for $N = 15$ (List-I) with its value (List-II).
\par\smallskip\noindent{\renewcommand{\arraystretch}{1.35}%
\begin{tabularx}{\linewidth}{@{}c>{\raggedright\arraybackslash}p{0.33\linewidth}@{\hspace{1.2em}}cL@{}}
\toprule
\multicolumn{2}{@{}l}{\textbf{List-I}} & \multicolumn{2}{l@{}}{\textbf{List-II}} \\ \midrule
\textbf{(P)} & $\varphi(15)$ & \textbf{(1)} & $13$ \\
\textbf{(Q)} & order of $4$ modulo $15$ & \textbf{(2)} & $8$ \\
\textbf{(R)} & Carmichael value $\lambda(15) = \operatorname{lcm}(\lambda(3), \lambda(5))$ & \textbf{(3)} & $2$ \\
\textbf{(S)} & $7^{-1} \bmod 15$ & \textbf{(4)} & $4$ \\
\bottomrule
\end{tabularx}}
\par\smallskip\noindent\begin{tabularx}{\linewidth}{@{}*{2}{L}@{}}
\optlabel{A}\,P$\to$4, Q$\to$3, R$\to$2, S$\to$1 & \optlabel{B}\,P$\to$2, Q$\to$3, R$\to$1, S$\to$4 \\[3pt]
\correctlabel{C}\,P$\to$2, Q$\to$3, R$\to$4, S$\to$1 & \optlabel{D}\,P$\to$2, Q$\to$4, R$\to$3, S$\to$1
\end{tabularx}
\end{qitem}
\begin{solbox}{C}
\textbf{\color{brand}Correct answer.} \textcolor{okgreen}{\textbf{(C)}} P$\to$2, Q$\to$3, R$\to$4, S$\to$1

\textbf{\color{brand}Concepts required.} $\varphi$, $\lambda$, orders and inverses are basic invariants of $\mathbb Z_N^*$; for $N = 15$ they can be computed by hand.

\textbf{\color{brand}Relevant formulae.} $\varphi(15) = 2\cdot4 = 8$; $4^2 = 16 \equiv 1$ so $\operatorname{ord}(4) = 2$; $\lambda(15) = \operatorname{lcm}(2, 4) = 4$; $7\cdot13 = 91 \equiv 1$.

\textbf{\color{brand}Step-by-step explanation.} \begin{itemize}\itemsep0pt
\item $\varphi(15) = 2\cdot4 = 8$ (2).
\item $4^2 = 16 \equiv 1$: order $2$ (3).
\item $\lambda(15) = \operatorname{lcm}(2, 4) = 4$ (4): every order divides it.
\item $7\cdot13 = 91 = 6\cdot15 + 1$: $7^{-1} = 13$ (1).
\end{itemize}
Hence P$\to$2, Q$\to$3, R$\to$4, S$\to$1.

\textbf{\color{brand}Why the other options are incorrect.}
\begin{itemize}[leftmargin=2.4em, itemsep=1pt, topsep=2pt]
\item[\textbf{(A)}] This assigns $4$ to $\varphi(15)$; $\varphi(15) = 8$ counts $1, 2, 4, 7, 8, 11, 13, 14$.
\item[\textbf{(B)}] This gives $\lambda(15) = 13$ and $7^{-1} = 4$; but $7\cdot4 = 28 \equiv 13 \not\equiv 1$.
\item[\textbf{(D)}] This swaps the order of $4$ and $\lambda(15)$: $4^2 \equiv 1$, so the order is $2$, while $\lambda(15) = 4$.
\end{itemize}
\textbf{\color{brand}Key takeaway.} Orders divide $\lambda(N)$, which divides $\varphi(N)$: here $2 \mid 4 \mid 8$.
\end{solbox}
\qsep

\needspace{13\baselineskip}
\begin{qitem}{16}
\qmeta{Matching lists}{Classical factoring complexity}
Match each classical algorithm (List-I) with its running time for an $n$-bit input $N$ (List-II).
\par\smallskip\noindent{\renewcommand{\arraystretch}{1.35}%
\begin{tabularx}{\linewidth}{@{}c>{\raggedright\arraybackslash}p{0.33\linewidth}@{\hspace{1.2em}}cL@{}}
\toprule
\multicolumn{2}{@{}l}{\textbf{List-I}} & \multicolumn{2}{l@{}}{\textbf{List-II}} \\ \midrule
\textbf{(P)} & trial division & \textbf{(1)} & $\exp\big((64/9)^{1/3}(\ln N)^{1/3}(\ln\ln N)^{2/3}(1 + o(1))\big)$ (sub-exponential) \\
\textbf{(Q)} & Pollard's rho & \textbf{(2)} & $O(N^{1/4})$ expected, i.e.\ $O(2^{n/4})$ (heuristic) \\
\textbf{(R)} & general number field sieve & \textbf{(3)} & $O(n)$ division steps (polynomial) \\
\textbf{(S)} & Euclidean $\gcd(a, N)$ & \textbf{(4)} & $O(\sqrt N) = O(2^{n/2})$ divisions \\
\bottomrule
\end{tabularx}}
\par\smallskip\noindent\begin{tabularx}{\linewidth}{@{}*{2}{L}@{}}
\optlabel{A}\,P$\to$2, Q$\to$4, R$\to$1, S$\to$3 & \optlabel{B}\,P$\to$4, Q$\to$1, R$\to$2, S$\to$3 \\[3pt]
\optlabel{C}\,P$\to$4, Q$\to$2, R$\to$3, S$\to$1 & \correctlabel{D}\,P$\to$4, Q$\to$2, R$\to$1, S$\to$3
\end{tabularx}
\end{qitem}
\begin{solbox}{D}
\textbf{\color{brand}Correct answer.} \textcolor{okgreen}{\textbf{(D)}} P$\to$4, Q$\to$2, R$\to$1, S$\to$3

\textbf{\color{brand}Concepts required.} Classical factoring methods range from exponential (trial division, rho) to sub-exponential (GNFS); gcd is polynomial.

\textbf{\color{brand}Relevant formulae.} Trial division $O(2^{n/2})$; Pollard rho $O(2^{n/4})$; GNFS $L_N[1/3, (64/9)^{1/3}]$; Euclid $O(n)$ steps.

\textbf{\color{brand}Step-by-step explanation.} \begin{itemize}\itemsep0pt
\item trial division tests divisors up to $\sqrt N$ (4);
\item Pollard's rho finds a factor $p$ in about $\sqrt p \le N^{1/4}$ steps (2);
\item the GNFS is the fastest known general method, sub-exponential but super-polynomial (1);
\item the gcd is polynomial (3) --- not a factoring algorithm on its own.
\end{itemize}
Hence P$\to$4, Q$\to$2, R$\to$1, S$\to$3. Not all classical factoring is exponential.

\textbf{\color{brand}Why the other options are incorrect.}
\begin{itemize}[leftmargin=2.4em, itemsep=1pt, topsep=2pt]
\item[\textbf{(A)}] This swaps trial division and Pollard's rho; rho finds a factor $p$ in about $\sqrt p \le N^{1/4}$ steps, faster than trial division.
\item[\textbf{(B)}] This gives Pollard's rho the GNFS running time and vice versa; GNFS is the fastest known general method.
\item[\textbf{(C)}] This gives the gcd the GNFS time; the Euclidean algorithm is polynomial, $O(n)$ division steps.
\end{itemize}
\textbf{\color{brand}Key takeaway.} The best classical method is sub-exponential; Shor's algorithm is polynomial --- on a fault-tolerant machine.
\end{solbox}
\qsep

\needspace{13\baselineskip}
\begin{qitem}{17}
\qmeta{Assertion--Reason}{Non-trivial square root $\Rightarrow$ factor}
\textbf{Assertion (A):} If $x^2 \equiv 1 \pmod N$ but $x \not\equiv \pm1 \pmod N$, then $\gcd(x - 1, N)$ is a non-trivial factor of $N$.\\[1mm]
\textbf{Reason (R):} In that case $N$ divides $(x - 1)(x + 1)$ but divides neither $x - 1$ nor $x + 1$.
\begin{itemize}[leftmargin=2.6em, labelsep=0.4em, itemsep=2pt, topsep=3pt]
\item[\correctlabel{A}] Both Assertion (A) and Reason (R) are true, and R is the correct explanation of A
\item[\optlabel{B}] Both A and R are true, but R is \emph{not} the correct explanation of A
\item[\optlabel{C}] A is true, but R is false
\item[\optlabel{D}] A is false, but R is true
\end{itemize}
\end{qitem}
\begin{solbox}{A}
\textbf{\color{brand}Correct answer.} \textcolor{okgreen}{\textbf{(A)}} Both Assertion (A) and Reason (R) are true, and R is the correct explanation of A

\textbf{\color{brand}Concepts required.} If $N \mid (x - 1)(x + 1)$ but $N$ divides neither factor, each prime factor of $N$ is split between them.

\textbf{\color{brand}Relevant formulae.} $x^2 \equiv 1 \Rightarrow N \mid (x - 1)(x + 1)$; $x \not\equiv \pm1 \Rightarrow 1 < \gcd(x \pm 1, N) < N$.

\textbf{\color{brand}Step-by-step explanation.} \textbf{R is true}: $x^2 - 1 = (x-1)(x+1) \equiv 0$, while $x \not\equiv \pm1$ means $N \nmid x \mp 1$. \textbf{A follows from R}: if $\gcd(x - 1, N)$ were
$1$, then $N \mid (x-1)(x+1)$ would force $N \mid x + 1$; if it were $N$, then $N \mid x - 1$ --- both excluded. So it is a proper divisor. This is exactly
Shor's factor-extraction step with $x = a^{r/2}$. Answer: (A).

\textbf{\color{brand}Why the other options are incorrect.}
\begin{itemize}[leftmargin=2.4em, itemsep=1pt, topsep=2pt]
\item[\textbf{(B)}] is wrong because R is exactly the argument that proves A: a divisor of a product that divides neither factor must share a proper factor with each.
\item[\textbf{(C)}] is wrong because R is true: $x^2 - 1 = (x - 1)(x + 1) \equiv 0$, and $x \not\equiv \pm1$ means $N$ divides neither factor.
\item[\textbf{(D)}] is wrong because A is true: e.g.\ $x = 4$, $N = 15$ gives $\gcd(3, 15) = 3$.
\end{itemize}
\textbf{\color{brand}Key takeaway.} This is the heart of the reduction: $x = a^{r/2}$ is a square root of $1$.
\end{solbox}
\qsep

\needspace{13\baselineskip}
\begin{qitem}{18}
\qmeta{Statement I / II}{Euler's theorem vs order}
\textbf{Statement I:} If $\gcd(a, N) = 1$, then $a^{\varphi(N)} \equiv 1 \pmod N$.\\[1mm]
\textbf{Statement II:} The multiplicative order of every $a$ coprime to $N$ equals $\varphi(N)$.
\begin{itemize}[leftmargin=2.6em, labelsep=0.4em, itemsep=2pt, topsep=3pt]
\item[\optlabel{A}] Both Statement I and Statement II are true
\item[\optlabel{B}] Both Statement I and Statement II are false
\item[\correctlabel{C}] Statement I is true, but Statement II is false
\item[\optlabel{D}] Statement I is false, but Statement II is true
\end{itemize}
\end{qitem}
\begin{solbox}{C}
\textbf{\color{brand}Correct answer.} \textcolor{okgreen}{\textbf{(C)}} Statement I is true, but Statement II is false

\textbf{\color{brand}Concepts required.} Euler's theorem gives an exponent that always works; the order is the smallest one and is usually a proper divisor of $\varphi(N)$.

\textbf{\color{brand}Relevant formulae.} $a^{\varphi(N)} \equiv 1$; $\operatorname{ord}(a) \mid \lambda(N) \mid \varphi(N)$; e.g.\ $\operatorname{ord}_{15}(4) = 2 \ne 8$.

\textbf{\color{brand}Step-by-step explanation.} \textbf{I is true} (Euler's theorem). \textbf{II is false}: the order only \emph{divides} $\varphi(N)$; e.g.\ $4$ has order $2$ modulo $15$ while
$\varphi(15) = 8$. Finding the actual order is the hard problem that Shor's quantum period finding solves. Answer: (C).

\textbf{\color{brand}Why the other options are incorrect.}
\begin{itemize}[leftmargin=2.4em, itemsep=1pt, topsep=2pt]
\item[\textbf{(A)}] is wrong because Statement II is false: $\operatorname{ord}_{15}(4) = 2$ while $\varphi(15) = 8$.
\item[\textbf{(B)}] is wrong because Statement I is Euler's theorem, which holds for every $a$ coprime to $N$.
\item[\textbf{(D)}] is wrong because Statement I is true and Statement II false, not the other way round.
\end{itemize}
\textbf{\color{brand}Key takeaway.} If orders were always $\varphi(N)$, knowing them would directly reveal $\varphi(N)$; in reality they only divide it.
\end{solbox}
\qsep

\needspace{20\baselineskip}
\begin{passagebox}{Paragraph for Questions 19 and 20}
Shor's reduction is applied by hand to $N = 91 = 7\times13$ with the base $a = 3$ (note $\gcd(3, 91) = 1$). Recall that the
order $r$ of $a$ modulo $N$ is the least $r > 0$ with $a^r \equiv 1 \pmod N$, and that if $r$ is even and $a^{r/2} \not\equiv -1 \pmod N$,
then $\gcd(a^{r/2} \pm 1, N)$ are non-trivial factors.
\end{passagebox}
\needspace{13\baselineskip}
\begin{qitem}{19}
\qmeta{Paragraph based}{Order of $3$ mod $91$}
The order $r$ of $3$ modulo $91$ is
\par\smallskip\noindent\begin{tabularx}{\linewidth}{@{}*{4}{L}@{}}
\optlabel{A}\,$72$ & \correctlabel{B}\,$6$ & \optlabel{C}\,$12$ & \optlabel{D}\,$3$
\end{tabularx}
\end{qitem}
\begin{solbox}{B}
\textbf{\color{brand}Correct answer.} \textcolor{okgreen}{\textbf{(B)}} $6$

\textbf{\color{brand}Concepts required.} The order is the first exponent where the powers return to $1$.

\textbf{\color{brand}Relevant formulae.} $3, 9, 27, 81, 243 \equiv 61, 183 \equiv 1 \pmod{91}$ ($729 = 8\cdot91 + 1$): $r = 6$.

\textbf{\color{brand}Step-by-step explanation.} By CRT, $r = \operatorname{lcm}(r_7, r_{13})$. Mod 7: $3, 2, 6, 4, 5, 1$: order $6$. Mod 13: $3, 9, 27 \equiv 1$: order $3$. So $r = \operatorname{lcm}(6, 3) = 6$.
\emph{Check}: $3^6 = 729 = 8\cdot91 + 1$. ($72 = \varphi(91)$ is only a multiple of the order.)

\textbf{\color{brand}Why the other options are incorrect.}
\begin{itemize}[leftmargin=2.4em, itemsep=1pt, topsep=2pt]
\item[\textbf{(A)}] $72 = \varphi(91)$ is an exponent that works, not the least one.
\item[\textbf{(C)}] $12$ is a multiple of the order ($3^{12} \equiv 1$) but not the least one.
\item[\textbf{(D)}] $3^3 = 27 \ne 1$; $3$ is half of the order.
\end{itemize}
\textbf{\color{brand}Key takeaway.} Order $=$ smallest exponent; multiples and $\varphi(N)$ also satisfy $a^x \equiv 1$.
\end{solbox}
\qsep

\needspace{13\baselineskip}
\begin{qitem}{20}
\qmeta{Paragraph based}{Factors of $91$ by gcd}
The values of $\gcd(3^{r/2} - 1, 91)$ and $\gcd(3^{r/2} + 1, 91)$ are, respectively,
\par\smallskip\noindent\begin{tabularx}{\linewidth}{@{}*{4}{L}@{}}
\optlabel{A}\,$26$ and $28$ & \optlabel{B}\,$1$ and $91$ & \correctlabel{C}\,$13$ and $7$ & \optlabel{D}\,$7$ and $13$
\end{tabularx}
\end{qitem}
\begin{solbox}{C}
\textbf{\color{brand}Correct answer.} \textcolor{okgreen}{\textbf{(C)}} $13$ and $7$

\textbf{\color{brand}Concepts required.} With $r$ even and $a^{r/2} \not\equiv -1$, the gcds with $a^{r/2} \mp 1$ give the factors.

\textbf{\color{brand}Relevant formulae.} $3^{r/2} = 27$; $\gcd(26, 91) = 13$, $\gcd(28, 91) = 7$; $7\cdot13 = 91$.

\textbf{\color{brand}Step-by-step explanation.} $r = 6$ is even and $3^3 = 27 \not\equiv -1 \equiv 90$. $\gcd(26, 91) = 13$ (since $91 = 7\cdot13$, $26 = 2\cdot13$) and $\gcd(28, 91) = 7$
($28 = 4\cdot7$). Both factors are found: $91 = 7\times13$. The order of the two numbers matters: $a^{r/2} - 1 = 26$ shares $13$ with $91$.

\textbf{\color{brand}Why the other options are incorrect.}
\begin{itemize}[leftmargin=2.4em, itemsep=1pt, topsep=2pt]
\item[\textbf{(A)}] $26$ and $28$ are $a^{r/2} \mp 1$ themselves; the gcd with $91$ is still needed.
\item[\textbf{(B)}] $1$ and $91$ would be the trivial result for $a^{r/2} \equiv -1$; but $27 \not\equiv 90$.
\item[\textbf{(D)}] $7$ and $13$ are the correct factors in the wrong order: $\gcd(26, 91) = 13$ and $\gcd(28, 91) = 7$.
\end{itemize}
\textbf{\color{brand}Key takeaway.} The quantum part ends with $r$; the factors come from two classical gcds.
\end{solbox}
\qsep

\newpage

\begin{sectionbanner}SECTION B \enspace$\cdot$\enspace Quantum Computing Foundations \hfill Q.21--40\end{sectionbanner}

\needspace{13\baselineskip}
\begin{qitem}{21}
\qmeta{Single-correct MCQ}{Uniform superposition}
After $H^{\otimes8}$ acts on $\ket{0}^{\otimes8}$ (the counting register of the notebook for $N = 15$), the amplitude of each basis state $\ket x$ is
\par\smallskip\noindent\begin{tabularx}{\linewidth}{@{}*{4}{L}@{}}
\optlabel{A}\,$1/256$ & \correctlabel{B}\,$1/16$ & \optlabel{C}\,$1/8$ & \optlabel{D}\,$1/\sqrt8$
\end{tabularx}
\end{qitem}
\begin{solbox}{B}
\textbf{\color{brand}Correct answer.} \textcolor{okgreen}{\textbf{(B)}} $1/16$

\textbf{\color{brand}Concepts required.} The Hadamard transform of $\ket0^{\otimes t}$ is the uniform superposition of all $2^t$ basis states.

\textbf{\color{brand}Relevant formulae.} $H^{\otimes t}\ket0^{\otimes t} = 2^{-t/2}\sum_x\ket x$; $t = 8$: amplitude $1/16$, probability $1/256$.

\textbf{\color{brand}Step-by-step explanation.} $H^{\otimes t}\ket{0}^{\otimes t} = \frac{1}{\sqrt{2^t}}\sum_{x=0}^{2^t-1}\ket x$. With $t = 8$: $\frac{1}{\sqrt{256}} = \frac1{16}$. The probabilities are $\frac{1}{256}$ each.
This superposition is the input on which the modular exponentiation acts ``for all $x$ at once''.

\textbf{\color{brand}Why the other options are incorrect.}
\begin{itemize}[leftmargin=2.4em, itemsep=1pt, topsep=2pt]
\item[\textbf{(A)}] $1/256$ is the probability $|1/16|^2$, not the amplitude.
\item[\textbf{(C)}] $1/8$ would be the amplitude for $t = 6$ ($2^{-3}$); for $t = 8$ it is $2^{-4}$.
\item[\textbf{(D)}] $1/\sqrt8$ is the amplitude for $t = 3$ qubits.
\end{itemize}
\textbf{\color{brand}Key takeaway.} Amplitudes are $2^{-t/2}$; probabilities $2^{-t}$.
\end{solbox}
\qsep

\needspace{13\baselineskip}
\begin{qitem}{22}
\qmeta{Single-correct MCQ}{Measuring the work register}
The state $\frac{1}{\sqrt8}\sum_{x=0}^{7}\ket x\ket{2^x \bmod 15}$ is prepared and the \emph{second} register is measured with result $4$. The probability
of this result and the post-measurement state of the first register are
\par\smallskip\noindent\begin{tabularx}{\linewidth}{@{}*{2}{L}@{}}
\optlabel{A}\,$1/2$; $\ \frac{1}{\sqrt2}(\ket2 + \ket6)$ & \optlabel{B}\,$1/4$; $\ \frac{1}{\sqrt2}(\ket2 - \ket6)$ \\[3pt]
\optlabel{C}\,$1/8$; $\ \ket2$ & \correctlabel{D}\,$1/4$; $\ \frac{1}{\sqrt2}(\ket2 + \ket6)$
\end{tabularx}
\end{qitem}
\begin{solbox}{D}
\textbf{\color{brand}Correct answer.} \textcolor{okgreen}{\textbf{(D)}} $1/4$; $\ \frac{1}{\sqrt2}(\ket2 + \ket6)$

\textbf{\color{brand}Concepts required.} Measuring one register of an entangled state keeps only the compatible terms of the other, renormalised.

\textbf{\color{brand}Relevant formulae.} $2^x \bmod 15$: $1, 2, 4, 8, 1, 2, 4, 8$; value $4$ at $x = 2, 6$: $P = 2/8$, state $(\ket2 + \ket6)/\sqrt2$.

\textbf{\color{brand}Step-by-step explanation.} $2^x \bmod 15$ for $x = 0, \dots, 7$: $1, 2, 4, 8, 1, 2, 4, 8$. The value $4$ occurs for $x = 2$ and $x = 6$, so
$P = 2\times\frac18 = \frac14$, and the first register collapses to $\frac{1}{\sqrt2}(\ket2 + \ket6)$ (equal positive amplitudes, renormalised). It is periodic
with period $r = 4$ and an unknown offset --- the QFT removes the offset from the probabilities.

\textbf{\color{brand}Why the other options are incorrect.}
\begin{itemize}[leftmargin=2.4em, itemsep=1pt, topsep=2pt]
\item[\textbf{(A)}] $1/2$ counts two of the four values as matching; only $x = 2, 6$ of $8$ terms give $4$.
\item[\textbf{(B)}] The amplitudes before the measurement are all $+1/\sqrt8$; no minus sign can appear.
\item[\textbf{(C)}] $1/8$ and $\ket2$ forget that the function is periodic: $x = 6$ gives $2^6 = 64 \equiv 4$ as well.
\end{itemize}
\textbf{\color{brand}Key takeaway.} The collapsed first register is periodic with period $r$ and a random offset.
\end{solbox}
\qsep

\needspace{13\baselineskip}
\begin{qitem}{23}
\qmeta{Single-correct MCQ}{Unitarity of $U_a$}
Why does the operator $U_a\ket y = \ket{ay \bmod N}$ (for $0 \le y < N$; identity for $y \ge N$) require $\gcd(a, N) = 1$ to be a valid quantum gate?
\begin{itemize}[leftmargin=2.6em, labelsep=0.4em, itemsep=2pt, topsep=3pt]
\item[\optlabel{A}] because the eigenvalues of $U_a$ must be real
\item[\optlabel{B}] because a quantum gate must map $\ket0$ to $\ket0$
\item[\correctlabel{C}] only then is $y \mapsto ay \bmod N$ a permutation of $\{0, \dots, N-1\}$, so $U_a$ is a permutation matrix and hence unitary
\item[\optlabel{D}] because otherwise $U_a$ would not be Hermitian
\end{itemize}
\end{qitem}
\begin{solbox}{C}
\textbf{\color{brand}Correct answer.} \textcolor{okgreen}{\textbf{(C)}} only then is $y \mapsto ay \bmod N$ a permutation of $\{0, \dots, N-1\}$, so $U_a$ is a permutation matrix and hence unitary

\textbf{\color{brand}Concepts required.} A permutation matrix is unitary; a non-injective map of basis states is not.

\textbf{\color{brand}Relevant formulae.} $\gcd(a, N) = 1 \iff a$ invertible mod $N \iff y \mapsto ay \bmod N$ bijective on $\{0, \dots, N-1\}$.

\textbf{\color{brand}Step-by-step explanation.} If $\gcd(a, N) = d > 1$, two inputs collide: e.g.\ $a = 3$, $N = 15$: $3\cdot0 \equiv 3\cdot5 \equiv 0$. A non-injective map cannot be unitary.
For $\gcd(a, N) = 1$, $a$ is invertible modulo $N$, the map is a bijection, and a permutation matrix satisfies $U^\dagger U = I$. Gates need not be
Hermitian, and their eigenvalues are complex phases.

\textbf{\color{brand}Why the other options are incorrect.}
\begin{itemize}[leftmargin=2.4em, itemsep=1pt, topsep=2pt]
\item[\textbf{(A)}] Unitary eigenvalues are phases $e^{i\phi}$; $U_a$ has complex eigenvalues $e^{2\pi is/r}$.
\item[\textbf{(B)}] Gates need not fix $\ket0$; $U_a$ happens to fix $\ket0$, but that is not the requirement.
\item[\textbf{(D)}] Quantum gates must be unitary, not Hermitian; $U_a$ is generally not Hermitian (it cycles states).
\end{itemize}
\textbf{\color{brand}Key takeaway.} Coprimality makes modular multiplication reversible, hence implementable as a quantum gate.
\end{solbox}
\qsep

\needspace{13\baselineskip}
\begin{qitem}{24}
\qmeta{Single-correct MCQ}{Eigenvalues of $U_a$}
On the subspace spanned by $\{\ket1, \ket2, \ket4, \ket8\}$, the operator $U_2\ket y = \ket{2y \bmod 15}$ has the eigenvalues
\par\smallskip\noindent\begin{tabularx}{\linewidth}{@{}*{2}{L}@{}}
\optlabel{A}\,$1, 2, 4, 8$ & \optlabel{B}\,$1, -1$ (each twice) \\[3pt]
\optlabel{C}\,$e^{2\pi i/15}, e^{4\pi i/15}, e^{8\pi i/15}, e^{16\pi i/15}$ & \correctlabel{D}\,$1, i, -1, -i$
\end{tabularx}
\end{qitem}
\begin{solbox}{D}
\textbf{\color{brand}Correct answer.} \textcolor{okgreen}{\textbf{(D)}} $1, i, -1, -i$

\textbf{\color{brand}Concepts required.} A cyclic permutation of length $r$ has the $r$-th roots of unity as eigenvalues.

\textbf{\color{brand}Relevant formulae.} $U_2$: $1 \to 2 \to 4 \to 8 \to 1$, $r = 4$; eigenvalues $e^{2\pi is/4} = 1, i, -1, -i$.

\textbf{\color{brand}Step-by-step explanation.} $U_2$ cycles $\ket1 \to \ket2 \to \ket4 \to \ket8 \to \ket1$: a cyclic shift of length $r = 4$, so $U_2^4 = I$ there and the eigenvalues are the fourth
roots of unity $e^{2\pi is/4}$, $s = 0, 1, 2, 3$. In general the eigenvalues on the orbit of $\ket1$ are $e^{2\pi is/r}$: the order $r$ is encoded in the
eigenphases $s/r$ --- this is what phase estimation reads out.

\textbf{\color{brand}Why the other options are incorrect.}
\begin{itemize}[leftmargin=2.4em, itemsep=1pt, topsep=2pt]
\item[\textbf{(A)}] $1, 2, 4, 8$ are the basis labels on the orbit, not eigenvalues; unitary eigenvalues have modulus $1$.
\item[\textbf{(B)}] $\pm1$ would require $U_2^2 = I$, but $U_2^2\ket1 = \ket4 \ne \ket1$.
\item[\textbf{(C)}] Phases with denominator $15$ use $N$ instead of the order $r = 4$.
\end{itemize}
\textbf{\color{brand}Key takeaway.} The order is hidden in the eigenphases $s/r$ of $U_a$.
\end{solbox}
\qsep

\needspace{13\baselineskip}
\begin{qitem}{25}
\qmeta{Single-correct MCQ}{$\ket1$ as sum of eigenvectors}
With $\ket{u_s} = \frac{1}{\sqrt r}\sum_{k=0}^{r-1}e^{-2\pi isk/r}\ket{a^k \bmod N}$, the superposition $\frac{1}{\sqrt r}\sum_{s=0}^{r-1}\ket{u_s}$ equals
\par\smallskip\noindent\begin{tabularx}{\linewidth}{@{}*{2}{L}@{}}
\optlabel{A}\,$\ket a$ & \correctlabel{B}\,$\ket1$ \\[3pt]
\optlabel{C}\,$\ket0$ & \optlabel{D}\,$\frac{1}{\sqrt N}\sum_{y=0}^{N-1}\ket y$
\end{tabularx}
\end{qitem}
\begin{solbox}{B}
\textbf{\color{brand}Correct answer.} \textcolor{okgreen}{\textbf{(B)}} $\ket1$

\textbf{\color{brand}Concepts required.} The eigenvectors $\ket{u_s}$ are Fourier transforms of the orbit; summing over $s$ inverts the transform.

\textbf{\color{brand}Relevant formulae.} $\sum_{s=0}^{r-1}e^{-2\pi isk/r} = r\,\delta_{k,0}$, so $\frac{1}{\sqrt r}\sum_s\ket{u_s} = \ket{a^0} = \ket1$.

\textbf{\color{brand}Step-by-step explanation.} $\frac{1}{\sqrt r}\sum_s\ket{u_s} = \frac1r\sum_k\Big(\sum_se^{-2\pi isk/r}\Big)\ket{a^k} = \sum_k\delta_{k,0}\ket{a^k} = \ket{a^0} = \ket1$, because the sum of
$r$-th roots of unity vanishes unless $k \equiv 0$. That is why the work register is initialised to $\ket1$: it is an equal superposition of all eigenvectors,
and QPE returns a random $s/r$.

\textbf{\color{brand}Why the other options are incorrect.}
\begin{itemize}[leftmargin=2.4em, itemsep=1pt, topsep=2pt]
\item[\textbf{(A)}] $\ket a = \ket{a^1}$ would need the phase factors $e^{2\pi is/r}$ inserted before summing.
\item[\textbf{(C)}] $\ket0$ is outside the orbit $\{a^k \bmod N\}$ of $\ket1$, so it cannot be a combination of the $\ket{u_s}$.
\item[\textbf{(D)}] The uniform superposition over all residues is not in the span of one orbit (it includes non-coprime values).
\end{itemize}
\textbf{\color{brand}Key takeaway.} Initialising the work register to $\ket1$ prepares all eigenvectors at once, so QPE yields a random $s/r$.
\end{solbox}
\qsep

\needspace{13\baselineskip}
\begin{qitem}{26}
\qmeta{Single-correct MCQ}{Phase kickback}
A control qubit in $\ket+$ controls a unitary $U$ acting on its eigenvector $\ket u$, $U\ket u = e^{i\phi}\ket u$, with $\phi = 2\pi/3$. Afterwards, the
expectation value $\langle X\rangle$ of the control qubit is
\par\smallskip\noindent\begin{tabularx}{\linewidth}{@{}*{4}{L}@{}}
\optlabel{A}\,$1/2$ & \optlabel{B}\,$1$ & \optlabel{C}\,$-\sqrt3/2$ & \correctlabel{D}\,$-1/2$
\end{tabularx}
\end{qitem}
\begin{solbox}{D}
\textbf{\color{brand}Correct answer.} \textcolor{okgreen}{\textbf{(D)}} $-1/2$

\textbf{\color{brand}Concepts required.} Phase kickback moves the eigenphase of the target onto the $\ket1$ branch of the control.

\textbf{\color{brand}Relevant formulae.} $\ket+\ket u \mapsto \frac{1}{\sqrt2}(\ket0 + e^{i\phi}\ket1)\ket u$; $\langle X\rangle = \cos\phi$, $\langle Y\rangle = \sin\phi$.

\textbf{\color{brand}Step-by-step explanation.} Phase kickback: $\frac{1}{\sqrt2}(\ket0 + \ket1)\ket u \mapsto \frac{1}{\sqrt2}(\ket0 + e^{i\phi}\ket1)\ket u$ --- the work register is unchanged, the phase moves to
the control. For $\frac{1}{\sqrt2}(\ket0 + e^{i\phi}\ket1)$, $\langle X\rangle = \cos\phi = \cos\frac{2\pi}{3} = -\frac12$ (and $\langle Y\rangle = \sin\phi = \frac{\sqrt3}{2}$).

\textbf{\color{brand}Why the other options are incorrect.}
\begin{itemize}[leftmargin=2.4em, itemsep=1pt, topsep=2pt]
\item[\textbf{(A)}] $1/2 = \cos(\pi/3)$ uses half of the phase.
\item[\textbf{(B)}] $\langle X\rangle = 1$ is the value before the controlled gate; the kickback rotates the state.
\item[\textbf{(C)}] $-\sqrt3/2$ mixes up sine and cosine; $\sin(2\pi/3) = +\sqrt3/2$ is $\langle Y\rangle$.
\end{itemize}
\textbf{\color{brand}Key takeaway.} The control qubit records the phase; the eigenvector target is untouched.
\end{solbox}
\qsep

\needspace{13\baselineskip}
\begin{qitem}{27}
\qmeta{Single-correct MCQ}{Phases of $U^{2^j}$}
In phase estimation, counting qubit $j$ controls $U^{2^j}$. If $U\ket u = e^{2\pi i/6}\ket u$, the phase kicked back onto counting qubit $j = 2$ is
\par\smallskip\noindent\begin{tabularx}{\linewidth}{@{}*{4}{L}@{}}
\optlabel{A}\,$2\pi$ & \optlabel{B}\,$\pi/3$ & \correctlabel{C}\,$4\pi/3$ & \optlabel{D}\,$2\pi/3$
\end{tabularx}
\end{qitem}
\begin{solbox}{C}
\textbf{\color{brand}Correct answer.} \textcolor{okgreen}{\textbf{(C)}} $4\pi/3$

\textbf{\color{brand}Concepts required.} Applying $U^{2^j}$ to an eigenvector multiplies its phase by $2^j$.

\textbf{\color{brand}Relevant formulae.} Phase on qubit $j$: $2\pi\varphi\cdot2^j$; $\varphi = 1/6$, $j = 2$: $2\pi\cdot4/6 = 4\pi/3$.

\textbf{\color{brand}Step-by-step explanation.} $U^{2^j}\ket u = e^{2\pi i\cdot2^j/6}\ket u$; for $j = 2$: $2\pi\cdot\frac46 = \frac{4\pi}{3}$ (equivalently $-\frac{2\pi}{3}$). Each counting qubit receives the phase
doubled relative to the previous one --- the binary-fraction pattern that the inverse QFT decodes.

\textbf{\color{brand}Why the other options are incorrect.}
\begin{itemize}[leftmargin=2.4em, itemsep=1pt, topsep=2pt]
\item[\textbf{(A)}] $2\pi$ would correspond to $U^6$; $2^j = 4$, not $6$.
\item[\textbf{(B)}] $\pi/3 = 2\pi/6$ is the phase for $j = 0$ ($U$ itself).
\item[\textbf{(D)}] $2\pi/3$ is the phase for $j = 1$ ($U^2$).
\end{itemize}
\textbf{\color{brand}Key takeaway.} Counting qubit $j$ carries the binary digit pattern $2^j\varphi \bmod 1$.
\end{solbox}
\qsep

\needspace{13\baselineskip}
\begin{qitem}{28}
\qmeta{Single-correct MCQ}{Hadamard transform signs}
$(H\otimes H)\ket{11}$ equals
\par\smallskip\noindent\begin{tabularx}{\linewidth}{@{}*{2}{L}@{}}
\correctlabel{A}\,$\frac12(\ket{00} - \ket{01} - \ket{10} + \ket{11})$ & \optlabel{B}\,$\frac{1}{\sqrt2}(\ket{00} - \ket{11})$ \\[3pt]
\optlabel{C}\,$\frac12(\ket{00} - \ket{01} - \ket{10} - \ket{11})$ & \optlabel{D}\,$\frac12(\ket{00} + \ket{01} + \ket{10} + \ket{11})$
\end{tabularx}
\end{qitem}
\begin{solbox}{A}
\textbf{\color{brand}Correct answer.} \textcolor{okgreen}{\textbf{(A)}} $\frac12(\ket{00} - \ket{01} - \ket{10} + \ket{11})$

\textbf{\color{brand}Concepts required.} $H^{\otimes n}\ket x = 2^{-n/2}\sum_y(-1)^{x\cdot y}\ket y$.

\textbf{\color{brand}Relevant formulae.} $x = 11$: signs $(-1)^{y_1 + y_2}$: $+, -, -, +$ for $00, 01, 10, 11$.

\textbf{\color{brand}Step-by-step explanation.} $H\ket1 = \frac{1}{\sqrt2}(\ket0 - \ket1)$, so $H\ket1\otimes H\ket1 = \frac12(\ket0 - \ket1)(\ket0 - \ket1) = \frac12(\ket{00} - \ket{01} - \ket{10} + \ket{11})$.
In general $H^{\otimes n}\ket x = 2^{-n/2}\sum_y(-1)^{x\cdot y}\ket y$: signs depend on the parity of $x\cdot y$.

\textbf{\color{brand}Why the other options are incorrect.}
\begin{itemize}[leftmargin=2.4em, itemsep=1pt, topsep=2pt]
\item[\textbf{(B)}] A Bell-type state cannot arise from local Hadamards on a product state.
\item[\textbf{(C)}] The last sign is $(-1)^{1 + 1} = +1$, not $-1$.
\item[\textbf{(D)}] All plus signs is $(H\otimes H)\ket{00}$.
\end{itemize}
\textbf{\color{brand}Key takeaway.} Hadamard transforms encode parities in signs --- the simplest Fourier transform.
\end{solbox}
\qsep

\needspace{13\baselineskip}
\begin{qitem}{29}
\qmeta{Single-correct MCQ}{Two-path interference}
Two computational paths lead to the same outcome with amplitudes $\frac12$ and $\frac12e^{i\theta}$. For $\theta = 2\pi/3$ the probability of that outcome is
\par\smallskip\noindent\begin{tabularx}{\linewidth}{@{}*{4}{L}@{}}
\optlabel{A}\,$0$ & \optlabel{B}\,$1/2$ & \correctlabel{C}\,$1/4$ & \optlabel{D}\,$3/4$
\end{tabularx}
\end{qitem}
\begin{solbox}{C}
\textbf{\color{brand}Correct answer.} \textcolor{okgreen}{\textbf{(C)}} $1/4$

\textbf{\color{brand}Concepts required.} Probabilities come from the squared modulus of the \emph{sum} of amplitudes of indistinguishable paths.

\textbf{\color{brand}Relevant formulae.} $\big|\frac12(1 + e^{i\theta})\big|^2 = \cos^2(\theta/2)$; $\theta = 2\pi/3$: $\cos^2(\pi/3) = 1/4$.

\textbf{\color{brand}Step-by-step explanation.} Amplitudes add before squaring: $\big|\frac12(1 + e^{i\theta})\big|^2 = \frac14|1 + e^{i\theta}|^2 = \frac14(2 + 2\cos\theta) = \cos^2\frac\theta2 = \cos^2\frac\pi3 = \frac14$.
Interference concentrates probability where the phases agree; in Shor's algorithm the inverse QFT arranges this for $y \approx k2^t/r$.

\textbf{\color{brand}Why the other options are incorrect.}
\begin{itemize}[leftmargin=2.4em, itemsep=1pt, topsep=2pt]
\item[\textbf{(A)}] $0$ needs $\theta = \pi$ (complete destructive interference).
\item[\textbf{(B)}] $1/2$ adds probabilities ($\frac14 + \frac14$) instead of amplitudes, ignoring interference.
\item[\textbf{(D)}] $3/4 = \sin^2(\pi/3)$ is the probability of the other output port.
\end{itemize}
\textbf{\color{brand}Key takeaway.} Interference, not parallelism alone, makes the right answer likely.
\end{solbox}
\qsep

\needspace{13\baselineskip}
\begin{qitem}{30}
\qmeta{Single-correct MCQ}{Reversible modular exponentiation}
Why does Shor's circuit compute $\ket x\ket1 \mapsto \ket x\ket{a^x \bmod N}$ instead of $\ket x \mapsto \ket{a^x \bmod N}$ in place?
\begin{itemize}[leftmargin=2.6em, labelsep=0.4em, itemsep=2pt, topsep=3pt]
\item[\correctlabel{A}] the map $x \mapsto a^x \bmod N$ is not injective ($x$ and $x + r$ give the same value), so it cannot be a unitary on a single register; keeping $x$ makes it reversible
\item[\optlabel{B}] because a single register cannot hold numbers larger than $N$
\item[\optlabel{C}] because the Hadamard gates must act on a different register than the QFT
\item[\optlabel{D}] because measurement requires two registers
\end{itemize}
\end{qitem}
\begin{solbox}{A}
\textbf{\color{brand}Correct answer.} \textcolor{okgreen}{\textbf{(A)}} the map $x \mapsto a^x \bmod N$ is not injective ($x$ and $x + r$ give the same value), so it cannot be a unitary on a single register; keeping $x$ makes it reversible

\textbf{\color{brand}Concepts required.} Unitary operations are invertible; a many-to-one function must keep its input to be reversible.

\textbf{\color{brand}Relevant formulae.} $a^{x + r} \equiv a^x$; $\ket x\ket y \mapsto \ket x\ket{y\,a^x \bmod N}$ is reversible.

\textbf{\color{brand}Step-by-step explanation.} Quantum gates are unitary, hence invertible. Since $a^{x + r} \equiv a^x$, an in-place map would send different inputs to the same output. Keeping the input
$x$ and writing the result into a second register (here by multiplying the work register $\ket1$ by $a^x$) gives a reversible map and entangles the
registers: $\sum_x\ket x\ket{a^x \bmod N}$.

\textbf{\color{brand}Why the other options are incorrect.}
\begin{itemize}[leftmargin=2.4em, itemsep=1pt, topsep=2pt]
\item[\textbf{(B)}] Register size is not the issue: $a^x \bmod N < N$ always fits in $n$ qubits.
\item[\textbf{(C)}] The Hadamards and the QFT both act on the counting register; this is not why two registers are used.
\item[\textbf{(D)}] Measurement can be performed on any register; it does not require a second one.
\end{itemize}
\textbf{\color{brand}Key takeaway.} Reversibility forces the input to be kept, which entangles the two registers.
\end{solbox}
\qsep

\needspace{13\baselineskip}
\begin{qitem}{31}
\qmeta{Single-correct MCQ}{CNOT cost of a Fredkin gate}
A controlled-SWAP (Fredkin) gate is compiled as $\mathrm{CNOT}_{b\to a}\cdot\mathrm{Toffoli}_{c,a\to b}\cdot\mathrm{CNOT}_{b\to a}$, and each Toffoli uses the
standard decomposition with 6 CNOTs. The number of CNOTs per controlled-SWAP is
\par\smallskip\noindent\begin{tabularx}{\linewidth}{@{}*{4}{L}@{}}
\optlabel{A}\,$9$ & \correctlabel{B}\,$8$ & \optlabel{C}\,$3$ & \optlabel{D}\,$6$
\end{tabularx}
\end{qitem}
\begin{solbox}{B}
\textbf{\color{brand}Correct answer.} \textcolor{okgreen}{\textbf{(B)}} $8$

\textbf{\color{brand}Concepts required.} A Fredkin gate can be built from two CNOTs around a Toffoli; a Toffoli needs 6 CNOTs in the standard decomposition.

\textbf{\color{brand}Relevant formulae.} $\mathrm{CSWAP}_{c;a,b} = \mathrm{CX}_{b\to a}\,\mathrm{CCX}_{c,a\to b}\,\mathrm{CX}_{b\to a}$: $1 + 6 + 1 = 8$.

\textbf{\color{brand}Step-by-step explanation.} $1 + 6 + 1 = 8$ CNOTs (plus single-qubit $H$, $T$, $T^\dagger$ gates). A plain SWAP needs $3$ CNOTs, but the control turns the middle CNOT into a Toffoli.
Gate-level modular multipliers are built from many such gates, which is why compiled Shor circuits are much larger than the notebook's matrix gates.

\textbf{\color{brand}Why the other options are incorrect.}
\begin{itemize}[leftmargin=2.4em, itemsep=1pt, topsep=2pt]
\item[\textbf{(A)}] $9 = 3\times3$ treats every CNOT of a SWAP as a 3-CNOT gate; only the middle CNOT must become a (6-CNOT) Toffoli.
\item[\textbf{(C)}] $3$ CNOTs make an uncontrolled SWAP; the control requires the Toffoli.
\item[\textbf{(D)}] $6$ counts only the Toffoli and forgets the two outer CNOTs.
\end{itemize}
\textbf{\color{brand}Key takeaway.} Controlled arithmetic is costly: every controlled SWAP is a Toffoli plus two CNOTs.
\end{solbox}
\qsep

\needspace{13\baselineskip}
\begin{qitem}{32}
\qmeta{Single-correct MCQ}{Implicit measurement}
\texttt{Shor.ipynb} measures only the counting register. Compared with also measuring the work register first, the distribution of the counting
outcome $y$ is
\begin{itemize}[leftmargin=2.6em, labelsep=0.4em, itemsep=2pt, topsep=3pt]
\item[\optlabel{A}] different: the peaks become twice as wide
\item[\optlabel{B}] undefined, because entangled registers cannot be measured separately
\item[\correctlabel{C}] exactly the same: a measurement of the other register that is never used cannot change the marginal statistics of $y$
\item[\optlabel{D}] different: without measuring the work register no periodic structure appears
\end{itemize}
\end{qitem}
\begin{solbox}{C}
\textbf{\color{brand}Correct answer.} \textcolor{okgreen}{\textbf{(C)}} exactly the same: a measurement of the other register that is never used cannot change the marginal statistics of $y$

\textbf{\color{brand}Concepts required.} Measuring a subsystem that is never used again does not change the reduced state of the rest.

\textbf{\color{brand}Relevant formulae.} $\rho_{\text{count}} = \operatorname{Tr}_{\text{work}}\ket\Psi\!\bra\Psi = \sum_wP(w)\rho_{\text{count}|w}$.

\textbf{\color{brand}Step-by-step explanation.} The work register is not acted on after the controlled multiplications. Measuring it (or not) only changes which conditional state the counting
register is in; averaged over the unseen result, the reduced state of the counting register --- and hence $P(y)$ --- is identical (principle of implicit
measurement). Textbook derivations ``measure the second register'' only to simplify the algebra.

\textbf{\color{brand}Why the other options are incorrect.}
\begin{itemize}[leftmargin=2.4em, itemsep=1pt, topsep=2pt]
\item[\textbf{(A)}] Peak widths depend on $t$ and $r$, not on whether the idle work register is measured.
\item[\textbf{(B)}] Entangled registers can be measured separately; the marginal statistics are well defined.
\item[\textbf{(D)}] The periodic structure is in the entangled state already; the measurement of the work register only simplifies textbook algebra.
\end{itemize}
\textbf{\color{brand}Key takeaway.} Principle of implicit (deferred) measurement: unused registers may be measured or ignored.
\end{solbox}
\qsep

\needspace{13\baselineskip}
\begin{qitem}{33}
\qmeta{Single-correct MCQ}{Notebook qubit count}
In \texttt{Shor.ipynb} the work register has $n = $ \texttt{N.bit\_length()} qubits and the counting register $t = 2n$. The total number of qubits for
$N = 143$ is
\par\smallskip\noindent\begin{tabularx}{\linewidth}{@{}*{4}{L}@{}}
\correctlabel{A}\,$24$ & \optlabel{B}\,$21$ & \optlabel{C}\,$22$ & \optlabel{D}\,$16$
\end{tabularx}
\end{qitem}
\begin{solbox}{A}
\textbf{\color{brand}Correct answer.} \textcolor{okgreen}{\textbf{(A)}} $24$

\textbf{\color{brand}Concepts required.} The notebook uses $n = \lceil\log_2(N + 1)\rceil$ work qubits and $t = 2n$ counting qubits: $3n$ in total.

\textbf{\color{brand}Relevant formulae.} $143 = 10001111_2$: $n = 8$, $3n = 24$.

\textbf{\color{brand}Step-by-step explanation.} $143 = 10001111_2$ has $8$ bits, so $n = 8$, $t = 16$ and $t + n = 24$ (the notebook's comment lists ``143 -> 24''). $16$ forgets the work register;
$22 = 2\lceil\log_2 143\rceil + 6$ has no basis.

\textbf{\color{brand}Why the other options are incorrect.}
\begin{itemize}[leftmargin=2.4em, itemsep=1pt, topsep=2pt]
\item[\textbf{(B)}] $21$ would be $3\cdot7$, but $143 > 127 = 2^7 - 1$ needs $8$ bits.
\item[\textbf{(C)}] $22$ does not follow from $t = 2n$ with $n = 8$; no formula of the notebook gives it.
\item[\textbf{(D)}] $16$ is only the counting register $t = 2n$.
\end{itemize}
\textbf{\color{brand}Key takeaway.} Qubits scale as $3n$ here: three per bit of $N$.
\end{solbox}
\qsep

\needspace{13\baselineskip}
\begin{qitem}{34}
\qmeta{Single-correct MCQ}{State-vector memory}
The memory of a double-precision state vector (16 bytes per amplitude) for the 24-qubit circuit of the previous question is about
\par\smallskip\noindent\begin{tabularx}{\linewidth}{@{}*{4}{L}@{}}
\optlabel{A}\,$16.8$ MB & \optlabel{B}\,$384$ bytes & \optlabel{C}\,$4.3$ GB & \correctlabel{D}\,$268$ MB
\end{tabularx}
\end{qitem}
\begin{solbox}{D}
\textbf{\color{brand}Correct answer.} \textcolor{okgreen}{\textbf{(D)}} $268$ MB

\textbf{\color{brand}Concepts required.} A state vector of $q$ qubits has $2^q$ complex amplitudes of $16$ bytes (complex128).

\textbf{\color{brand}Relevant formulae.} $M = 16\cdot2^q$ bytes; $q = 24$: $268\,435\,456$ B $\approx 268$ MB.

\textbf{\color{brand}Step-by-step explanation.} $16\times2^{24} = 268\,435\,456$ bytes $\approx 268$ MB (the notebook's table). $384 = 16\times24$ is linear; $16.8$ MB is $2^{20}\times16$ (20 qubits);
$4.3$ GB corresponds to 28 qubits. Each extra bit of $N$ adds 3 qubits, i.e.\ multiplies the memory by 8.

\textbf{\color{brand}Why the other options are incorrect.}
\begin{itemize}[leftmargin=2.4em, itemsep=1pt, topsep=2pt]
\item[\textbf{(A)}] $16.8$ MB is $16\cdot2^{20}$, the memory of $20$ qubits.
\item[\textbf{(B)}] $384$ bytes grows linearly with the number of qubits; state vectors grow exponentially.
\item[\textbf{(C)}] $4.3$ GB is $16\cdot2^{28}$, the memory of $28$ qubits.
\end{itemize}
\textbf{\color{brand}Key takeaway.} Each extra qubit doubles the memory; each extra bit of $N$ multiplies it by $8$.
\end{solbox}
\qsep

\needspace{13\baselineskip}
\begin{qitem}{35}
\qmeta{Matching lists}{Roles of circuit blocks}
Match each operation of the order-finding circuit (List-I) with its role (List-II).
\par\smallskip\noindent{\renewcommand{\arraystretch}{1.35}%
\begin{tabularx}{\linewidth}{@{}c>{\raggedright\arraybackslash}p{0.33\linewidth}@{\hspace{1.2em}}cL@{}}
\toprule
\multicolumn{2}{@{}l}{\textbf{List-I}} & \multicolumn{2}{l@{}}{\textbf{List-II}} \\ \midrule
\textbf{(P)} & $H$ on every counting qubit & \textbf{(1)} & turns the phases $2\pi\cdot2^js/r$ into a binary estimate $y \approx 2^ts/r$ \\
\textbf{(Q)} & controlled $U_{a^{2^j} \bmod N}$ & \textbf{(2)} & prepares the work register in $\ket1$, an equal superposition of the eigenvectors $\ket{u_s}$ \\
\textbf{(R)} & inverse QFT on the counting register & \textbf{(3)} & creates the uniform superposition over all exponents $x$ \\
\textbf{(S)} & $X$ on the last work qubit & \textbf{(4)} & performs modular exponentiation / kicks the eigenphases back onto the counting qubits \\
\bottomrule
\end{tabularx}}
\par\smallskip\noindent\begin{tabularx}{\linewidth}{@{}*{2}{L}@{}}
\optlabel{A}\,P$\to$3, Q$\to$1, R$\to$4, S$\to$2 & \correctlabel{B}\,P$\to$3, Q$\to$4, R$\to$1, S$\to$2 \\[3pt]
\optlabel{C}\,P$\to$3, Q$\to$4, R$\to$2, S$\to$1 & \optlabel{D}\,P$\to$4, Q$\to$3, R$\to$1, S$\to$2
\end{tabularx}
\end{qitem}
\begin{solbox}{B}
\textbf{\color{brand}Correct answer.} \textcolor{okgreen}{\textbf{(B)}} P$\to$3, Q$\to$4, R$\to$1, S$\to$2

\textbf{\color{brand}Concepts required.} Order finding: superposition, controlled modular exponentiation (kickback), inverse QFT, measurement.

\textbf{\color{brand}Relevant formulae.} $\frac{1}{\sqrt{2^t}}\sum_x\ket x\ket1 \to \frac{1}{\sqrt{2^t}}\sum_x\ket x\ket{a^x} \to$ peaks at $y \approx k2^t/r$.

\textbf{\color{brand}Step-by-step explanation.} \begin{itemize}\itemsep0pt
\item Hadamards: superposition over $x$ (3);
\item controlled powers of $U_a$: $\ket x\ket1 \mapsto \ket x\ket{a^x}$, equivalently phase kickback (4);
\item inverse QFT: interference that converts the periodic phase pattern into peaks at $y \approx k2^t/r$ (1);
\item $X$ on the least significant work qubit: $\ket{0001} = \ket1$ (2).
\end{itemize}
Hence P$\to$3, Q$\to$4, R$\to$1, S$\to$2.

\textbf{\color{brand}Why the other options are incorrect.}
\begin{itemize}[leftmargin=2.4em, itemsep=1pt, topsep=2pt]
\item[\textbf{(A)}] This gives the controlled multiplications the decoding role of the inverse QFT and the QFT the exponentiation.
\item[\textbf{(C)}] This swaps the roles of the inverse QFT and the $X$ gate; $X$ only prepares $\ket1$ in the work register.
\item[\textbf{(D)}] This gives the Hadamards the role of modular exponentiation and vice versa; the Hadamards only create the superposition.
\end{itemize}
\textbf{\color{brand}Key takeaway.} Know what each block does: it is the basis for every resource and error discussion.
\end{solbox}
\qsep

\needspace{13\baselineskip}
\begin{qitem}{36}
\qmeta{Matching lists}{Phase gates $R_1(\lambda)$}
The phase gate is $R_1(\lambda) = \operatorname{diag}(1, e^{i\lambda})$ (controlled versions appear in the QFT). Match each angle (List-I) with the named gate (List-II).
\par\smallskip\noindent{\renewcommand{\arraystretch}{1.35}%
\begin{tabularx}{\linewidth}{@{}c>{\raggedright\arraybackslash}p{0.33\linewidth}@{\hspace{1.2em}}cL@{}}
\toprule
\multicolumn{2}{@{}l}{\textbf{List-I}} & \multicolumn{2}{l@{}}{\textbf{List-II}} \\ \midrule
\textbf{(P)} & $R_1(\pi)$ & \textbf{(1)} & $S^\dagger$ \\
\textbf{(Q)} & $R_1(\pi/2)$ & \textbf{(2)} & $T$ \\
\textbf{(R)} & $R_1(\pi/4)$ & \textbf{(3)} & $Z$ \\
\textbf{(S)} & $R_1(-\pi/2)$ & \textbf{(4)} & $S$ \\
\bottomrule
\end{tabularx}}
\par\smallskip\noindent\begin{tabularx}{\linewidth}{@{}*{2}{L}@{}}
\optlabel{A}\,P$\to$3, Q$\to$4, R$\to$1, S$\to$2 & \correctlabel{B}\,P$\to$3, Q$\to$4, R$\to$2, S$\to$1 \\[3pt]
\optlabel{C}\,P$\to$3, Q$\to$2, R$\to$4, S$\to$1 & \optlabel{D}\,P$\to$4, Q$\to$3, R$\to$2, S$\to$1
\end{tabularx}
\end{qitem}
\begin{solbox}{B}
\textbf{\color{brand}Correct answer.} \textcolor{okgreen}{\textbf{(B)}} P$\to$3, Q$\to$4, R$\to$2, S$\to$1

\textbf{\color{brand}Concepts required.} $R_1(\lambda) = \operatorname{diag}(1, e^{i\lambda})$ contains $Z$, $S$, $T$ and their inverses as special cases.

\textbf{\color{brand}Relevant formulae.} $e^{i\pi} = -1$, $e^{i\pi/2} = i$, $e^{i\pi/4}$, $e^{-i\pi/2} = -i$.

\textbf{\color{brand}Step-by-step explanation.} $R_1(\pi) = \operatorname{diag}(1, -1) = Z$ (3); $R_1(\pi/2) = \operatorname{diag}(1, i) = S$ (4); $R_1(\pi/4) = \operatorname{diag}(1, e^{i\pi/4}) = T$ (2);
$R_1(-\pi/2) = \operatorname{diag}(1, -i) = S^\dagger$ (1). The inverse QFT of the notebook uses \texttt{cr1(-pi/2**(k-j))}: controlled $S^\dagger$, $T^\dagger$, \dots

\textbf{\color{brand}Why the other options are incorrect.}
\begin{itemize}[leftmargin=2.4em, itemsep=1pt, topsep=2pt]
\item[\textbf{(A)}] This exchanges $T$ and $S^\dagger$; $R_1(-\pi/2)$ has $-i$ on the diagonal, which is $S^\dagger$.
\item[\textbf{(C)}] This exchanges $S$ and $T$: $R_1(\pi/2) = \operatorname{diag}(1, i) = S$, while $T$ has $e^{i\pi/4}$.
\item[\textbf{(D)}] This assigns $R_1(\pi)$ to $S$; $R_1(\pi) = \operatorname{diag}(1, -1) = Z$.
\end{itemize}
\textbf{\color{brand}Key takeaway.} The inverse QFT uses controlled $R_1(-\pi/2^k)$: controlled $S^\dagger$, $T^\dagger$ and finer rotations.
\end{solbox}
\qsep

\needspace{13\baselineskip}
\begin{qitem}{37}
\qmeta{Assertion--Reason}{Kickback and controls}
\textbf{Assertion (A):} When the target register is in an eigenvector of $U$, a controlled-$U$ gate leaves the target register unchanged.\\[1mm]
\textbf{Reason (R):} Controlled gates never change the state of their control qubit.
\begin{itemize}[leftmargin=2.6em, labelsep=0.4em, itemsep=2pt, topsep=3pt]
\item[\optlabel{A}] Both Assertion (A) and Reason (R) are true, and R is the correct explanation of A
\item[\optlabel{B}] Both A and R are true, but R is \emph{not} the correct explanation of A
\item[\correctlabel{C}] A is true, but R is false
\item[\optlabel{D}] A is false, but R is true
\end{itemize}
\end{qitem}
\begin{solbox}{C}
\textbf{\color{brand}Correct answer.} \textcolor{okgreen}{\textbf{(C)}} A is true, but R is false

\textbf{\color{brand}Concepts required.} Phase kickback: a controlled-$U$ on an eigenvector changes the control (in superposition), not the target.

\textbf{\color{brand}Relevant formulae.} $(\alpha\ket0 + \beta\ket1)\ket u \mapsto (\alpha\ket0 + \beta e^{i\phi}\ket1)\ket u$.

\textbf{\color{brand}Step-by-step explanation.} \textbf{A is true}: $\ket c\ket u \mapsto \ket c\,e^{ic\phi}\ket u$ --- the target stays $\ket u$, only a phase appears. \textbf{R is false}: that phase belongs to the
control branch $\ket1$, so a control in superposition changes, $\ket+ \mapsto \frac{1}{\sqrt2}(\ket0 + e^{i\phi}\ket1)$ (phase kickback). This is precisely how
phase estimation and Shor's algorithm extract information. Answer: (C).

\textbf{\color{brand}Why the other options are incorrect.}
\begin{itemize}[leftmargin=2.4em, itemsep=1pt, topsep=2pt]
\item[\textbf{(A)}] is wrong because R is false: the control's relative phase changes, as phase kickback shows.
\item[\textbf{(B)}] is wrong for the same reason: R is false, so it cannot be true but irrelevant.
\item[\textbf{(D)}] is wrong because A is true: the target remains the eigenvector $\ket u$.
\end{itemize}
\textbf{\color{brand}Key takeaway.} In QPE the information flows from target to control --- the opposite of the classical intuition.
\end{solbox}
\qsep

\needspace{13\baselineskip}
\begin{qitem}{38}
\qmeta{Statement I / II}{Parallelism vs interference}
\textbf{Statement I:} Because the circuit evaluates $a^x \bmod N$ for all $x$ in superposition, measuring the registers directly reveals the period $r$.\\[1mm]
\textbf{Statement II:} Interference produced by the inverse QFT is needed to turn the periodic structure into high-probability measurement outcomes.
\begin{itemize}[leftmargin=2.6em, labelsep=0.4em, itemsep=2pt, topsep=3pt]
\item[\optlabel{A}] Both Statement I and Statement II are true
\item[\optlabel{B}] Both Statement I and Statement II are false
\item[\optlabel{C}] Statement I is true, but Statement II is false
\item[\correctlabel{D}] Statement I is false, but Statement II is true
\end{itemize}
\end{qitem}
\begin{solbox}{D}
\textbf{\color{brand}Correct answer.} \textcolor{okgreen}{\textbf{(D)}} Statement I is false, but Statement II is true

\textbf{\color{brand}Concepts required.} Quantum parallelism produces a superposition, but a measurement returns one sample; interference extracts global properties.

\textbf{\color{brand}Relevant formulae.} Measuring $\sum_x\ket x\ket{f(x)}$ gives one pair $(x, f(x))$; the inverse QFT concentrates amplitude at $y \approx k2^t/r$.

\textbf{\color{brand}Step-by-step explanation.} \textbf{I is false}: a measurement returns one random $x$ and one value $a^x \bmod N$ --- no more information than a single classical evaluation. Quantum
parallelism alone does not solve factoring. \textbf{II is true}: the inverse QFT makes the amplitudes interfere constructively near $y = k2^t/r$, so that one
measurement carries information about $r$. Answer: (D).

\textbf{\color{brand}Why the other options are incorrect.}
\begin{itemize}[leftmargin=2.4em, itemsep=1pt, topsep=2pt]
\item[\textbf{(A)}] is wrong because Statement I is false: a direct measurement yields one random evaluation, not the period.
\item[\textbf{(B)}] is wrong because Statement II is true: interference via the inverse QFT is the essential step.
\item[\textbf{(C)}] has the truth values reversed: I is false and II is true.
\end{itemize}
\textbf{\color{brand}Key takeaway.} Parallelism alone does not solve factoring; interference does the work.
\end{solbox}
\qsep

\needspace{20\baselineskip}
\begin{passagebox}{Paragraph for Questions 39 and 40}
Quantum phase estimation (QPE) with $t = 3$ counting qubits is applied to a unitary $U$ and an eigenvector $\ket u$ with
$U\ket u = e^{2\pi i\varphi}\ket u$, $0 \le \varphi < 1$. After the Hadamards, the controlled-$U^{2^j}$ gates and the inverse QFT, the counting register is
measured and gives $y \in \{0, \dots, 7\}$ with probability
\[ P(y) = \Big|\frac18\sum_{x=0}^{7}e^{2\pi ix(\varphi - y/8)}\Big|^2 = \frac{\sin^2\big(8\pi(\varphi - y/8)\big)}{64\,\sin^2\big(\pi(\varphi - y/8)\big)}\quad(P = 1\text{ if }\varphi = y/8). \]
The estimate of the phase is $y/8$; $y$ is written in binary with the most significant bit first.
\end{passagebox}
\needspace{13\baselineskip}
\begin{qitem}{39}
\qmeta{Paragraph based}{QPE: exact phase}
For $\varphi = 3/8$ the measurement gives
\par\smallskip\noindent\begin{tabularx}{\linewidth}{@{}*{2}{L}@{}}
\correctlabel{A}\,$y = 3$ (bits $011$) with probability $1$ & \optlabel{B}\,$y = 6$ (bits $110$) with probability $1$ \\[3pt]
\optlabel{C}\,$y = 3$ with probability $1/2$ & \optlabel{D}\,$y = 3$ or $y = 4$ with probability $1/2$ each
\end{tabularx}
\end{qitem}
\begin{solbox}{A}
\textbf{\color{brand}Correct answer.} \textcolor{okgreen}{\textbf{(A)}} $y = 3$ (bits $011$) with probability $1$

\textbf{\color{brand}Concepts required.} QPE is exact when the phase has a $t$-bit binary expansion.

\textbf{\color{brand}Relevant formulae.} $\varphi = 3/8 = 0.011_2$; $P(y) = 1$ for $y = 8\varphi = 3$.

\textbf{\color{brand}Step-by-step explanation.} $\varphi = 3/8 = 0.011_2$ is exactly representable with three bits, so all terms in the sum have phase $0$ for $y = 3$: $P(3) = 1$, and $P(y) = 0$ for every
other $y$ (the geometric sum of eighth roots of unity vanishes). $110$ is the bit-reversed string --- the inverse QFT's final SWAPs make the most significant
bit come first.

\textbf{\color{brand}Why the other options are incorrect.}
\begin{itemize}[leftmargin=2.4em, itemsep=1pt, topsep=2pt]
\item[\textbf{(B)}] $110$ is the bit-reversed string; the notebook's inverse QFT includes the SWAPs, so $y = 3$ is read MSB first.
\item[\textbf{(C)}] For an exactly representable phase all amplitude interferes constructively at $y = 3$; probability $1$, not $1/2$.
\item[\textbf{(D)}] Splitting between $3$ and $4$ happens for $\varphi$ halfway between grid points ($7/16$), not for $3/8$.
\end{itemize}
\textbf{\color{brand}Key takeaway.} If $r$ divides $2^t$, Shor's outcomes are exactly $s2^t/r$ (as for $N = 15$).
\end{solbox}
\qsep

\needspace{13\baselineskip}
\begin{qitem}{40}
\qmeta{Paragraph based}{QPE: inexact phase}
For $\varphi = 1/3$ the most likely outcome and its probability are
\par\smallskip\noindent\begin{tabularx}{\linewidth}{@{}*{2}{L}@{}}
\correctlabel{A}\,$y = 3$ with $P \approx 0.69$ & \optlabel{B}\,$y = 3$ with $P \approx 0.41$ \\[3pt]
\optlabel{C}\,$y = 2$ or $3$ with $P = 1/2$ each & \optlabel{D}\,$y = 3$ with $P = 1$
\end{tabularx}
\end{qitem}
\begin{solbox}{A}
\textbf{\color{brand}Correct answer.} \textcolor{okgreen}{\textbf{(A)}} $y = 3$ with $P \approx 0.69$

\textbf{\color{brand}Concepts required.} For phases that are not multiples of $2^{-t}$ the outcome distribution is a peak around $2^t\varphi$.

\textbf{\color{brand}Relevant formulae.} $\delta = 1/3 - 3/8 = -1/24$; $P(3) = \sin^2(\pi/3)/(64\sin^2(\pi/24)) \approx 0.688$.

\textbf{\color{brand}Step-by-step explanation.} $8\varphi = 2.67$, nearest integer $y = 3$, $\delta = \varphi - 3/8 = -\frac{1}{24}$. $P(3) = \frac{\sin^2(8\pi/24)}{64\sin^2(\pi/24)} = \frac{\sin^2(\pi/3)}{64\sin^2(7.5^\circ)}
= \frac{0.75}{64\times0.017037} = \frac{0.75}{1.0904} \approx 0.688$. Because $1/3$ has no finite binary expansion, the probability spreads over neighbouring $y$;
$\frac{4}{\pi^2} \approx 0.41$ is only the worst-case lower bound for the nearest outcome.

\textbf{\color{brand}Why the other options are incorrect.}
\begin{itemize}[leftmargin=2.4em, itemsep=1pt, topsep=2pt]
\item[\textbf{(B)}] $4/\pi^2 \approx 0.405$ is only the worst-case lower bound for the nearest outcome, attained when $\varphi$ lies halfway.
\item[\textbf{(C)}] $P(2) \approx 0.17$ and $P(3) \approx 0.69$: $8\varphi = 2.67$ is closer to $3$ than to $2$, so the two are far from equal.
\item[\textbf{(D)}] Probability $1$ needs an exactly representable phase; $1/3$ has the infinite binary expansion $0.0101\ldots_2$.
\end{itemize}
\textbf{\color{brand}Key takeaway.} Peaks have finite width; continued fractions turn the nearest $y$ into $s/r$.
\end{solbox}
\qsep

\newpage

\begin{sectionbanner}SECTION C \enspace$\cdot$\enspace Shor's Algorithm Fundamentals \hfill Q.41--60\end{sectionbanner}

\needspace{13\baselineskip}
\begin{qitem}{41}
\qmeta{Single-correct MCQ}{Steps of Shor's algorithm}
Which sequence describes Shor's factoring algorithm for an odd composite $N$ that is not a prime power?
\begin{itemize}[leftmargin=2.6em, labelsep=0.4em, itemsep=2pt, topsep=3pt]
\item[\optlabel{A}] find the order $r$ by quantum phase estimation; the factors are $r$ and $N/r$
\item[\optlabel{B}] apply the QFT to $\ket N$; measure; the outcome is a factor of $N$
\item[\optlabel{C}] pick random $a$; measure $a^x \bmod N$ in superposition; the most frequent value is a factor
\item[\correctlabel{D}] pick random $a$; if $\gcd(a, N) > 1$ stop (factor found); find the order $r$ of $a$ by quantum phase estimation and continued fractions; if $r$ is even and $a^{r/2} \not\equiv -1$, output $\gcd(a^{r/2} \pm 1, N)$; otherwise repeat
\end{itemize}
\end{qitem}
\begin{solbox}{D}
\textbf{\color{brand}Correct answer.} \textcolor{okgreen}{\textbf{(D)}} pick random $a$; if $\gcd(a, N) > 1$ stop (factor found); find the order $r$ of $a$ by quantum phase estimation and continued fractions; if $r$ is even and $a^{r/2} \not\equiv -1$, output $\gcd(a^{r/2} \pm 1, N)$; otherwise repeat

\textbf{\color{brand}Concepts required.} Shor's algorithm is a classical reduction from factoring to order finding, with a quantum subroutine for the order only.

\textbf{\color{brand}Relevant formulae.} $r = \operatorname{ord}_N(a)$; if $2 \mid r$ and $a^{r/2} \not\equiv -1$: $p = \gcd(a^{r/2} - 1, N)$, $q = \gcd(a^{r/2} + 1, N)$.

\textbf{\color{brand}Step-by-step explanation.} The quantum computer is used only for \emph{order finding}; everything else --- the $\gcd$ test, continued fractions, the checks on $r$ and the
final $\gcd(a^{r/2} \pm 1, N)$ --- is classical. Option 3 confuses the order with a factor (for $N = 15$, $a = 7$: $r = 4$ does not divide $15$).

\textbf{\color{brand}Why the other options are incorrect.}
\begin{itemize}[leftmargin=2.4em, itemsep=1pt, topsep=2pt]
\item[\textbf{(A)}] The order is not a factor: for $N = 15$, $a = 7$, $r = 4 \nmid 15$; the factors come from $\gcd(a^{r/2} \pm 1, N)$.
\item[\textbf{(B)}] A QFT of the number $N$ itself produces no factor; Shor applies the QFT to the counting register after modular exponentiation.
\item[\textbf{(C)}] Measuring $a^x \bmod N$ gives one random value of the function; no interference step and no factor appear.
\end{itemize}
\textbf{\color{brand}Key takeaway.} Quantum part: period finding. Classical part: everything else.
\end{solbox}
\qsep

\needspace{13\baselineskip}
\begin{qitem}{42}
\qmeta{Single-correct MCQ}{$N = 15$, $a = 2$}
For $N = 15$ and $a = 2$, the order $r$ and the factors obtained from $\gcd(a^{r/2} \pm 1, N)$ are
\par\smallskip\noindent\begin{tabularx}{\linewidth}{@{}*{2}{L}@{}}
\optlabel{A}\,$r = 4$; $\ 1$ and $15$ & \correctlabel{B}\,$r = 4$; $\ \gcd(3, 15) = 3$, $\gcd(5, 15) = 5$ \\[3pt]
\optlabel{C}\,$r = 2$; $\ 3$ and $5$ & \optlabel{D}\,$r = 8$; $\ 3$ and $5$
\end{tabularx}
\end{qitem}
\begin{solbox}{B}
\textbf{\color{brand}Correct answer.} \textcolor{okgreen}{\textbf{(B)}} $r = 4$; $\ \gcd(3, 15) = 3$, $\gcd(5, 15) = 5$

\textbf{\color{brand}Concepts required.} For a small $N$ the order can be read off the list of powers; the gcds then give the factors.

\textbf{\color{brand}Relevant formulae.} $2, 4, 8, 16 \equiv 1 \pmod{15}$: $r = 4$; $2^2 = 4$: $\gcd(3, 15) = 3$, $\gcd(5, 15) = 5$.

\textbf{\color{brand}Step-by-step explanation.} $2^1, 2^2, 2^3, 2^4 \equiv 2, 4, 8, 1$, so $r = 4$. $2^{r/2} = 4$, $4 \not\equiv -1 \equiv 14$. $\gcd(4 - 1, 15) = 3$ and $\gcd(4 + 1, 15) = 5$, and $3\times5 = 15$.
$r = 8$ is a multiple of the order, not the order (the smallest positive exponent).

\textbf{\color{brand}Why the other options are incorrect.}
\begin{itemize}[leftmargin=2.4em, itemsep=1pt, topsep=2pt]
\item[\textbf{(A)}] $1$ and $15$ are the trivial gcds that occur only for $a^{r/2} \equiv -1$; here $4 \not\equiv 14$.
\item[\textbf{(C)}] $2^2 = 4 \ne 1$, so $r \ne 2$; the order is $4$ (the factors are right, but for the wrong reason).
\item[\textbf{(D)}] $8$ is a multiple of the order; the order is the smallest exponent with $a^r \equiv 1$.
\end{itemize}
\textbf{\color{brand}Key takeaway.} $N = 15$, $a = 2$ behaves like the notebook default $a = 7$: $r = 4$, factors $3 \times 5$.
\end{solbox}
\qsep

\needspace{13\baselineskip}
\begin{qitem}{43}
\qmeta{Single-correct MCQ}{Failure for $a = N - 1$}
For $N = 15$ and $a = 14$, the order finding succeeds perfectly. The classical post-processing then
\begin{itemize}[leftmargin=2.6em, labelsep=0.4em, itemsep=2pt, topsep=3pt]
\item[\optlabel{A}] succeeds with factors $3$ and $5$
\item[\optlabel{B}] fails because $r$ is odd
\item[\correctlabel{C}] fails: $r = 2$ but $14^{1} \equiv -1 \pmod{15}$, so $\gcd(14 \pm 1, 15)$ gives only the trivial factors $15$ and $1$; another $a$ is needed
\item[\optlabel{D}] fails because $\gcd(14, 15) = 1$
\end{itemize}
\end{qitem}
\begin{solbox}{C}
\textbf{\color{brand}Correct answer.} \textcolor{okgreen}{\textbf{(C)}} fails: $r = 2$ but $14^{1} \equiv -1 \pmod{15}$, so $\gcd(14 \pm 1, 15)$ gives only the trivial factors $15$ and $1$; another $a$ is needed

\textbf{\color{brand}Concepts required.} The reduction fails when $a^{r/2} \equiv -1 \pmod N$, even if the order is found correctly.

\textbf{\color{brand}Relevant formulae.} $14 \equiv -1$, $14^2 \equiv 1$: $r = 2$, $a^{r/2} = 14 \equiv -1$; $\gcd(13, 15) = 1$, $\gcd(15, 15) = 15$.

\textbf{\color{brand}Step-by-step explanation.} $14^2 = 196 = 13\cdot15 + 1$, so $r = 2$ (even). But $a^{r/2} = 14 \equiv -1$: $\gcd(13, 15) = 1$ and $\gcd(15, 15) = 15$ --- trivial. For any $N$, $a = N - 1$ is always such a
failing base. $\gcd(a, N) = 1$ is required, not a reason for failure.

\textbf{\color{brand}Why the other options are incorrect.}
\begin{itemize}[leftmargin=2.4em, itemsep=1pt, topsep=2pt]
\item[\textbf{(A)}] $3$ and $5$ cannot appear: with $a^{r/2} \equiv -1$ both gcds are trivial.
\item[\textbf{(B)}] $r = 2$ is even; the odd-order failure is a different case (e.g.\ $a = 4$ modulo $21$).
\item[\textbf{(D)}] $\gcd(a, N) = 1$ is the normal case required for order finding, not a cause of failure.
\end{itemize}
\textbf{\color{brand}Key takeaway.} $a = N - 1$ always fails: a failed attempt must be reported, then another base tried.
\end{solbox}
\qsep

\needspace{13\baselineskip}
\begin{qitem}{44}
\qmeta{Single-correct MCQ}{Odd order: retry}
For $N = 21$ and $a = 4$, the outcome of the classical checks is
\begin{itemize}[leftmargin=2.6em, labelsep=0.4em, itemsep=2pt, topsep=3pt]
\item[\correctlabel{A}] $r = 3$ is odd, so $a^{r/2}$ is not an integer power; the attempt fails and a new $a$ is chosen
\item[\optlabel{B}] $r = 3$, factors $\gcd(4^1 - 1, 21) = 3$ and $\gcd(4^1 + 1, 21) = 1$
\item[\optlabel{C}] $\gcd(4, 21) = 7$, a lucky factor
\item[\optlabel{D}] $r = 6$, factors $3$ and $7$
\end{itemize}
\end{qitem}
\begin{solbox}{A}
\textbf{\color{brand}Correct answer.} \textcolor{okgreen}{\textbf{(A)}} $r = 3$ is odd, so $a^{r/2}$ is not an integer power; the attempt fails and a new $a$ is chosen

\textbf{\color{brand}Concepts required.} An odd order gives no square root $a^{r/2}$ in the integers; the attempt fails.

\textbf{\color{brand}Relevant formulae.} $4, 16, 64 \equiv 1 \pmod{21}$: $r = 3$.

\textbf{\color{brand}Step-by-step explanation.} $4^1, 4^2, 4^3 \equiv 4, 16, 64 = 3\cdot21 + 1$, so $r = 3$, odd. The reduction needs $r$ even (then $(a^{r/2} - 1)(a^{r/2} + 1) \equiv 0$). A failed attempt must be
reported as a failure, not as a success; one simply repeats with another base ($a = 2$ gives $r = 6$ and the factors $3$, $7$).

\textbf{\color{brand}Why the other options are incorrect.}
\begin{itemize}[leftmargin=2.4em, itemsep=1pt, topsep=2pt]
\item[\textbf{(B)}] The reduction is not applied for odd $r$; ``$4^{1}$'' is not $a^{r/2}$ for $r = 3$, and a gcd of $1$ is no factor anyway.
\item[\textbf{(C)}] $\gcd(4, 21) = 1$; $4$ shares no factor with $21 = 3\cdot7$.
\item[\textbf{(D)}] $4^3 \equiv 1$, so the order is $3$, not $6$; $6$ is the order of $a = 2$ modulo $21$.
\end{itemize}
\textbf{\color{brand}Key takeaway.} Odd order $\Rightarrow$ retry with a new base; never present it as a success.
\end{solbox}
\qsep

\needspace{13\baselineskip}
\begin{qitem}{45}
\qmeta{Single-correct MCQ}{Lucky base}
For $N = 21$ the randomly chosen base is $a = 6$. What happens?
\begin{itemize}[leftmargin=2.6em, labelsep=0.4em, itemsep=2pt, topsep=3pt]
\item[\optlabel{A}] $U_6$ is a valid unitary and phase estimation returns $r = 2$
\item[\optlabel{B}] order finding gives $r = 6$ and the factors $3$, $7$
\item[\correctlabel{C}] $\gcd(6, 21) = 3$ is already a non-trivial factor; no quantum computation is needed
\item[\optlabel{D}] the algorithm fails because $a$ is even
\end{itemize}
\end{qitem}
\begin{solbox}{C}
\textbf{\color{brand}Correct answer.} \textcolor{okgreen}{\textbf{(C)}} $\gcd(6, 21) = 3$ is already a non-trivial factor; no quantum computation is needed

\textbf{\color{brand}Concepts required.} The first step of the algorithm is the classical test $\gcd(a, N)$.

\textbf{\color{brand}Relevant formulae.} $\gcd(6, 21) = 3$, $21/3 = 7$.

\textbf{\color{brand}Step-by-step explanation.} The first classical step is $\gcd(a, N)$. Here $\gcd(6, 21) = 3$, so $21 = 3\times7$ immediately (a ``lucky'' base). Moreover $U_6$ would not be unitary
($6\cdot0 \equiv 6\cdot7 \equiv 0 \pmod{21}$), and $6$ has no order modulo $21$. Even bases are fine in general (e.g.\ $a = 2$).

\textbf{\color{brand}Why the other options are incorrect.}
\begin{itemize}[leftmargin=2.4em, itemsep=1pt, topsep=2pt]
\item[\textbf{(A)}] $U_6$ maps $\ket0$ and $\ket7$ to $\ket0$: it is not unitary, so phase estimation cannot be applied.
\item[\textbf{(B)}] Order finding is not defined for $a$ not coprime to $N$: no power of $6$ is $\equiv 1 \pmod{21}$.
\item[\textbf{(D)}] Even bases are allowed; $a = 2$ works for $N = 21$. The point is the common factor $3$.
\end{itemize}
\textbf{\color{brand}Key takeaway.} Always compute $\gcd(a, N)$ first --- sometimes it already solves the problem.
\end{solbox}
\qsep

\needspace{13\baselineskip}
\begin{qitem}{46}
\qmeta{Single-correct MCQ}{Register sizes for $N = 33$}
With the register sizes of \texttt{Shor.ipynb} ($n = $ \texttt{N.bit\_length()}, $t = 2n$), the circuit for $N = 33$ uses
\par\smallskip\noindent\begin{tabularx}{\linewidth}{@{}*{2}{L}@{}}
\correctlabel{A}\,$18$ qubits ($n = 6$, $t = 12$) & \optlabel{B}\,$33$ qubits \\[3pt]
\optlabel{C}\,$12$ qubits & \optlabel{D}\,$15$ qubits ($n = 5$, $t = 10$)
\end{tabularx}
\end{qitem}
\begin{solbox}{A}
\textbf{\color{brand}Correct answer.} \textcolor{okgreen}{\textbf{(A)}} $18$ qubits ($n = 6$, $t = 12$)

\textbf{\color{brand}Concepts required.} Register sizes: $n = $ bit length of $N$, $t = 2n$ counting qubits, $3n$ in total.

\textbf{\color{brand}Relevant formulae.} $33 = 100001_2$: $n = 6$, $t = 12$, total $18$.

\textbf{\color{brand}Step-by-step explanation.} $33 = 100001_2$ has $6$ bits, so $n = 6$, $t = 12$ and $t + n = 18$. Five bits only reach $31 < 33$, so $n = 5$ is too small. The count is $3n$ in this construction
--- far from the $N$ qubits of option 4.

\textbf{\color{brand}Why the other options are incorrect.}
\begin{itemize}[leftmargin=2.4em, itemsep=1pt, topsep=2pt]
\item[\textbf{(B)}] The qubit count grows with $\log_2N$, not with $N$; $33$ qubits would be $3\times11$ bits.
\item[\textbf{(C)}] $12$ is only the counting register $t = 2n$.
\item[\textbf{(D)}] $n = 5$ bits only represent $0, \dots, 31$, which does not contain all residues modulo $33$.
\end{itemize}
\textbf{\color{brand}Key takeaway.} $3n$ qubits in the notebook's construction; better constructions need about $2n$.
\end{solbox}
\qsep

\needspace{13\baselineskip}
\begin{qitem}{47}
\qmeta{Single-correct MCQ}{$N = 15$, $a = 11$ outcomes}
For $N = 15$, $a = 11$ and $t = 8$ counting qubits, the ideal outcome distribution and the result of post-processing $y = 128$ are
\begin{itemize}[leftmargin=2.6em, labelsep=0.4em, itemsep=2pt, topsep=3pt]
\item[\optlabel{A}] peaks at $y \in \{0, 128\}$; $r = 2$, but $11 \equiv -1$, so the factors are trivial
\item[\optlabel{B}] a uniform distribution, because $11$ is prime
\item[\correctlabel{C}] peaks at $y \in \{0, 128\}$ with probability $1/2$ each; $128/256 = 1/2 \Rightarrow r = 2$, $\gcd(10, 15) = 5$, $\gcd(12, 15) = 3$
\item[\optlabel{D}] peaks at $y \in \{0, 64, 128, 192\}$; $r = 4$
\end{itemize}
\end{qitem}
\begin{solbox}{C}
\textbf{\color{brand}Correct answer.} \textcolor{okgreen}{\textbf{(C)}} peaks at $y \in \{0, 128\}$ with probability $1/2$ each; $128/256 = 1/2 \Rightarrow r = 2$, $\gcd(10, 15) = 5$, $\gcd(12, 15) = 3$

\textbf{\color{brand}Concepts required.} For $r \mid 2^t$ the outcomes are exactly $k2^t/r$, each with probability $1/r$.

\textbf{\color{brand}Relevant formulae.} $11^2 \equiv 1$: $r = 2$, peaks at $0$ and $128$; $\gcd(11 \mp 1, 15) = 5, 3$.

\textbf{\color{brand}Step-by-step explanation.} $11^2 = 121 = 8\cdot15 + 1$, so $r = 2$ and the peaks are at $k\cdot256/2 = 0, 128$, each with probability $\frac12$ (exact, since $r$ divides $2^t$). $y = 128$
gives $s/r = 1/2$, $r = 2$; $11 \not\equiv -1 \equiv 14$, and $\gcd(10, 15) = 5$, $\gcd(12, 15) = 3$. $y = 0$ gives no information.

\textbf{\color{brand}Why the other options are incorrect.}
\begin{itemize}[leftmargin=2.4em, itemsep=1pt, topsep=2pt]
\item[\textbf{(A)}] $11 \not\equiv -1 \equiv 14 \pmod{15}$; the gcds $5$ and $3$ are non-trivial.
\item[\textbf{(B)}] Primality of the base is irrelevant; the distribution depends only on the order $r$.
\item[\textbf{(D)}] Four peaks would mean $r = 4$; but $11^2 = 121 \equiv 1$, so $r = 2$.
\end{itemize}
\textbf{\color{brand}Key takeaway.} The number of peaks is the order; their spacing is $2^t/r$.
\end{solbox}
\qsep

\needspace{13\baselineskip}
\begin{qitem}{48}
\qmeta{Single-correct MCQ}{Continued fractions, $N = 33$}
For $N = 33$, $a = 5$ and $t = 12$, the measurement gives $y = 1229$. The continued-fraction step and the factors are
\begin{itemize}[leftmargin=2.6em, labelsep=0.4em, itemsep=2pt, topsep=3pt]
\item[\optlabel{A}] $1229/4096 \approx 1/3 \Rightarrow r = 3$; the attempt fails
\item[\optlabel{B}] $1229/4096 \approx 3/10 \Rightarrow r = 3$; factors $3$ and $11$
\item[\optlabel{C}] $1229/4096 \approx 3/10 \Rightarrow r = 10$; $5^5 \equiv -1$, trivial factors
\item[\correctlabel{D}] $1229/4096 \approx 3/10 \Rightarrow r = 10$; $5^5 \equiv 23$, $\gcd(22, 33) = 11$, $\gcd(24, 33) = 3$
\end{itemize}
\end{qitem}
\begin{solbox}{D}
\textbf{\color{brand}Correct answer.} \textcolor{okgreen}{\textbf{(D)}} $1229/4096 \approx 3/10 \Rightarrow r = 10$; $5^5 \equiv 23$, $\gcd(22, 33) = 11$, $\gcd(24, 33) = 3$

\textbf{\color{brand}Concepts required.} Continued fractions find the best approximation $s/r$ with $r \le N$ to $y/2^t$.

\textbf{\color{brand}Relevant formulae.} $1229/4096 = 0.30005 \approx 3/10$; $5^{10} \equiv 1$, $5^5 \equiv 23$; $\gcd(22, 33) = 11$, $\gcd(24, 33) = 3$.

\textbf{\color{brand}Step-by-step explanation.} $1229/4096 = 0.30005$; \texttt{limit\_denominator(33)} gives $3/10$ (within $\frac{1}{2\cdot4096}$). $5^{10} \equiv 1 \pmod{33}$ and no smaller power is $1$, so $r = 10$.
$5^5 = 3125 = 94\cdot33 + 23$; $23 \ne 32$; $\gcd(22, 33) = 11$, $\gcd(24, 33) = 3$. The numerator $3$ is $s$, not the order.

\textbf{\color{brand}Why the other options are incorrect.}
\begin{itemize}[leftmargin=2.4em, itemsep=1pt, topsep=2pt]
\item[\textbf{(A)}] $1/3 = 0.333$ is much farther from $0.30005$ than $3/10$; continued fractions give $3/10$.
\item[\textbf{(B)}] The numerator $3$ is the random index $s$; the order is the denominator $10$.
\item[\textbf{(C)}] $5^5 \equiv 23$, and $-1 \equiv 32$; so the factors are non-trivial.
\end{itemize}
\textbf{\color{brand}Key takeaway.} Read the order from the denominator, then verify $a^r \equiv 1$.
\end{solbox}
\qsep

\needspace{13\baselineskip}
\begin{qitem}{49}
\qmeta{Single-correct MCQ}{Multiples test}
For $N = 33$, $a = 5$, $t = 12$, the outcome is $y = 2048$. The procedure of \texttt{candidate\_orders} returns
\begin{itemize}[leftmargin=2.6em, labelsep=0.4em, itemsep=2pt, topsep=3pt]
\item[\optlabel{A}] the order $2$, because the denominator is always the order
\item[\correctlabel{B}] the fraction $1/2$ and, after testing multiples of $2$, the order $10$ (since $5^2, 5^4, 5^6, 5^8 \not\equiv 1$, $5^{10} \equiv 1$)
\item[\optlabel{C}] no candidate, because $y = 2^{t-1}$ is never useful
\item[\optlabel{D}] the order $5$
\end{itemize}
\end{qitem}
\begin{solbox}{B}
\textbf{\color{brand}Correct answer.} \textcolor{okgreen}{\textbf{(B)}} the fraction $1/2$ and, after testing multiples of $2$, the order $10$ (since $5^2, 5^4, 5^6, 5^8 \not\equiv 1$, $5^{10} \equiv 1$)

\textbf{\color{brand}Concepts required.} If $\gcd(s, r) > 1$ the reduced fraction has a denominator that divides $r$; testing multiples recovers $r$.

\textbf{\color{brand}Relevant formulae.} $2048/4096 = 1/2 = 5/10$; test $2k$: $5^{10} \equiv 1$ first at $k = 5$.

\textbf{\color{brand}Step-by-step explanation.} $2048/4096 = 1/2$: the measured $s/r = 5/10$ has $\gcd(s, r) = 5$, so it reduces to $1/2$. Testing $2k$: $5^2 \equiv 25$, $5^4 \equiv 31$, $5^6 \equiv 16$, $5^8 \equiv 4$,
$5^{10} \equiv 1$ --- the first multiple that works is $10$. Without the multiples test this outcome would be lost.

\textbf{\color{brand}Why the other options are incorrect.}
\begin{itemize}[leftmargin=2.4em, itemsep=1pt, topsep=2pt]
\item[\textbf{(A)}] The denominator is the order only if $\gcd(s, r) = 1$; here $s = 5$, $r = 10$ share the factor $5$.
\item[\textbf{(C)}] $y = 2^{t-1}$ is useful whenever $r$ is even: it gives $1/2$, and the multiples test finds $r$.
\item[\textbf{(D)}] $5^5 \equiv 23 \ne 1$, so $5$ is not the order; also $5$ is not a multiple of the denominator $2$.
\end{itemize}
\textbf{\color{brand}Key takeaway.} The multiples test turns many ``partial'' outcomes into successes.
\end{solbox}
\qsep

\needspace{13\baselineskip}
\begin{qitem}{50}
\qmeta{Single-correct MCQ}{Condition $2^t \ge N^2$}
The textbook condition $2^t \ge N^2$ guarantees that continued fractions recover $s/r$ from the nearest outcome. For $N = 33$, the smallest such
$t$ and the value used by \texttt{Shor.ipynb} are
\par\smallskip\noindent\begin{tabularx}{\linewidth}{@{}*{2}{L}@{}}
\optlabel{A}\,$t = 12$; the notebook uses $11$ & \correctlabel{B}\,$t = 11$; the notebook uses $t = 2n = 12$ \\[3pt]
\optlabel{C}\,$t = 10$; the notebook uses $12$ & \optlabel{D}\,$t = 6$; the notebook uses $12$
\end{tabularx}
\end{qitem}
\begin{solbox}{B}
\textbf{\color{brand}Correct answer.} \textcolor{okgreen}{\textbf{(B)}} $t = 11$; the notebook uses $t = 2n = 12$

\textbf{\color{brand}Concepts required.} With $2^t \ge N^2$, the outcome nearest $s2^t/r$ lies within $\frac{1}{2r^2}$ of $s/r$, so continued fractions find it.

\textbf{\color{brand}Relevant formulae.} $N^2 = 1089$; $2^{10} = 1024 < 1089 \le 2048 = 2^{11}$.

\textbf{\color{brand}Step-by-step explanation.} $N^2 = 1089$: $2^{10} = 1024 < 1089 \le 2^{11} = 2048$, so $t = 11$. The notebook uses $t = 2\cdot$\texttt{bit\_length} $= 12$, which always satisfies $2^{2n} \ge N^2$
because $N < 2^n$. Then two distinct fractions with denominators $\le N$ differ by more than $1/N^2 \ge 1/2^t$.

\textbf{\color{brand}Why the other options are incorrect.}
\begin{itemize}[leftmargin=2.4em, itemsep=1pt, topsep=2pt]
\item[\textbf{(A)}] This swaps the two values: the minimum is $11$, and the notebook's $t = 2n$ gives $12$.
\item[\textbf{(C)}] $2^{10} = 1024 < 1089 = N^2$, so $t = 10$ is too small.
\item[\textbf{(D)}] $t = 6$ is the work register size $n$; $2^6 = 64 \ll N^2$.
\end{itemize}
\textbf{\color{brand}Key takeaway.} $t = 2n$ always satisfies $2^t \ge N^2$, with up to one qubit to spare.
\end{solbox}
\qsep

\needspace{13\baselineskip}
\begin{qitem}{51}
\qmeta{Single-correct MCQ}{Exact success, notebook default}
For the notebook's default ($N = 15$, $a = 7$, $t = 8$) the ideal outcomes are $y \in \{0, 64, 128, 192\}$, each with probability $1/4$. Using the
multiples test, the exact probability that one shot yields the factors is
\par\smallskip\noindent\begin{tabularx}{\linewidth}{@{}*{2}{L}@{}}
\optlabel{A}\,$1/4$ --- only $y = 64$ succeeds & \optlabel{B}\,$1$ \\[3pt]
\optlabel{C}\,$1/2$ --- only $y = 64$ and $192$ succeed & \correctlabel{D}\,$3/4$ --- only $y = 0$ fails
\end{tabularx}
\end{qitem}
\begin{solbox}{D}
\textbf{\color{brand}Correct answer.} \textcolor{okgreen}{\textbf{(D)}} $3/4$ --- only $y = 0$ fails

\textbf{\color{brand}Concepts required.} Count the outcomes whose post-processing yields the factors, weighted by their probabilities.

\textbf{\color{brand}Relevant formulae.} $y = 64, 192 \to r = 4$; $y = 128 \to 1/2 \to$ multiples $\to 4$; $y = 0 \to$ nothing: $3/4$.

\textbf{\color{brand}Step-by-step explanation.} $y = 64$, $192$ give $1/4$, $3/4$, i.e.\ $r = 4$ directly; $y = 128$ gives $1/2$, and the multiples test finds $4$; $y = 0$ gives no denominator. $7^2 = 49 \equiv 4$,
$\gcd(3, 15) = 3$, $\gcd(5, 15) = 5$. Hence $3/4$; the recorded run gave $0.766$ (finite-shot fluctuation around $0.75$).

\textbf{\color{brand}Why the other options are incorrect.}
\begin{itemize}[leftmargin=2.4em, itemsep=1pt, topsep=2pt]
\item[\textbf{(A)}] $1/4$ counts one successful outcome; three of the four peaks succeed.
\item[\textbf{(B)}] Success probability $1$ is impossible: the outcome $y = 0$ occurs with probability $1/4$ and carries no information.
\item[\textbf{(C)}] $1/2$ forgets that $y = 128$ succeeds through the multiples test.
\end{itemize}
\textbf{\color{brand}Key takeaway.} Recorded $0.766$ vs exact $0.75$: the difference is shot noise, not a better algorithm.
\end{solbox}
\qsep

\needspace{13\baselineskip}
\begin{qitem}{52}
\qmeta{Single-correct MCQ}{Probability $\varphi(r)/r$}
Suppose $r = 12$ divides $2^t$, so every outcome is exactly $y = s2^t/r$ with $s$ uniform in $\{0, \dots, 11\}$. \emph{Without} the multiples test,
the fraction of shots whose reduced denominator equals $r$ is
\par\smallskip\noindent\begin{tabularx}{\linewidth}{@{}*{4}{L}@{}}
\correctlabel{A}\,$\varphi(12)/12 = 1/3$ & \optlabel{B}\,$1/2$ & \optlabel{C}\,$11/12$ & \optlabel{D}\,$1/12$
\end{tabularx}
\end{qitem}
\begin{solbox}{A}
\textbf{\color{brand}Correct answer.} \textcolor{okgreen}{\textbf{(A)}} $\varphi(12)/12 = 1/3$

\textbf{\color{brand}Concepts required.} The reduced fraction $s/r$ keeps the denominator $r$ iff $\gcd(s, r) = 1$.

\textbf{\color{brand}Relevant formulae.} $P = \varphi(r)/r$; $\varphi(12) = |\{1, 5, 7, 11\}| = 4$, $4/12 = 1/3$.

\textbf{\color{brand}Step-by-step explanation.} The reduced fraction $s/12$ has denominator $12$ exactly when $\gcd(s, 12) = 1$: $s \in \{1, 5, 7, 11\}$, i.e.\ $\varphi(12) = 4$ of $12$ values. In general the
probability is $\varphi(r)/r = \Omega(1/\log\log r)$, so $O(\log\log N)$ repetitions (or the multiples/lcm trick) suffice.

\textbf{\color{brand}Why the other options are incorrect.}
\begin{itemize}[leftmargin=2.4em, itemsep=1pt, topsep=2pt]
\item[\textbf{(B)}] $1/2$ is the fraction of odd $s$; $s = 3, 9$ are odd but share the factor $3$ with $12$.
\item[\textbf{(C)}] $11/12$ counts every $s \ne 0$; most of them share a factor with $12$.
\item[\textbf{(D)}] $1/12$ counts only $s = 1$; $s = 5, 7, 11$ are also coprime to $12$.
\end{itemize}
\textbf{\color{brand}Key takeaway.} $\varphi(r)/r = \Omega(1/\log\log r)$: a few repetitions suffice.
\end{solbox}
\qsep

\needspace{13\baselineskip}
\begin{qitem}{53}
\qmeta{Single-correct MCQ}{Role of the QFT in Shor}
What is the role of the (inverse) quantum Fourier transform \emph{within} Shor's algorithm?
\begin{itemize}[leftmargin=2.6em, labelsep=0.4em, itemsep=2pt, topsep=3pt]
\item[\correctlabel{A}] it is a subroutine of phase estimation that turns the periodic phase pattern on the counting register into peaks near $y = k2^t/r$; it does not factor $N$ by itself
\item[\optlabel{B}] it finds the factors $p$ and $q$ directly from $\ket N$
\item[\optlabel{C}] it computes $a^x \bmod N$
\item[\optlabel{D}] it is the whole of Shor's algorithm; the remaining steps are optional
\end{itemize}
\end{qitem}
\begin{solbox}{A}
\textbf{\color{brand}Correct answer.} \textcolor{okgreen}{\textbf{(A)}} it is a subroutine of phase estimation that turns the periodic phase pattern on the counting register into peaks near $y = k2^t/r$; it does not factor $N$ by itself

\textbf{\color{brand}Concepts required.} The QFT is a component of phase estimation; Shor's algorithm is phase estimation of $U_a$ plus classical processing.

\textbf{\color{brand}Relevant formulae.} $\mathrm{QFT}^{-1}: \frac{1}{\sqrt{2^t}}\sum_xe^{2\pi ix\varphi}\ket x \mapsto \approx\ket{2^t\varphi}$.

\textbf{\color{brand}Step-by-step explanation.} Shor's algorithm $=$ classical reduction $+$ order finding (modular exponentiation, QFT, measurement) $+$ classical post-processing. The QFT is one
component; the dominant cost of real implementations is the modular exponentiation, not the QFT ($O(n^2)$ gates, or fewer when approximated).

\textbf{\color{brand}Why the other options are incorrect.}
\begin{itemize}[leftmargin=2.4em, itemsep=1pt, topsep=2pt]
\item[\textbf{(B)}] No circuit acts on $\ket N$; $N$ enters only through the classical choice of the multipliers.
\item[\textbf{(C)}] $a^x \bmod N$ is computed by the controlled modular multiplications, not by the QFT.
\item[\textbf{(D)}] The QFT alone factors nothing; modular exponentiation and the classical steps are essential.
\end{itemize}
\textbf{\color{brand}Key takeaway.} Never confuse Shor with the QFT: the QFT is one subroutine.
\end{solbox}
\qsep

\needspace{13\baselineskip}
\begin{qitem}{54}
\qmeta{Single-correct MCQ}{Skipped identity multipliers}
\texttt{build\_shor\_kernel} computes $c_j = a^{2^j} \bmod N$ classically and skips $c_j = 1$. For $N = 15$, $a = 7$, $t = 8$, the number of controlled
multiplications actually applied is
\par\smallskip\noindent\begin{tabularx}{\linewidth}{@{}*{2}{L}@{}}
\optlabel{A}\,$4$ & \optlabel{B}\,$1$ \\[3pt]
\correctlabel{C}\,$2$ ($c_0 = 7$, $c_1 = 4$; $c_j = 1$ for $j \ge 2$) & \optlabel{D}\,$8$
\end{tabularx}
\end{qitem}
\begin{solbox}{C}
\textbf{\color{brand}Correct answer.} \textcolor{okgreen}{\textbf{(C)}} $2$ ($c_0 = 7$, $c_1 = 4$; $c_j = 1$ for $j \ge 2$)

\textbf{\color{brand}Concepts required.} The multiplier constants $c_j = a^{2^j} \bmod N$ are computed classically; $c_j = 1$ is the identity.

\textbf{\color{brand}Relevant formulae.} $7, 49 \equiv 4, 16 \equiv 1, 1, \dots$: two non-identity multipliers.

\textbf{\color{brand}Step-by-step explanation.} $7^1 = 7$, $7^2 = 49 \equiv 4$, $7^4 \equiv 16 \equiv 1$, and then $1$ forever. Only two multipliers remain --- matching the recorded gate counts
(\texttt{ccustom\_unitary\_constant}: 2; operations \texttt{mulmod\_7\_mod\_15}, \texttt{mulmod\_4\_mod\_15}). This shortcut uses the known order and is
specific to the tiny example.

\textbf{\color{brand}Why the other options are incorrect.}
\begin{itemize}[leftmargin=2.4em, itemsep=1pt, topsep=2pt]
\item[\textbf{(A)}] $4$ is the order $r$, not the number of multipliers.
\item[\textbf{(B)}] $1$ forgets that $c_1 = 4 \ne 1$ (it appears as \texttt{mulmod\_4\_mod\_15}).
\item[\textbf{(D)}] $8$ would be the count without skipping the identity multipliers; the notebook skips $c_j = 1$.
\end{itemize}
\textbf{\color{brand}Key takeaway.} Skipping identity multipliers exploits $r \mid 2^j$ --- a shortcut available only for toy examples.
\end{solbox}
\qsep

\needspace{13\baselineskip}
\begin{qitem}{55}
\qmeta{Matching lists}{Functions of \texttt{Shor.ipynb}}
Match each function of \texttt{Shor.ipynb} (List-I) with what it does (List-II).
\par\smallskip\noindent{\renewcommand{\arraystretch}{1.35}%
\begin{tabularx}{\linewidth}{@{}c>{\raggedright\arraybackslash}p{0.33\linewidth}@{\hspace{1.2em}}cL@{}}
\toprule
\multicolumn{2}{@{}l}{\textbf{List-I}} & \multicolumn{2}{l@{}}{\textbf{List-II}} \\ \midrule
\textbf{(P)} & \texttt{mulmod\_matrix} & \textbf{(1)} & brute-force order of $a$ modulo $m$ as a classical reference (feasible only for tiny $m$) \\
\textbf{(Q)} & \texttt{build\_shor\_kernel} & \textbf{(2)} & continued-fraction denominator of $y/2^t$ and the smallest multiple $k r'$ with $a^{kr'} \equiv 1$ \\
\textbf{(R)} & \texttt{candidate\_orders} & \textbf{(3)} & builds the CUDA-Q circuit: Hadamards, controlled $U_{a^{2^j}}$, inverse QFT, measurement \\
\textbf{(S)} & \texttt{classical\_order} & \textbf{(4)} & permutation matrix $\ket y \mapsto \ket{cy \bmod N}$ for $y < N$, identity above \\
\bottomrule
\end{tabularx}}
\par\smallskip\noindent\begin{tabularx}{\linewidth}{@{}*{2}{L}@{}}
\optlabel{A}\,P$\to$3, Q$\to$4, R$\to$2, S$\to$1 & \correctlabel{B}\,P$\to$4, Q$\to$3, R$\to$2, S$\to$1 \\[3pt]
\optlabel{C}\,P$\to$4, Q$\to$3, R$\to$1, S$\to$2 & \optlabel{D}\,P$\to$4, Q$\to$2, R$\to$3, S$\to$1
\end{tabularx}
\end{qitem}
\begin{solbox}{B}
\textbf{\color{brand}Correct answer.} \textcolor{okgreen}{\textbf{(B)}} P$\to$4, Q$\to$3, R$\to$2, S$\to$1

\textbf{\color{brand}Concepts required.} The notebook separates matrix construction, circuit construction, post-processing and classical checks.

\textbf{\color{brand}Relevant formulae.} \texttt{mulmod\_matrix} $\to$ \texttt{cudaq.register\_operation} $\to$ \texttt{build\_shor\_kernel} $\to$ \texttt{cudaq.sample} $\to$ \texttt{candidate\_orders}.

\textbf{\color{brand}Step-by-step explanation.} \texttt{mulmod\_matrix} builds the unitary (4); \texttt{build\_shor\_kernel} assembles the circuit with \texttt{cudaq.register\_operation} gates (3);
\texttt{candidate\_orders} is the classical post-processing (2); \texttt{classical\_order} is the brute-force check (1). Hence P$\to$4, Q$\to$3, R$\to$2, S$\to$1.

\textbf{\color{brand}Why the other options are incorrect.}
\begin{itemize}[leftmargin=2.4em, itemsep=1pt, topsep=2pt]
\item[\textbf{(A)}] This swaps the matrix and the kernel builder: \texttt{mulmod\_matrix} returns a NumPy permutation matrix, not a circuit.
\item[\textbf{(C)}] This swaps post-processing and the classical reference: \texttt{candidate\_orders} uses $y$ and continued fractions; \texttt{classical\_order} uses brute force.
\item[\textbf{(D)}] This swaps the kernel builder and the post-processing: \texttt{build\_shor\_kernel} builds the circuit and never sees $y$.
\end{itemize}
\textbf{\color{brand}Key takeaway.} Know which part of the code is quantum (the kernel) and which is classical (everything around it).
\end{solbox}
\qsep

\needspace{13\baselineskip}
\begin{qitem}{56}
\qmeta{Matching lists}{Outcomes to fractions}
For $N = 33$, $a = 5$ ($r = 10$) and $t = 12$, match each outcome $y$ (List-I) with \texttt{Fraction(y, 4096).limit\_denominator(33)} (List-II).
\par\smallskip\noindent{\renewcommand{\arraystretch}{1.35}%
\begin{tabularx}{\linewidth}{@{}c>{\raggedright\arraybackslash}p{0.33\linewidth}@{\hspace{1.2em}}cL@{}}
\toprule
\multicolumn{2}{@{}l}{\textbf{List-I}} & \multicolumn{2}{l@{}}{\textbf{List-II}} \\ \midrule
\textbf{(P)} & $410$ & \textbf{(1)} & $2/5$ \\
\textbf{(Q)} & $819$ & \textbf{(2)} & $7/10$ \\
\textbf{(R)} & $1638$ & \textbf{(3)} & $1/10$ \\
\textbf{(S)} & $2867$ & \textbf{(4)} & $1/5$ \\
\bottomrule
\end{tabularx}}
\par\smallskip\noindent\begin{tabularx}{\linewidth}{@{}*{2}{L}@{}}
\optlabel{A}\,P$\to$3, Q$\to$4, R$\to$2, S$\to$1 & \optlabel{B}\,P$\to$3, Q$\to$1, R$\to$4, S$\to$2 \\[3pt]
\optlabel{C}\,P$\to$4, Q$\to$3, R$\to$1, S$\to$2 & \correctlabel{D}\,P$\to$3, Q$\to$4, R$\to$1, S$\to$2
\end{tabularx}
\end{qitem}
\begin{solbox}{D}
\textbf{\color{brand}Correct answer.} \textcolor{okgreen}{\textbf{(D)}} P$\to$3, Q$\to$4, R$\to$1, S$\to$2

\textbf{\color{brand}Concepts required.} \texttt{limit\_denominator(N)} returns the closest fraction with denominator at most $N$.

\textbf{\color{brand}Relevant formulae.} $410/4096 = 0.1001$, $819/4096 = 0.19995$, $1638/4096 = 0.39990$, $2867/4096 = 0.69995$.

\textbf{\color{brand}Step-by-step explanation.} $410/4096 = 0.1001 \approx 1/10$; $819/4096 = 0.19995 \approx 1/5$; $1638/4096 = 0.3999 \approx 2/5$; $2867/4096 = 0.69995 \approx 7/10$. Only $410$ and $2867$ give the
denominator $10$ directly ($s = 1, 7$ coprime to $10$); $819$ and $1638$ give $5$, and the multiples test then finds $10$.

\textbf{\color{brand}Why the other options are incorrect.}
\begin{itemize}[leftmargin=2.4em, itemsep=1pt, topsep=2pt]
\item[\textbf{(A)}] This exchanges $2/5$ and $7/10$: $1638/4096 \approx 0.4$ and $2867/4096 \approx 0.7$.
\item[\textbf{(B)}] This assigns $1/10$ to $1638$; but $1638/4096 \approx 0.4 = 2/5$.
\item[\textbf{(C)}] This exchanges $1/10$ and $1/5$: $410/4096 \approx 0.1$ and $819/4096 \approx 0.2$.
\end{itemize}
\textbf{\color{brand}Key takeaway.} Only fractions with $\gcd(s, r) = 1$ give $r$ directly; others need the multiples test.
\end{solbox}
\qsep

\needspace{13\baselineskip}
\begin{qitem}{57}
\qmeta{Assertion--Reason}{Prime powers}
\textbf{Assertion (A):} Shor's order-finding reduction is the method of choice for factoring $N = 81 = 3^4$.\\[1mm]
\textbf{Reason (R):} For an odd prime power $N = p^k$, the only square roots of $1$ modulo $N$ are $\pm1$.
\begin{itemize}[leftmargin=2.6em, labelsep=0.4em, itemsep=2pt, topsep=3pt]
\item[\optlabel{A}] Both Assertion (A) and Reason (R) are true, and R is the correct explanation of A
\item[\optlabel{B}] Both A and R are true, but R is \emph{not} the correct explanation of A
\item[\optlabel{C}] A is true, but R is false
\item[\correctlabel{D}] A is false, but R is true
\end{itemize}
\end{qitem}
\begin{solbox}{D}
\textbf{\color{brand}Correct answer.} \textcolor{okgreen}{\textbf{(D)}} A is false, but R is true

\textbf{\color{brand}Concepts required.} Prime powers are excluded from the quantum reduction and handled classically.

\textbf{\color{brand}Relevant formulae.} $\mathbb Z_{p^k}^*$ cyclic $\Rightarrow x^2 \equiv 1 \iff x \equiv \pm1$; $81 = 3^4$ found by integer roots.

\textbf{\color{brand}Step-by-step explanation.} \textbf{A is false}: prime powers are detected classically (integer $k$-th roots) --- \texttt{Shor.ipynb} does so with \texttt{prime\_power} before any quantum step.
\textbf{R is true}: $\mathbb Z_{p^k}^*$ is cyclic, so $x^2 \equiv 1$ has only the solutions $\pm1$; hence $a^{r/2}$ would always be $\pm1$ and $\gcd(a^{r/2} \pm 1, N)$ would
be trivial. Answer: (D).

\textbf{\color{brand}Why the other options are incorrect.}
\begin{itemize}[leftmargin=2.4em, itemsep=1pt, topsep=2pt]
\item[\textbf{(A)}] is wrong because A is false: prime powers are detected classically (\texttt{prime\_power} in the notebook).
\item[\textbf{(B)}] is wrong because A is false, so it cannot be a true assertion with an unrelated reason.
\item[\textbf{(C)}] is wrong because A is false and R is true, the opposite of (C).
\end{itemize}
\textbf{\color{brand}Key takeaway.} Shor needs at least two distinct prime factors; check prime powers classically first.
\end{solbox}
\qsep

\needspace{13\baselineskip}
\begin{qitem}{58}
\qmeta{Statement I / II}{Uninformative outcomes}
\textbf{Statement I:} A measured value $y = 0$ means that the order is $r = 1$.\\[1mm]
\textbf{Statement II:} When $r$ does not divide $2^t$, every measured $y$ yields the exact order through continued fractions.
\begin{itemize}[leftmargin=2.6em, labelsep=0.4em, itemsep=2pt, topsep=3pt]
\item[\optlabel{A}] Both Statement I and Statement II are true
\item[\correctlabel{B}] Both Statement I and Statement II are false
\item[\optlabel{C}] Statement I is true, but Statement II is false
\item[\optlabel{D}] Statement I is false, but Statement II is true
\end{itemize}
\end{qitem}
\begin{solbox}{B}
\textbf{\color{brand}Correct answer.} \textcolor{okgreen}{\textbf{(B)}} Both Statement I and Statement II are false

\textbf{\color{brand}Concepts required.} Measurement outcomes are random; only some of them reveal the order directly.

\textbf{\color{brand}Relevant formulae.} $y = 0 \leftrightarrow s = 0$ (always present); success requires $y \approx s2^t/r$ with $\gcd(s, r) = 1$.

\textbf{\color{brand}Step-by-step explanation.} \textbf{I is false}: $y = 0$ is the peak $s = 0$, present for every $r$; it carries no information (the notebook returns no candidate). \textbf{II is false}: only
outcomes close to $s2^t/r$ with $\gcd(s, r) = 1$ give $r$ directly; others give a divisor of $r$ or nothing, and the probability mass is spread over neighbouring $y$.
Answer: (B).

\textbf{\color{brand}Why the other options are incorrect.}
\begin{itemize}[leftmargin=2.4em, itemsep=1pt, topsep=2pt]
\item[\textbf{(A)}] is wrong because both statements are false.
\item[\textbf{(C)}] is wrong because Statement I is false: $y = 0$ occurs for every order and gives no denominator.
\item[\textbf{(D)}] is wrong because Statement II is false: for $N = 21$, $a = 2$ the outcome $y = 512$ gives $1/2$, and outcomes between peaks give nothing useful.
\end{itemize}
\textbf{\color{brand}Key takeaway.} Shor's algorithm is probabilistic; failures are expected and handled by repetition.
\end{solbox}
\qsep

\needspace{20\baselineskip}
\begin{passagebox}{Paragraph for Questions 59 and 60}
Shor's algorithm is applied to $N = 39$ with base $a = 2$, using the register sizes of \texttt{Shor.ipynb}: the work register has
$n = $ \texttt{N.bit\_length()} $= 6$ qubits and the counting register $t = 2n = 12$ qubits, so the peaks of the outcome distribution lie near
$y \approx k\cdot4096/r$. A measured value $y$ is post-processed with \texttt{Fraction(y, 2**t).limit\_denominator(N)}; the denominator $r'$ of the
result, or the smallest multiple $kr'$ with $a^{kr'} \equiv 1 \pmod N$, is taken as the order. The powers of $2$ modulo $39$ are
$2, 4, 8, 16, 32, 25, 11, 22, 5, 10, 20, 1$.
\end{passagebox}
\needspace{13\baselineskip}
\begin{qitem}{59}
\qmeta{Paragraph based}{$N = 39$: order and factors}
The order $r$ and the factors $\gcd(2^{r/2} - 1, 39)$, $\gcd(2^{r/2} + 1, 39)$ are
\par\smallskip\noindent\begin{tabularx}{\linewidth}{@{}*{2}{L}@{}}
\correctlabel{A}\,$r = 12$; $\ 3$ and $13$ & \optlabel{B}\,$r = 6$; $\ 3$ and $13$ \\[3pt]
\optlabel{C}\,$r = 12$; $\ 7$ and $13$ & \optlabel{D}\,$r = 12$; $\ 1$ and $39$
\end{tabularx}
\end{qitem}
\begin{solbox}{A}
\textbf{\color{brand}Correct answer.} \textcolor{okgreen}{\textbf{(A)}} $r = 12$; $\ 3$ and $13$

\textbf{\color{brand}Concepts required.} The order is the position of the first $1$ in the list of powers.

\textbf{\color{brand}Relevant formulae.} $2^{12} \equiv 1 \pmod{39}$, $2^6 \equiv 25$; $\gcd(24, 39) = 3$, $\gcd(26, 39) = 13$.

\textbf{\color{brand}Step-by-step explanation.} From the list, $2^{12} \equiv 1$ is the first $1$: $r = 12$, even. $2^6 \equiv 25 \not\equiv 38$. $\gcd(24, 39) = 3$ and $\gcd(26, 39) = 13$; $3\times13 = 39$.
$r = 6$ would require $2^6 \equiv 1$, but $2^6 \equiv 25$.

\textbf{\color{brand}Why the other options are incorrect.}
\begin{itemize}[leftmargin=2.4em, itemsep=1pt, topsep=2pt]
\item[\textbf{(B)}] $2^6 \equiv 25 \ne 1$, so $r \ne 6$ (and the factors would have to come from $2^3$).
\item[\textbf{(C)}] $7$ does not divide $39 = 3\cdot13$.
\item[\textbf{(D)}] $25 \not\equiv -1 \equiv 38$, so the gcds are non-trivial.
\end{itemize}
\textbf{\color{brand}Key takeaway.} $N = 39$, $a = 2$: $r = 12$, $39 = 3 \times 13$.
\end{solbox}
\qsep

\needspace{13\baselineskip}
\begin{qitem}{60}
\qmeta{Paragraph based}{$N = 39$: continued fraction}
The outcome $y = 1707$ is measured. The continued-fraction step gives
\par\smallskip\noindent\begin{tabularx}{\linewidth}{@{}*{2}{L}@{}}
\optlabel{A}\,$5/12$, so $r = 5$ & \optlabel{B}\,$1/2$, so the multiples test is needed \\[3pt]
\correctlabel{C}\,$5/12$, so $r = 12$ directly & \optlabel{D}\,$2/5$, so $r = 5$
\end{tabularx}
\end{qitem}
\begin{solbox}{C}
\textbf{\color{brand}Correct answer.} \textcolor{okgreen}{\textbf{(C)}} $5/12$, so $r = 12$ directly

\textbf{\color{brand}Concepts required.} A continued-fraction convergent within $\frac{1}{2r^2}$ of $y/2^t$ is the reduced form of $s/r$.

\textbf{\color{brand}Relevant formulae.} $|1707/4096 - 5/12| = 8.1\times10^{-5} < 1/(2\cdot39^2) = 3.3\times10^{-4}$.

\textbf{\color{brand}Step-by-step explanation.} $1707/4096 = 0.41675$, and $5/12 = 0.41667$; $|1707/4096 - 5/12| = 8.1\times10^{-5} < 1/(2\cdot39^2)$, so \texttt{limit\_denominator(39)} returns $5/12$.
$\gcd(5, 12) = 1$, so the denominator is the order $12$. The numerator $5$ is the random index $s$, not the order.

\textbf{\color{brand}Why the other options are incorrect.}
\begin{itemize}[leftmargin=2.4em, itemsep=1pt, topsep=2pt]
\item[\textbf{(A)}] The numerator $5$ is the random index $s$; the order is the denominator $12$.
\item[\textbf{(B)}] $1/2 = 0.5$ is far from $0.41675$; that fraction belongs to $y = 2048$.
\item[\textbf{(D)}] $2/5 = 0.4$ differs by $0.017$, while $5/12$ differs by only $8\times10^{-5}$.
\end{itemize}
\textbf{\color{brand}Key takeaway.} Denominator $=$ order when $\gcd(s, r) = 1$; always verify $a^r \equiv 1$.
\end{solbox}
\qsep

\newpage

\begin{sectionbanner}SECTION D \enspace$\cdot$\enspace Advanced Shor's Algorithm and Implementation \hfill Q.61--80\end{sectionbanner}

\needspace{13\baselineskip}
\begin{qitem}{61}
\qmeta{Single-correct MCQ}{Beauregard qubit count}
Beauregard's circuit for Shor's algorithm uses $2n + 3$ logical qubits. For RSA-2048 ($n = 2048$) this is
\par\smallskip\noindent\begin{tabularx}{\linewidth}{@{}*{4}{L}@{}}
\optlabel{A}\,$6147$ & \correctlabel{B}\,$4099$ & \optlabel{C}\,$4096$ & \optlabel{D}\,$2051$
\end{tabularx}
\end{qitem}
\begin{solbox}{B}
\textbf{\color{brand}Correct answer.} \textcolor{okgreen}{\textbf{(B)}} $4099$

\textbf{\color{brand}Concepts required.} Logical-qubit counts of Shor circuits are linear in the bit length $n$.

\textbf{\color{brand}Relevant formulae.} Beauregard: $2n + 3$; $n = 2048$: $4099$.

\textbf{\color{brand}Step-by-step explanation.} $2\times2048 + 3 = 4099$. $6147 = 3\times2048 + 3$ resembles the $3n$ register count of \texttt{Shor.ipynb} ($6144$). These are \emph{logical}
(error-corrected) qubits; the physical count is orders of magnitude larger.

\textbf{\color{brand}Why the other options are incorrect.}
\begin{itemize}[leftmargin=2.4em, itemsep=1pt, topsep=2pt]
\item[\textbf{(A)}] $6147 = 3n + 3$; the $3n$ register count belongs to the notebook's construction, not Beauregard's.
\item[\textbf{(C)}] $4096 = 2n$ forgets the three extra qubits (carry/ancilla and the recycled control).
\item[\textbf{(D)}] $2051 = n + 3$ omits one of the two $n$-qubit registers (work register and adder register).
\end{itemize}
\textbf{\color{brand}Key takeaway.} Thousands of logical qubits --- each made of hundreds to thousands of physical qubits.
\end{solbox}
\qsep

\needspace{13\baselineskip}
\begin{qitem}{62}
\qmeta{Single-correct MCQ}{Published RSA-2048 estimates}
Which pair of published resource estimates for factoring RSA-2048 is correct?
\begin{itemize}[leftmargin=2.6em, labelsep=0.4em, itemsep=2pt, topsep=3pt]
\item[\optlabel{A}] Gidney (2025): about $20$ billion qubits in about $8$ hours
\item[\optlabel{B}] both estimates state that RSA-2048 was factored on existing hardware
\item[\correctlabel{C}] Gidney and Ekerå (2019): about $20$ million noisy qubits in about $8$ hours; Gidney (2025): fewer than one million noisy qubits in less than a week
\item[\optlabel{D}] Gidney and Ekerå (2019): $4099$ physical qubits in one second
\end{itemize}
\end{qitem}
\begin{solbox}{C}
\textbf{\color{brand}Correct answer.} \textcolor{okgreen}{\textbf{(C)}} Gidney and Ekerå (2019): about $20$ million noisy qubits in about $8$ hours; Gidney (2025): fewer than one million noisy qubits in less than a week

\textbf{\color{brand}Concepts required.} Published resource estimates are model-based projections, not experimental results.

\textbf{\color{brand}Relevant formulae.} GE2019: $\approx2\times10^7$ physical qubits, $\approx8$ h; Gidney 2025: $<10^6$ physical qubits, $<1$ week.

\textbf{\color{brand}Step-by-step explanation.} Both are \emph{estimates} under explicit assumptions (surface code, physical error rate $10^{-3}$, $\mu$s cycle times), not experiments. Gidney--Ekerå 2019:
$\approx 20$ million physical qubits, $\approx 8$ hours. Gidney 2025 (approximate residue arithmetic, yoked surface codes, magic-state cultivation): $< 1$
million qubits, $< 1$ week. No cryptographically relevant number has been factored by Shor's algorithm.

\textbf{\color{brand}Why the other options are incorrect.}
\begin{itemize}[leftmargin=2.4em, itemsep=1pt, topsep=2pt]
\item[\textbf{(A)}] The 2025 estimate \emph{lowered} the qubit count to under one million; $20$ billion is fabricated.
\item[\textbf{(B)}] No resource estimate claims an actual factorisation of RSA-2048; none has been performed.
\item[\textbf{(D)}] $4099$ is Beauregard's \emph{logical} count; GE2019 required about $20$ million physical qubits and hours, not seconds.
\end{itemize}
\textbf{\color{brand}Key takeaway.} Distinguish estimates from demonstrations: RSA-2048 remains unfactored.
\end{solbox}
\qsep

\needspace{13\baselineskip}
\begin{qitem}{63}
\qmeta{Single-correct MCQ}{$O(n^3)$ scaling}
With schoolbook arithmetic, modular exponentiation costs $O(n^3)$ gates. Doubling the bit length $n$ of $N$ multiplies the gate count by about
\par\smallskip\noindent\begin{tabularx}{\linewidth}{@{}*{4}{L}@{}}
\optlabel{A}\,$2$ & \optlabel{B}\,$2^{n}$ & \correctlabel{C}\,$8$ & \optlabel{D}\,$4$
\end{tabularx}
\end{qitem}
\begin{solbox}{C}
\textbf{\color{brand}Correct answer.} \textcolor{okgreen}{\textbf{(C)}} $8$

\textbf{\color{brand}Concepts required.} Polynomial scaling: $O(n^3)$ means doubling $n$ multiplies the cost by $2^3$.

\textbf{\color{brand}Relevant formulae.} $(2n)^3/n^3 = 8$.

\textbf{\color{brand}Step-by-step explanation.} $(2n)^3 = 8n^3$. The cost is polynomial: $t = 2n$ controlled multiplications, each $O(n^2)$ gates. Classical sub-exponential algorithms (GNFS) grow much faster,
which is the source of Shor's advantage --- but only on a fault-tolerant machine. $2^n$ is the exponential cost of trial division.

\textbf{\color{brand}Why the other options are incorrect.}
\begin{itemize}[leftmargin=2.4em, itemsep=1pt, topsep=2pt]
\item[\textbf{(A)}] $2$ would be linear scaling, $O(n)$.
\item[\textbf{(B)}] $2^n$ is exponential growth, the cost of brute-force methods, not of Shor's algorithm.
\item[\textbf{(D)}] $4$ would be quadratic scaling, $O(n^2)$ --- e.g.\ the QFT alone.
\end{itemize}
\textbf{\color{brand}Key takeaway.} Shor is polynomial ($O(n^3)$ or better), GNFS sub-exponential --- the gap grows with $n$.
\end{solbox}
\qsep

\needspace{13\baselineskip}
\begin{qitem}{64}
\qmeta{Single-correct MCQ}{Semiclassical QFT}
In the ``semiclassical'' (one-control-qubit) implementation of the inverse QFT used by Beauregard, the $t = 2n$ counting qubits are replaced by
\begin{itemize}[leftmargin=2.6em, labelsep=0.4em, itemsep=2pt, topsep=3pt]
\item[\optlabel{A}] $2n$ ancilla qubits that are never measured
\item[\optlabel{B}] $n$ qubits, because the QFT needs only half the bits
\item[\optlabel{C}] no qubits: the QFT is computed classically before the circuit
\item[\correctlabel{D}] a single qubit that is measured and reset $2n$ times; earlier measurement results classically control the phase corrections
\end{itemize}
\end{qitem}
\begin{solbox}{D}
\textbf{\color{brand}Correct answer.} \textcolor{okgreen}{\textbf{(D)}} a single qubit that is measured and reset $2n$ times; earlier measurement results classically control the phase corrections

\textbf{\color{brand}Concepts required.} The inverse QFT followed by measurement can be performed one qubit at a time with classical feed-forward.

\textbf{\color{brand}Relevant formulae.} Bit $y_k$ measured $\to$ rotations $R_1(-\pi\sum_{j}y_j2^{-(j-k)})$ on the next qubit; $3n \to 2n + 3$ qubits.

\textbf{\color{brand}Step-by-step explanation.} The inverse QFT followed by measurement can be done bit by bit (Griffiths--Niu): measure, then apply classically controlled phases to the next qubit.
One control qubit is recycled, so the count drops from $3n$ to about $2n + 3$ ($n$ work, $n + 2$ ancilla for the Fourier-space adder, $1$ control). The
measurement statistics are unchanged.

\textbf{\color{brand}Why the other options are incorrect.}
\begin{itemize}[leftmargin=2.4em, itemsep=1pt, topsep=2pt]
\item[\textbf{(A)}] Ancilla qubits that are never measured would not reduce the qubit count; the point is to recycle one qubit.
\item[\textbf{(B)}] The QFT needs all $2n$ bits of precision; they are measured sequentially on one qubit, not reduced to $n$ qubits.
\item[\textbf{(C)}] The QFT acts on quantum amplitudes, which cannot be computed classically in advance.
\end{itemize}
\textbf{\color{brand}Key takeaway.} Measurement plus classical control trades qubits for mid-circuit measurements.
\end{solbox}
\qsep

\needspace{13\baselineskip}
\begin{qitem}{65}
\qmeta{Single-correct MCQ}{Approximate QFT}
For $N = 21$, $a = 2$, $t = 10$, an approximate inverse QFT keeps only controlled phases between qubits at distance $k - j \le m$. With $m = 2$ the
number of controlled phases and the (re-executed, exact) success probability are
\begin{itemize}[leftmargin=2.6em, labelsep=0.4em, itemsep=2pt, topsep=3pt]
\item[\correctlabel{A}] $17$ phases; $0.689$ (full QFT: $45$ phases, $0.817$)
\item[\optlabel{B}] $45$ phases; $0.817$
\item[\optlabel{C}] $17$ phases; $0.817$ --- truncation never changes the result
\item[\optlabel{D}] $2$ phases; $0.689$
\end{itemize}
\end{qitem}
\begin{solbox}{A}
\textbf{\color{brand}Correct answer.} \textcolor{okgreen}{\textbf{(A)}} $17$ phases; $0.689$ (full QFT: $45$ phases, $0.817$)

\textbf{\color{brand}Concepts required.} The approximate QFT drops small-angle controlled phases; it reduces the gate count at a small cost in accuracy.

\textbf{\color{brand}Relevant formulae.} Pairs at distance $\le2$: $9 + 8 = 17$; full $\binom{10}{2} = 45$.

\textbf{\color{brand}Step-by-step explanation.} Distance $1$: $9$ pairs, distance $2$: $8$ pairs, total $17$ (full QFT: $\binom{10}{2} = 45$). The re-executed simulation gave success $0.689$ for $m = 2$, $0.787$ for $m = 3$,
$0.811$ for $m = 4$ and $0.817$ for the full QFT. Very small cut-offs lose accuracy; $m \approx \log_2 n$ plus a few suffices asymptotically.

\textbf{\color{brand}Why the other options are incorrect.}
\begin{itemize}[leftmargin=2.4em, itemsep=1pt, topsep=2pt]
\item[\textbf{(B)}] $45$ phases and $0.817$ describe the full QFT, not the truncation $m = 2$.
\item[\textbf{(C)}] Truncation changes the amplitudes; the re-executed simulation shows a drop from $0.817$ to $0.689$.
\item[\textbf{(D)}] $m = 2$ is the maximum distance, not the number of gates; each distance contributes $t - d$ pairs.
\end{itemize}
\textbf{\color{brand}Key takeaway.} Moderate cut-offs ($m \approx 4$) lose little; aggressive ones lose a lot.
\end{solbox}
\qsep

\needspace{13\baselineskip}
\begin{qitem}{66}
\qmeta{Single-correct MCQ}{Controlled phases in the QFT}
The full inverse QFT on $t = 8$ counting qubits (notebook default for $N = 15$) contains how many controlled phase gates?
\par\smallskip\noindent\begin{tabularx}{\linewidth}{@{}*{4}{L}@{}}
\optlabel{A}\,$8$ & \correctlabel{B}\,$28$ & \optlabel{C}\,$64$ & \optlabel{D}\,$56$
\end{tabularx}
\end{qitem}
\begin{solbox}{B}
\textbf{\color{brand}Correct answer.} \textcolor{okgreen}{\textbf{(B)}} $28$

\textbf{\color{brand}Concepts required.} The QFT on $t$ qubits uses one controlled phase per qubit pair.

\textbf{\color{brand}Relevant formulae.} $t(t - 1)/2 = 28$ for $t = 8$.

\textbf{\color{brand}Step-by-step explanation.} One for each pair $j < k$: $\binom82 = \frac{8\cdot7}{2} = 28$ --- the recorded gate count \texttt{cr1: 28}. In general $t(t - 1)/2 = O(t^2)$, which for
large $t$ is dwarfed by the $O(n^3)$ cost of the modular exponentiation.

\textbf{\color{brand}Why the other options are incorrect.}
\begin{itemize}[leftmargin=2.4em, itemsep=1pt, topsep=2pt]
\item[\textbf{(A)}] $8$ is the number of Hadamards in the inverse QFT, not of controlled phases.
\item[\textbf{(C)}] $64 = t^2$ counts ordered pairs including $j = k$.
\item[\textbf{(D)}] $56 = t(t - 1)$ counts each unordered pair twice.
\end{itemize}
\textbf{\color{brand}Key takeaway.} QFT cost is $O(t^2)$ --- small compared with modular exponentiation.
\end{solbox}
\qsep

\needspace{13\baselineskip}
\begin{qitem}{67}
\qmeta{Single-correct MCQ}{Recorded gate counts}
The recorded gate counts of \texttt{Shor.ipynb} for $N = 15$ are \texttt{x: 1, ccustom\_unitary\_constant: 2, swap: 4, h: 16, cr1: 28} (and
\texttt{mz: 8}). Why are there $16$ Hadamard gates?
\begin{itemize}[leftmargin=2.6em, labelsep=0.4em, itemsep=2pt, topsep=3pt]
\item[\correctlabel{A}] $8$ create the superposition on the counting qubits and $8$ belong to the inverse QFT
\item[\optlabel{B}] $16 = 2^4$ is the dimension of the work register
\item[\optlabel{C}] each controlled multiplication contains $8$ Hadamards
\item[\optlabel{D}] Hadamards are inserted before every measurement of the work register
\end{itemize}
\end{qitem}
\begin{solbox}{A}
\textbf{\color{brand}Correct answer.} \textcolor{okgreen}{\textbf{(A)}} $8$ create the superposition on the counting qubits and $8$ belong to the inverse QFT

\textbf{\color{brand}Concepts required.} Hadamards appear at the start (superposition) and in the inverse QFT (one per qubit).

\textbf{\color{brand}Relevant formulae.} $8 + 8 = 16$; total $1 + 2 + 4 + 16 + 28 = 51$ gates without \texttt{mz}.

\textbf{\color{brand}Step-by-step explanation.} $t = 8$ Hadamards start the algorithm, and the inverse QFT applies one $H$ per qubit ($8$ more). The multipliers are single custom gates (no internal $H$),
and the work register is not measured. Total gates without \texttt{mz}: $1 + 2 + 4 + 16 + 28 = 51$, depth $20$.

\textbf{\color{brand}Why the other options are incorrect.}
\begin{itemize}[leftmargin=2.4em, itemsep=1pt, topsep=2pt]
\item[\textbf{(B)}] The coincidence $16 = 2^4$ is irrelevant; the count is $t + t$.
\item[\textbf{(C)}] The multipliers are single \texttt{ccustom\_unitary\_constant} gates; no Hadamards are inside them.
\item[\textbf{(D)}] The work register is not measured in the notebook; only \texttt{mz(count)} appears.
\end{itemize}
\textbf{\color{brand}Key takeaway.} Gate counts should be traced back to the circuit's structure.
\end{solbox}
\qsep

\needspace{13\baselineskip}
\begin{qitem}{68}
\qmeta{Single-correct MCQ}{No hardware execution}
Why does \texttt{Shor.ipynb} refuse to run on hardware targets?
\begin{itemize}[leftmargin=2.6em, labelsep=0.4em, itemsep=2pt, topsep=3pt]
\item[\optlabel{A}] CUDA-Q has no hardware targets
\item[\optlabel{B}] the inverse QFT cannot be executed on hardware
\item[\optlabel{C}] Shor's algorithm is mathematically impossible on hardware
\item[\correctlabel{D}] its modular multipliers are exact custom unitaries (permutation matrices) registered with \texttt{cudaq.register\_operation}; hardware would need them decomposed into native gates, which the notebook does not do
\end{itemize}
\end{qitem}
\begin{solbox}{D}
\textbf{\color{brand}Correct answer.} \textcolor{okgreen}{\textbf{(D)}} its modular multipliers are exact custom unitaries (permutation matrices) registered with \texttt{cudaq.register\_operation}; hardware would need them decomposed into native gates, which the notebook does not do

\textbf{\color{brand}Concepts required.} Custom matrix gates are a simulator convenience; hardware executes only native gate sets.

\textbf{\color{brand}Relevant formulae.} \texttt{cudaq.register\_operation(name, matrix)}; hardware $\to$ \texttt{ValueError} in the notebook.

\textbf{\color{brand}Step-by-step explanation.} The notebook raises a \texttt{ValueError} for hardware and points to the resource estimates. A dense $2^n\times2^n$ matrix gate is fine for a simulator but has no
native implementation; compiling it costs $O(n^2)$--$O(n^3)$ gates per multiplier. The inverse QFT itself uses standard gates.

\textbf{\color{brand}Why the other options are incorrect.}
\begin{itemize}[leftmargin=2.4em, itemsep=1pt, topsep=2pt]
\item[\textbf{(A)}] CUDA-Q supports several hardware providers; the restriction is the notebook's custom gates.
\item[\textbf{(B)}] The inverse QFT uses $H$, controlled phases and SWAPs, which compile to native gates.
\item[\textbf{(C)}] Shor's algorithm is a well-defined circuit; small instances have run on hardware with compiled (often simplified) circuits.
\end{itemize}
\textbf{\color{brand}Key takeaway.} Running on hardware requires a gate-level multiplier --- the hard engineering part.
\end{solbox}
\qsep

\needspace{13\baselineskip}
\begin{qitem}{69}
\qmeta{Single-correct MCQ}{CNOTs of a controlled multiplier}
In the gate-level circuit, ``multiply by $7$ modulo $15$'' is built as three left rotations of the 4 work bits (each $3$ SWAPs) followed by $X$ on all
four bits. Its controlled version, with controlled SWAPs of $8$ CNOTs and controlled-$X$ $=$ CNOT, uses how many CNOTs?
\par\smallskip\noindent\begin{tabularx}{\linewidth}{@{}*{4}{L}@{}}
\optlabel{A}\,$9$ & \optlabel{B}\,$40$ & \correctlabel{C}\,$76$ & \optlabel{D}\,$72$
\end{tabularx}
\end{qitem}
\begin{solbox}{C}
\textbf{\color{brand}Correct answer.} \textcolor{okgreen}{\textbf{(C)}} $76$

\textbf{\color{brand}Concepts required.} Count CNOTs gate by gate: controlled SWAP $= 8$, controlled $X = 1$.

\textbf{\color{brand}Relevant formulae.} $9\times8 + 4 = 76$; whole circuit: $76 + 48 + 3 = 127$.

\textbf{\color{brand}Step-by-step explanation.} $9$ controlled SWAPs $\times\,8 = 72$ CNOTs, plus $4$ CNOTs for the controlled $X$ gates: $76$. (Multiplication by $7$ is valid on $y = 1, \dots, 14$:
$7y \bmod 15 = 15 - (y\cdot2^3 \bmod 15)$.) With the controlled $\times4$ ($6$ controlled SWAPs $= 48$) and the QFT SWAP ($3$) the circuit has $127$ CNOTs.

\textbf{\color{brand}Why the other options are incorrect.}
\begin{itemize}[leftmargin=2.4em, itemsep=1pt, topsep=2pt]
\item[\textbf{(A)}] $9$ counts controlled SWAPs, not CNOTs.
\item[\textbf{(B)}] $40$ would use $4$ CNOTs per controlled SWAP; a Toffoli alone needs $6$.
\item[\textbf{(D)}] $72$ forgets the four controlled $X$ gates.
\end{itemize}
\textbf{\color{brand}Key takeaway.} Even multiplying a 4-bit number by $7$ costs $76$ CNOTs when controlled.
\end{solbox}
\qsep

\needspace{13\baselineskip}
\begin{qitem}{70}
\qmeta{Single-correct MCQ}{Random baseline}
In the gate-level circuit ($t = 3$, $a = 7$, $N = 15$) the counting outcome becomes uniformly random in the fully depolarized limit. The success
probability with the notebook's post-processing is then
\begin{itemize}[leftmargin=2.6em, labelsep=0.4em, itemsep=2pt, topsep=3pt]
\item[\optlabel{A}] $1/8$
\item[\correctlabel{B}] $3/8$: $y = 2, 4, 6$ give $r = 4$ and the factors; $y = 0$ gives nothing; odd $y$ give $r' = 8$, and $7^4 \equiv 1$ gives only trivial factors
\item[\optlabel{C}] $3/4$, the same as without noise
\item[\optlabel{D}] $0$
\end{itemize}
\end{qitem}
\begin{solbox}{B}
\textbf{\color{brand}Correct answer.} \textcolor{okgreen}{\textbf{(B)}} $3/8$: $y = 2, 4, 6$ give $r = 4$ and the factors; $y = 0$ gives nothing; odd $y$ give $r' = 8$, and $7^4 \equiv 1$ gives only trivial factors

\textbf{\color{brand}Concepts required.} A uniformly random outcome is the zero-information baseline for any noisy run.

\textbf{\color{brand}Relevant formulae.} Successful $y \in \{2, 4, 6\}$: $3/8 = 0.375$; odd $y \to 8$, $7^4 \equiv 1 \to$ trivial gcds.

\textbf{\color{brand}Step-by-step explanation.} $y/8$: $2/8 = 1/4$, $6/8 = 3/4$ give $4$; $4/8 = 1/2$ gives $4$ via multiples. For odd $y$ the denominator $8$ satisfies $7^8 \equiv 1$, but $7^{4} \equiv 1$, so
$\gcd(0, 15) = 15$, $\gcd(2, 15) = 1$ --- trivial. Hence $3/8 = 0.375$: a noisy run can still ``succeed'' by luck, which is why a success rate must be compared
with this random baseline.

\textbf{\color{brand}Why the other options are incorrect.}
\begin{itemize}[leftmargin=2.4em, itemsep=1pt, topsep=2pt]
\item[\textbf{(A)}] $1/8$ counts only one of the three successful outcomes.
\item[\textbf{(C)}] $3/4$ is the noiseless success probability; full depolarization destroys the peaks.
\item[\textbf{(D)}] $0$ ignores that random outcomes can still give $r = 4$ by chance.
\end{itemize}
\textbf{\color{brand}Key takeaway.} Compare noisy success rates with the random baseline, never with zero.
\end{solbox}
\qsep

\needspace{13\baselineskip}
\begin{qitem}{71}
\qmeta{Single-correct MCQ}{Memory for $N = 1023$}
Using the notebook's $3n$-qubit construction, the state-vector memory (16 bytes per amplitude) for $N = 1023$ is about
\par\smallskip\noindent\begin{tabularx}{\linewidth}{@{}*{2}{L}@{}}
\optlabel{A}\,$1.1$ TB ($36$ qubits) & \optlabel{B}\,$16.4$ kB \\[3pt]
\optlabel{C}\,$268$ MB ($24$ qubits) & \correctlabel{D}\,$17.2$ GB ($30$ qubits)
\end{tabularx}
\end{qitem}
\begin{solbox}{D}
\textbf{\color{brand}Correct answer.} \textcolor{okgreen}{\textbf{(D)}} $17.2$ GB ($30$ qubits)

\textbf{\color{brand}Concepts required.} Memory of the full state vector: $16\cdot2^{3n}$ bytes in the notebook's construction.

\textbf{\color{brand}Relevant formulae.} $1023$: $n = 10$, $q = 30$, $M = 16\cdot2^{30} = 17.2$ GB.

\textbf{\color{brand}Step-by-step explanation.} $1023$ has $10$ bits, so $3n = 30$ qubits: $16\times2^{30} = 17.18\times10^9$ bytes. This matches the notebook's reverse estimate that $80$ GB allows at most
$32$ qubits, i.e.\ $N \le 1023$. $N = 1024$ would need $11$ bits and $33$ qubits ($137$ GB).

\textbf{\color{brand}Why the other options are incorrect.}
\begin{itemize}[leftmargin=2.4em, itemsep=1pt, topsep=2pt]
\item[\textbf{(A)}] $1.1$ TB corresponds to $36$ qubits, i.e.\ a $12$-bit $N$.
\item[\textbf{(B)}] $16.4$ kB is $16\cdot2^{10}$: it uses $n$ qubits instead of $3n$.
\item[\textbf{(C)}] $268$ MB corresponds to $24$ qubits, i.e.\ an $8$-bit $N$ such as $143$.
\end{itemize}
\textbf{\color{brand}Key takeaway.} Simulation memory explodes: a few more bits of $N$ exhaust any computer.
\end{solbox}
\qsep

\needspace{13\baselineskip}
\begin{qitem}{72}
\qmeta{Single-correct MCQ}{Toy RSA}
Toy RSA: $p = 5$, $q = 11$, $e = 7$. The private exponent $d$ and the ciphertext of $m = 2$ are
\par\smallskip\noindent\begin{tabularx}{\linewidth}{@{}*{4}{L}@{}}
\correctlabel{A}\,$d = 23$; $c = 18$ & \optlabel{B}\,$d = 3$; $c = 18$ & \optlabel{C}\,$d = 23$; $c = 14$ & \optlabel{D}\,$d = 7$; $c = 2$
\end{tabularx}
\end{qitem}
\begin{solbox}{A}
\textbf{\color{brand}Correct answer.} \textcolor{okgreen}{\textbf{(A)}} $d = 23$; $c = 18$

\textbf{\color{brand}Concepts required.} RSA: $ed \equiv 1 \pmod{\varphi(N)}$, $c = m^e \bmod N$, $m = c^d \bmod N$.

\textbf{\color{brand}Relevant formulae.} $\varphi(55) = 40$, $7\cdot23 = 161 \equiv 1$; $2^7 = 128 \equiv 18 \pmod{55}$.

\textbf{\color{brand}Step-by-step explanation.} $\varphi(55) = 4\cdot10 = 40$; $7d \equiv 1 \pmod{40}$: $7\cdot23 = 161 = 4\cdot40 + 1$, so $d = 23$. $c = 2^7 \bmod 55 = 128 - 110 = 18$; check $18^{23} \bmod 55 = 2$. Shor's
algorithm breaks RSA by factoring $N = 55$, which reveals $\varphi(N)$ and hence $d$.

\textbf{\color{brand}Why the other options are incorrect.}
\begin{itemize}[leftmargin=2.4em, itemsep=1pt, topsep=2pt]
\item[\textbf{(B)}] $7\cdot3 = 21 \not\equiv 1 \pmod{40}$, so $d \ne 3$.
\item[\textbf{(C)}] $c = 2^7 \bmod 55 = 18$, not $14$.
\item[\textbf{(D)}] $d = e$ would need $49 \equiv 1 \pmod{40}$; and encryption does change the message.
\end{itemize}
\textbf{\color{brand}Key takeaway.} Factoring $N$ gives $\varphi(N)$ and hence $d$: that is the threat.
\end{solbox}
\qsep

\needspace{13\baselineskip}
\begin{qitem}{73}
\qmeta{Single-correct MCQ}{Post-quantum standards}
Which standard specifies NIST's lattice-based key-encapsulation mechanism intended to replace RSA/ECDH key exchange?
\par\smallskip\noindent\begin{tabularx}{\linewidth}{@{}*{2}{L}@{}}
\correctlabel{A}\,FIPS 203 (ML-KEM) & \optlabel{B}\,SP 800-90A (random bit generators) \\[3pt]
\optlabel{C}\,FIPS 197 (AES) & \optlabel{D}\,FIPS 186 (DSA/ECDSA)
\end{tabularx}
\end{qitem}
\begin{solbox}{A}
\textbf{\color{brand}Correct answer.} \textcolor{okgreen}{\textbf{(A)}} FIPS 203 (ML-KEM)

\textbf{\color{brand}Concepts required.} NIST's post-quantum standards: FIPS 203 (ML-KEM), FIPS 204 (ML-DSA), FIPS 205 (SLH-DSA), 2024.

\textbf{\color{brand}Relevant formulae.} ML-KEM: module learning with errors; no known efficient quantum algorithm.

\textbf{\color{brand}Step-by-step explanation.} NIST published FIPS 203 (ML-KEM, from CRYSTALS-Kyber), FIPS 204 (ML-DSA, from Dilithium) and FIPS 205 (SLH-DSA, from SPHINCS+) in August 2024. AES (FIPS 197) is
symmetric and not broken by Shor; FIPS 186 contains the signature schemes Shor's algorithm does break.

\textbf{\color{brand}Why the other options are incorrect.}
\begin{itemize}[leftmargin=2.4em, itemsep=1pt, topsep=2pt]
\item[\textbf{(B)}] SP 800-90A specifies random bit generators, not key exchange.
\item[\textbf{(C)}] FIPS 197 is AES, a symmetric cipher; it is not a key-encapsulation mechanism.
\item[\textbf{(D)}] FIPS 186 specifies DSA/ECDSA/RSA signatures --- schemes broken by Shor's algorithm.
\end{itemize}
\textbf{\color{brand}Key takeaway.} Replace RSA/ECC key exchange by ML-KEM (often in hybrid mode).
\end{solbox}
\qsep

\needspace{13\baselineskip}
\begin{qitem}{74}
\qmeta{Single-correct MCQ}{Impact on RSA, ECC, AES}
Which statement about the cryptographic impact of a large fault-tolerant quantum computer is correct?
\begin{itemize}[leftmargin=2.6em, labelsep=0.4em, itemsep=2pt, topsep=3pt]
\item[\optlabel{A}] hash functions such as SHA-256 are broken by Shor's algorithm
\item[\optlabel{B}] only RSA is affected; elliptic-curve cryptography is immune
\item[\correctlabel{C}] Shor's algorithm breaks RSA and elliptic-curve cryptography; AES is affected only by Grover's quadratic speed-up, so AES-256 remains adequate
\item[\optlabel{D}] Shor's algorithm breaks AES-256 in polynomial time
\end{itemize}
\end{qitem}
\begin{solbox}{C}
\textbf{\color{brand}Correct answer.} \textcolor{okgreen}{\textbf{(C)}} Shor's algorithm breaks RSA and elliptic-curve cryptography; AES is affected only by Grover's quadratic speed-up, so AES-256 remains adequate

\textbf{\color{brand}Concepts required.} Shor solves the hidden subgroup problem in abelian groups: factoring and discrete logarithms, including on elliptic curves.

\textbf{\color{brand}Relevant formulae.} Grover: $O(\sqrt{2^k})$ for a $k$-bit key: AES-256 $\to \approx2^{128}$ work.

\textbf{\color{brand}Step-by-step explanation.} Shor's period finding solves factoring and discrete logarithms in any abelian group, including elliptic-curve groups (ECDH, ECDSA). Symmetric ciphers and hash
functions have no such structure; Grover search gives at most a square-root speed-up, mitigated by doubling key length.

\textbf{\color{brand}Why the other options are incorrect.}
\begin{itemize}[leftmargin=2.4em, itemsep=1pt, topsep=2pt]
\item[\textbf{(A)}] Hash functions are affected only by generic (Grover-type) speed-ups, not by Shor.
\item[\textbf{(B)}] ECC relies on discrete logarithms, which Shor's algorithm also solves (with fewer qubits than RSA of equal security).
\item[\textbf{(D)}] Shor's algorithm needs algebraic periodicity; AES has none, and Grover gives only a square-root speed-up.
\end{itemize}
\textbf{\color{brand}Key takeaway.} Public-key cryptography falls; symmetric cryptography needs at most larger keys.
\end{solbox}
\qsep

\needspace{13\baselineskip}
\begin{qitem}{75}
\qmeta{Matching lists}{Fault-tolerance terms}
Match each fault-tolerance term (List-I) with its meaning (List-II).
\par\smallskip\noindent{\renewcommand{\arraystretch}{1.35}%
\begin{tabularx}{\linewidth}{@{}c>{\raggedright\arraybackslash}p{0.33\linewidth}@{\hspace{1.2em}}cL@{}}
\toprule
\multicolumn{2}{@{}l}{\textbf{List-I}} & \multicolumn{2}{l@{}}{\textbf{List-II}} \\ \midrule
\textbf{(P)} & logical qubit & \textbf{(1)} & minimum number of physical errors that can cause an undetected logical error \\
\textbf{(Q)} & physical qubit & \textbf{(2)} & a special resource state that enables non-Clifford gates such as $T$ or Toffoli in the surface code \\
\textbf{(R)} & code distance $d$ & \textbf{(3)} & an error-corrected qubit encoded in many noisy qubits \\
\textbf{(S)} & magic state & \textbf{(4)} & a noisy hardware qubit with error rates around $10^{-3}$ \\
\bottomrule
\end{tabularx}}
\par\smallskip\noindent\begin{tabularx}{\linewidth}{@{}*{2}{L}@{}}
\optlabel{A}\,P$\to$3, Q$\to$4, R$\to$2, S$\to$1 & \optlabel{B}\,P$\to$3, Q$\to$1, R$\to$4, S$\to$2 \\[3pt]
\optlabel{C}\,P$\to$4, Q$\to$3, R$\to$1, S$\to$2 & \correctlabel{D}\,P$\to$3, Q$\to$4, R$\to$1, S$\to$2
\end{tabularx}
\end{qitem}
\begin{solbox}{D}
\textbf{\color{brand}Correct answer.} \textcolor{okgreen}{\textbf{(D)}} P$\to$3, Q$\to$4, R$\to$1, S$\to$2

\textbf{\color{brand}Concepts required.} Fault tolerance encodes logical qubits in many physical qubits and supplies non-Clifford gates via magic states.

\textbf{\color{brand}Relevant formulae.} Surface code: $\approx2d^2$ physical qubits per logical qubit; logical error $\sim(p/p_{\text{th}})^{(d+1)/2}$.

\textbf{\color{brand}Step-by-step explanation.} Logical qubit: encoded, error-corrected (3); physical qubit: the noisy hardware element (4); code distance: the number of errors needed for an undetected
logical error --- a surface-code patch uses about $2d^2$ physical qubits (1); magic states: distilled or cultivated resources for $T$/Toffoli (2). Shor's
arithmetic is dominated by Toffoli gates, so magic-state production dominates its cost.

\textbf{\color{brand}Why the other options are incorrect.}
\begin{itemize}[leftmargin=2.4em, itemsep=1pt, topsep=2pt]
\item[\textbf{(A)}] This swaps code distance and magic states; the distance is a number of errors, magic states are resources.
\item[\textbf{(B)}] This gives the physical qubit the meaning of code distance; a physical qubit is the noisy hardware element.
\item[\textbf{(C)}] This swaps logical and physical qubits; the logical qubit is the encoded one.
\end{itemize}
\textbf{\color{brand}Key takeaway.} Shor's cost on real hardware is dominated by error correction and magic states.
\end{solbox}
\qsep

\needspace{13\baselineskip}
\begin{qitem}{76}
\qmeta{Matching lists}{Qubit counts of constructions}
Match each construction for an $n$-bit $N$ (List-I) with its logical-qubit count (List-II).
\par\smallskip\noindent{\renewcommand{\arraystretch}{1.35}%
\begin{tabularx}{\linewidth}{@{}c>{\raggedright\arraybackslash}p{0.33\linewidth}@{\hspace{1.2em}}cL@{}}
\toprule
\multicolumn{2}{@{}l}{\textbf{List-I}} & \multicolumn{2}{l@{}}{\textbf{List-II}} \\ \midrule
\textbf{(P)} & \texttt{Shor.ipynb} ($t = 2n$ counting $+$ $n$ work) & \textbf{(1)} & $n + 1$ \\
\textbf{(Q)} & Beauregard (2003) & \textbf{(2)} & $3n + 0.002\,n\log_2 n$ \\
\textbf{(R)} & Gidney--Ekerå (2019) & \textbf{(3)} & $3n$ \\
\textbf{(S)} & one recycled counting qubit $+$ $n$ work qubits with exact (oracle) multipliers & \textbf{(4)} & $2n + 3$ \\
\bottomrule
\end{tabularx}}
\par\smallskip\noindent\begin{tabularx}{\linewidth}{@{}*{2}{L}@{}}
\optlabel{A}\,P$\to$3, Q$\to$4, R$\to$1, S$\to$2 & \optlabel{B}\,P$\to$3, Q$\to$2, R$\to$4, S$\to$1 \\[3pt]
\correctlabel{C}\,P$\to$3, Q$\to$4, R$\to$2, S$\to$1 & \optlabel{D}\,P$\to$4, Q$\to$3, R$\to$2, S$\to$1
\end{tabularx}
\end{qitem}
\begin{solbox}{C}
\textbf{\color{brand}Correct answer.} \textcolor{okgreen}{\textbf{(C)}} P$\to$3, Q$\to$4, R$\to$2, S$\to$1

\textbf{\color{brand}Concepts required.} Different constructions trade qubits for gates and measurements.

\textbf{\color{brand}Relevant formulae.} $3n$ (notebook), $2n + 3$ (Beauregard), $3n + 0.002n\log_2n$ (GE2019), $n + 1$ (idealised semiclassical).

\textbf{\color{brand}Step-by-step explanation.} Notebook: $2n + n = 3n$ (3). Beauregard: $2n + 3$ (4). Gidney--Ekerå: $3n + 0.002n\log_2 n$, about $6190$ for $n = 2048$ (2). With the semiclassical QFT and
idealised matrix multipliers one needs only $n + 1$ (1) --- but real multipliers need ancilla qubits.

\textbf{\color{brand}Why the other options are incorrect.}
\begin{itemize}[leftmargin=2.4em, itemsep=1pt, topsep=2pt]
\item[\textbf{(A)}] This swaps Gidney--Ekerå and the idealised semiclassical count.
\item[\textbf{(B)}] This swaps Beauregard and Gidney--Ekerå; GE2019 uses slightly more than $3n$ qubits but far fewer gates.
\item[\textbf{(D)}] This swaps the notebook and Beauregard: the notebook uses $t + n = 3n$.
\end{itemize}
\textbf{\color{brand}Key takeaway.} Qubit counts alone do not decide feasibility; gate counts and error rates matter too.
\end{solbox}
\qsep

\needspace{13\baselineskip}
\begin{qitem}{77}
\qmeta{Assertion--Reason}{Harvest now, decrypt later}
\textbf{Assertion (A):} Migrating to post-quantum cryptography is urgent even though no current quantum computer can run Shor's algorithm at
cryptographic sizes.\\[1mm]
\textbf{Reason (R):} Shor's algorithm also solves the discrete-logarithm problem.
\begin{itemize}[leftmargin=2.6em, labelsep=0.4em, itemsep=2pt, topsep=3pt]
\item[\optlabel{A}] Both Assertion (A) and Reason (R) are true, and R is the correct explanation of A
\item[\correctlabel{B}] Both A and R are true, but R is \emph{not} the correct explanation of A
\item[\optlabel{C}] A is true, but R is false
\item[\optlabel{D}] A is false, but R is true
\end{itemize}
\end{qitem}
\begin{solbox}{B}
\textbf{\color{brand}Correct answer.} \textcolor{okgreen}{\textbf{(B)}} Both A and R are true, but R is \emph{not} the correct explanation of A

\textbf{\color{brand}Concepts required.} Data with long confidentiality lifetimes is already at risk from future quantum computers.

\textbf{\color{brand}Relevant formulae.} Urgent if (data lifetime) $+$ (migration time) $>$ (time to a cryptographically relevant quantum computer).

\textbf{\color{brand}Step-by-step explanation.} \textbf{A is true}: encrypted data recorded today can be decrypted later (``harvest now, decrypt later''), and migrations take years. \textbf{R is true}: period
finding solves discrete logarithms, so Diffie--Hellman and ECC fall too. But R does not explain the \emph{urgency} --- that comes from data lifetime and
migration time. Answer: (B).

\textbf{\color{brand}Why the other options are incorrect.}
\begin{itemize}[leftmargin=2.4em, itemsep=1pt, topsep=2pt]
\item[\textbf{(A)}] is wrong because R, although true, does not explain the urgency; that comes from harvest-now-decrypt-later and slow migration.
\item[\textbf{(C)}] is wrong because R is true: Shor's algorithm solves discrete logarithms.
\item[\textbf{(D)}] is wrong because A is true.
\end{itemize}
\textbf{\color{brand}Key takeaway.} The threat is about the future decryption of today's recorded data.
\end{solbox}
\qsep

\needspace{13\baselineskip}
\begin{qitem}{78}
\qmeta{Statement I / II}{Qubits and repeated squaring}
\textbf{Statement I:} In the notebook's construction, a $32$-bit $N$ needs $96$ qubits.\\[1mm]
\textbf{Statement II:} Repeated squaring lets $a^x \bmod N$ be implemented with $t$ controlled multiplications by the classically precomputed constants $a^{2^j} \bmod N$.
\begin{itemize}[leftmargin=2.6em, labelsep=0.4em, itemsep=2pt, topsep=3pt]
\item[\correctlabel{A}] Both Statement I and Statement II are true
\item[\optlabel{B}] Both Statement I and Statement II are false
\item[\optlabel{C}] Statement I is true, but Statement II is false
\item[\optlabel{D}] Statement I is false, but Statement II is true
\end{itemize}
\end{qitem}
\begin{solbox}{A}
\textbf{\color{brand}Correct answer.} \textcolor{okgreen}{\textbf{(A)}} Both Statement I and Statement II are true

\textbf{\color{brand}Concepts required.} Qubit count $3n$ and the repeated-squaring structure of modular exponentiation.

\textbf{\color{brand}Relevant formulae.} $3\cdot32 = 96$; $a^x = \prod_{j=0}^{t-1}(a^{2^j})^{x_j}$.

\textbf{\color{brand}Step-by-step explanation.} \textbf{I is true}: $3n = 96$ (the notebook's reverse estimate: $100$ qubits allow $33$ bits). \textbf{II is true}: $a^x = \prod_j(a^{2^j})^{x_j}$, so counting qubit $j$
controls the multiplication by $a^{2^j} \bmod N$, computed classically. Answer: (A).

\textbf{\color{brand}Why the other options are incorrect.}
\begin{itemize}[leftmargin=2.4em, itemsep=1pt, topsep=2pt]
\item[\textbf{(B)}] is wrong because both statements are true.
\item[\textbf{(C)}] is wrong because Statement II is true: the controlled multipliers by $a^{2^j}$ are exactly the notebook's construction.
\item[\textbf{(D)}] is wrong because Statement I is true: $t + n = 64 + 32 = 96$.
\end{itemize}
\textbf{\color{brand}Key takeaway.} $t$ controlled multiplications with classically precomputed constants implement $a^x \bmod N$.
\end{solbox}
\qsep

\needspace{20\baselineskip}
\begin{passagebox}{Paragraph for Questions 79 and 80}
A gate-level version of the order-finding circuit for $N = 15$, $a = 7$ uses $t = 3$ counting qubits and $4$ work qubits. The
multiplications by $7$ and $4 = 7^2 \bmod 15$ are compiled into controlled SWAPs (each CNOT--Toffoli--CNOT, Toffoli $=$ 6 CNOTs) and CNOTs; the
inverse QFT uses $3$ CNOTs (one SWAP) and $3$ controlled phases. In total the circuit has $127$ CNOTs. It was simulated with the CUDA-Q
\texttt{density-matrix-cpu} target and a depolarizing noise model: a one-qubit channel with parameter $p$ after every $H$, $R_z$ and $X$, and a
two-qubit channel with parameter $p$ after every CNOT and controlled phase. The probability that one shot yields the factors $3 \times 5$ was
(re-executed, exact from the density matrix):
\begin{center}\begin{tabular}{lccccc}\toprule
$p$ & $0$ & $0.001$ & $0.003$ & $0.01$ & $0.03$\\\midrule
success probability & $0.750$ & $0.740$ & $0.726$ & $0.703$ & $0.665$\\\bottomrule
\end{tabular}\end{center}
If the counting register were completely random (uniform $y$), the success probability would be $0.375$. These are simulator results, not
measurements on a quantum processor.
\end{passagebox}
\needspace{13\baselineskip}
\begin{qitem}{79}
\qmeta{Paragraph based}{Error-free probability vs success}
At $p = 0.01$, the probability that \emph{all} $127$ CNOTs are free of error is about $0.99^{127}$. Compare it with the success probability.
\begin{itemize}[leftmargin=2.6em, labelsep=0.4em, itemsep=2pt, topsep=3pt]
\item[\optlabel{A}] $0.99^{127} \approx 0.88$, equal to the success probability
\item[\correctlabel{B}] $0.99^{127} \approx 0.28$, yet the success probability is $0.703$: many errors do not change $y$ or still lead to a correct order
\item[\optlabel{C}] $0.99^{127} \approx 0.01$, so the result is pure noise
\item[\optlabel{D}] $0.99^{127} \approx 0.28$, so the success probability must be at most $0.28$
\end{itemize}
\end{qitem}
\begin{solbox}{B}
\textbf{\color{brand}Correct answer.} \textcolor{okgreen}{\textbf{(B)}} $0.99^{127} \approx 0.28$, yet the success probability is $0.703$: many errors do not change $y$ or still lead to a correct order

\textbf{\color{brand}Concepts required.} The error-free probability bounds nothing: many errors are harmless for the final answer.

\textbf{\color{brand}Relevant formulae.} $0.99^{127} = e^{-1.276} = 0.279$; measured (simulated) success $0.703$.

\textbf{\color{brand}Step-by-step explanation.} $0.99^{127} = e^{127\ln0.99} = e^{-1.276} = 0.279$. The success probability is not bounded by the error-free probability: errors on work qubits after their last
use, phase errors that commute with measurement, or errors that shift $y$ to another successful value do not cause failure. Also, a random $y$ still succeeds
with probability $0.375$.

\textbf{\color{brand}Why the other options are incorrect.}
\begin{itemize}[leftmargin=2.4em, itemsep=1pt, topsep=2pt]
\item[\textbf{(A)}] $0.99^{127} \approx 0.28$, not $0.88$; $0.88 \approx 0.99^{13}$.
\item[\textbf{(C)}] $0.99^{127} \approx 0.28$, not $0.01$; and $0.703$ is well above the random baseline $0.375$.
\item[\textbf{(D)}] Errors that do not change $y$, or move it to another successful value, do not cause failure; $0.703 > 0.28$ contradicts this option.
\end{itemize}
\textbf{\color{brand}Key takeaway.} Estimate error impact by simulation, not by the naive product of gate fidelities.
\end{solbox}
\qsep

\needspace{13\baselineskip}
\begin{qitem}{80}
\qmeta{Paragraph based}{Interpreting noisy results}
Which conclusion is supported by the table?
\begin{itemize}[leftmargin=2.6em, labelsep=0.4em, itemsep=2pt, topsep=3pt]
\item[\optlabel{A}] the results show that a quantum processor factored $15$
\item[\optlabel{B}] noise increases the success probability
\item[\optlabel{C}] at $p = 0.03$ the circuit is completely random
\item[\correctlabel{D}] the success probability decreases monotonically from $0.750$ towards the random baseline $0.375$; at $p = 0.03$ it is $0.665$, still above the baseline
\end{itemize}
\end{qitem}
\begin{solbox}{D}
\textbf{\color{brand}Correct answer.} \textcolor{okgreen}{\textbf{(D)}} the success probability decreases monotonically from $0.750$ towards the random baseline $0.375$; at $p = 0.03$ it is $0.665$, still above the baseline

\textbf{\color{brand}Concepts required.} Noise interpolates between the ideal distribution and the uniform baseline.

\textbf{\color{brand}Relevant formulae.} Ideal $0.750$; baseline $0.375$; $p = 0.03$: $0.665$.

\textbf{\color{brand}Step-by-step explanation.} The values fall $0.750 \to 0.740 \to 0.726 \to 0.703 \to 0.665$; the gap to the uniform baseline ($0.375$) shrinks but has not vanished. These are
density-matrix \emph{simulations} with an idealised noise model --- no quantum processor was used, and factoring $15$ with a compiled circuit that uses the known
order is a demonstration, not evidence of quantum advantage.

\textbf{\color{brand}Why the other options are incorrect.}
\begin{itemize}[leftmargin=2.4em, itemsep=1pt, topsep=2pt]
\item[\textbf{(A)}] These are simulations on \texttt{density-matrix-cpu}; no quantum processor was used.
\item[\textbf{(B)}] The values decrease with $p$; noise never helped here.
\item[\textbf{(C)}] $0.665$ is far above the random baseline $0.375$; the circuit is not yet completely random.
\end{itemize}
\textbf{\color{brand}Key takeaway.} Label simulated results as simulations; compare them with the right baseline.
\end{solbox}
\qsep

\end{document}